% 本文件由 LaTeX 排版 Agent 根据有效 translation.json 自动生成。
% author=John Chrysostom; work_id=2302; title=斐理伯书讲道集

\chapter{斐理伯书讲道集}

% structure: p0001 source=h1 target=omitted reason=page-title-already-rendered-by-parent

引言\newline{}第一篇讲道\newline{}第二篇讲道\newline{}第三篇讲道\newline{}第四篇讲道\newline{}第五篇讲道\newline{}第六篇讲道\newline{}第七篇讲道\newline{}第八篇讲道\newline{}第九篇讲道\newline{}第十篇讲道\newline{}第十一篇讲道\newline{}第十二篇讲道\newline{}第十三篇讲道\newline{}第十四篇讲道\newline{}第十五篇讲道\par

\SourceRecord{Translated by John A. Broadus.；Nicene and Post-Nicene Fathers, First Series；Vol. 13.；Edited by Philip Schaff.；Buffalo, NY: Christian Literature Publishing Co.,；1889.；<http://www.newadvent.org/fathers/2302.htm>.}

% structure: p0001 source=h1 target=section reason=page-title

\section{\WorkTitle{斐理伯书讲道}}

斐理伯人属马其顿的一座城，这座城是一个殖民地，正如路加所说。在这里，那位卖紫布的妇人归化了，她是一位罕有的虔敬和留心的妇女。在这里，会堂长信了。在这里，保禄与息拉一同受了鞭打。在这里，官长请求他们离开，并惧怕他们，宣讲有了一个辉煌的开端。而他自己给他们作了许多崇高的见证，称他们为自己的冠冕，说他们受了许多苦。因为他说：\BibleQuote{给你们，蒙天主恩赐，不仅为相信祂，也为祂受苦。} \BibleRef{斐 1:29} 但当他给他们写信时，他正被捆绑。因此他说：\Quote{以致我身上的锁链在基督内显明了整个御营，} 他将尼禄的宫殿称为御营。但他曾被捆绑，又被释放，他向弟茂德表明此事，说：\BibleQuote{在我初次辩护时，没有人支持我，众人都离弃了我；愿这罪不归于他们。但主站在我旁，加给了我力量。} \BibleRef{弟后 4:16} 他所说的锁链，是指那次辩护之前的捆绑。因为弟茂德当时不在场，这是明显的：他说：\BibleQuote{在我初次辩护时，没有人支持我；} 他若早已知道，就不会这样写信给他了。但当他写这封信时，弟茂德与他在一起。他藉着自己所说的话表明此事：\BibleQuote{但我在主耶稣内希望不久给你们派弟茂德去。} \BibleRef{斐 2:19} 又说：\Quote{我一旦看清了自己的事，就希望立刻派他去。} 因为他被解开了锁链，在到他们那里之后，又被捆绑了。但如果他说：\Quote{是的，我正被奠祭在你们信德的祭献和服务上，} 这并不是说这事已成事实，而是等于说：\Quote{每当这事发生，我就喜乐，} 以提升他们因他的锁链而有的沮丧。因为他当时不会死，从他所说的话可以明显看出：\BibleQuote{但我在主内希望，我自己也快要到你们那里去。} \BibleRef{斐 2:24} 又说：\Quote{怀着这种信心，我知道我必生存，并且与你们众人一同存留。}\par

2. 但斐理伯人曾派厄帕夫洛狄托到他那里，给他带钱，并打听他的事情，因为他们对他非常友爱。因为他们派了他，听他本人说：\Quote{我一切全有，并且富足；我已满足，因从厄帕夫洛狄托收到了你们送来的东西。} 同时他们也派他去打听这些事。因为他也派他去打听这事，他在书信开头就写自己的事，说：\BibleQuote{但我愿意你们知道，我所遭遇的事反而更促进了福音的进展。} \BibleRef{斐 1:12} 又说：\Quote{我希望不久给你们派弟茂德去，好使我得知你们的情况，获得安慰。} 这个\Quote{我也，}好像他意指\Quote{正如你们为了完全确信派人来打听我的事一样，同样\InnerQuote{我也，}好使我得知你们的事而获得安慰。} 既然他们已经很长时间没有派人来（他藉此证明说：\BibleQuote{如今你们终于又对我表示关心} \BibleRef{斐 4:10}），而后来他们听说他被捆绑\BibleRef{斐 2:26}；因为，如果他们听说厄帕夫洛狄托病了，而他并不像保禄那样是非常显著的人，那么他们更会听说保禄的事，他们感到不安是合理的；因此，在书信的开头，他给了他们许多关于他锁链的安慰，表明他们不仅不应该不安，甚至应该喜乐。然后他给他们关于同心合意和谦卑的劝告，教导他们这是他们最大的安全，这样他们就能轻易战胜敌人。因为令你们的老师痛苦的，不是被捆绑，而是他们的弟子不同心。因为前者甚至促进福音，而后者却造成分裂。\par

3. 因此，在劝勉他们同心合意，并指出同心源自谦卑，然后又向那些以基督教为幌子到处败坏教义的犹太人射出一箭，称他们为\BibleQuote{狗}和\BibleQuote{作恶的工人}\BibleRef{斐 3:2}，并劝诫远离他们，教导应当听从谁，又详细谈论道德问题，整顿他们，使他们回归自己，说\BibleQuote{主已经近了}\BibleRef{斐 4:5}之后，他以惯有的智慧也提到了所送来的东西，然后给予他们丰富的安慰。他在写作中似乎给予他们特别的尊荣，从未在什么地方写任何责备的话，这证明他们的德行，因为他们没有给他们的导师任何把柄，而他给他们写信也不是责备的口吻，而始终是鼓励的口吻。正如我一开始也说过，这城市对信仰表现出极大的热忱；连狱卒（你们知道那是一个充满邪恶的行当）在见到一个奇迹后，立刻跑到他们那里，并同全家受了洗。因为他独自看到了所发生的奇迹，但他所获得的益处却不独属于他，而是同他的妻子和全家共享。不仅如此，那些鞭打他的官长们似乎也是出于一时冲动而非恶意，因为他们立刻派人释放他，后来又感到害怕。他不仅在他们信德的方面、在危险的方面为他们作证，也在善行的方面为他们作证，他说：\BibleQuote{在福音开始时，你们就一次两次地送东西来供应我的需要}\BibleRef{斐 4:15}，当时别人都没有这样做；因为他说：\Quote{没有教会与我合伙，有授受的关系}；而他们的间断乃是由于缺乏机会，而非出于选择，他说：\BibleQuote{你们原没有忽略我，只是缺乏机会。}\BibleRef{斐 4:10} 我们既知道这些事，又拥有如此多的榜样，以及他对他们的爱---因为他非常爱他们，这表现在他的这句话上：\BibleQuote{我没有别人像他那样，会真心实意地关心你们的事}\BibleRef{斐 2:20}；又说：\Quote{因为我把你们放在心上，也在我的锁链中}---\par

4. 我们既知道这些事，也使自己配得上这样的榜样，预备为基督受苦。但如今迫害已不再。所以，如果没有别的，就让我们效法他们行善的热忱，不要以为自己若已经捐献了一次两次就完成了全部。因为我们必须终生行善。因为我们不是一次讨天主的喜悦，而是恒常不断。赛跑的人如果跑了十圈，却放弃最后一圈，就前功尽弃了；同样，我们如果开始行善，后来又懈怠，就丧失了一切，败坏了一切。请听那有益的训诫说：\BibleQuote{不要使仁慈和忠信离开你}\BibleRef{箴 3:3}。他并不是说做一次、两次、三次、十次、百次，而是持续地：\BibleQuote{不要使它们离开你。}他也没有说\Quote{你不要离开它们}，而是说\BibleQuote{不要使它们离开你}，显示我们需要它们，而非它们需要我们；教导我们应当尽力留住它们。又说：\BibleQuote{要把它系在你的颈项上。}因为富贵人家的孩子颈项上挂着金饰，从不摘下，因为这显明他们高贵的出身；同样，我们也应常把仁慈佩戴在身上，显明我们是那慈悲之父的子女，\BibleQuote{祂使太阳升起，光照恶人，也光照善人}\BibleRef{玛 5:45}。你说：\Quote{但那些不信的人并不相信。}那么我说，如果我们这样行，他们就因此而相信了。如果他们看到我们怜悯众人，并登记在祂门下以祂为师，他们就知道我们这样行是效法祂。因为经上说：\Quote{仁慈和真实的信德。}他称\Quote{真实}很好。因为祂不愿这信德来自抢夺或欺诈，因为那就不是\Quote{信}，也不是\Quote{真实。}所以他说，你不要这样，但要以仁慈持守信德。\par

让我们佩上这装饰。让我们为自己的灵魂编一条金链，我指的是怜悯之链，趁我们还在此世。因为若这时代过去，我们便不能再使用它了。为何？因为那里没有穷人，那里没有财富，那里不再有匮乏。趁我们仍是孩童时，不要剥夺自己这装饰。因为如同孩童，若他们成了成人，这些就被取下，而他们得以换上别的装饰；我们也是如此。那里不再有金钱的施舍，而是别的、远为高贵的。那么，不要使我们自己失去这个！使我们的灵魂显得美丽！施舍是伟大的，美丽的，可敬的，这恩赐是伟大的，但良善更为伟大。若我们学着轻视财富，我们还会学到其他事。看，从这生出多少好事！那施舍得当的人，学会轻视财富。那学会轻视财富的人，已砍断了万恶之根。因此他所得的好处并不止于施舍的报偿与酬报，还在于他的灵魂变得哲思、高超而富有。那施舍的人受教导，不再羡慕财富与黄金。这教训一旦铭刻于他心，他便向攀登天堂迈出了一大步，并砍除了无数争斗、争执、嫉妒与沮丧的缘由。因为你们知道，你们也知道，万物都为财富而作，无数战争为财富而起。但那学会轻视财富的人，将自己置身于宁静的港湾，他不再害怕损失。因为施舍教给了他这些。他不再贪图邻人的东西；他怎能呢？既然他舍弃自己的并给予？他不再嫉妒富人；他怎能呢？既然他甘愿变穷？他擦亮了自己灵魂的眼睛。而这些都是在此世。但在来世，他所获得的福分无法尽述。他不会与愚拙的童女一起留在外面，而会与那些明智的童女一同进入，与新郎在一起，他的灯光明亮。虽然那些童女在守贞中忍受了艰辛，但尚未尝过这些艰辛的他，将比她们更好。这就是怜悯之德的力量。她带着她的乳儿，以极大的勇气进入。因为她认识天上的守门者，那些守护婚筵之门者，不仅认识，而且受尊敬；凡她所认识为敬重她的人，她会以极大的勇气领入，无人会反对，全都让路。因为若她将天主自天降下，并说服祂成为人，那么她将更能把一个人提升至天上；因为她的权能是伟大的。如果因怜悯与慈爱，天主成为人，并说服自己成为仆人，那么更将把祂的仆人带入祂自己的家。让我们爱她，让我们倾心于她，不是一天，也不是两天，而是我们的一生，好使她承认我们。若她承认我们，主也必承认我们。若她否认我们，主也必否认我们，并说：\Quote{我不认识你们。}但愿我们不会听到这声音，而是那有福的话：\BibleQuote{你们蒙我父降福的，来继承自创世以来为你们预备的国度吧。}\BibleRef{玛 25:34}愿我们众人获得这福分，藉着祂的恩宠与慈爱，在我们的主耶稣基督内，愿光荣、权能、尊崇归于祂——及圣父及圣神，自今而始，直至永远，永世无尽。阿们。\par

\SourceRecord{Translated by John A. Broadus.；Nicene and Post-Nicene Fathers, First Series；Vol. 13.；Edited by Philip Schaff.；Buffalo, NY: Christian Literature Publishing Co.,；1889.；<http://www.newadvent.org/fathers/230200.htm>.}

% structure: p0001 source=h1 target=section reason=page-title

\section{斐理伯书讲道辞（一）}

\BibleRef{斐 1:1,2}\par

\BibleQuote{保禄和弟茂德，耶稣基督的仆人，致书给在斐理伯、在基督耶稣内的众位圣人，以及同为主教者和执事们：愿恩宠与平安，由我们的父天主和主耶稣基督赐与你们。}\par

这里，他写信给同等尊荣的人，并未列出他师傅的品位，而是用了另一个，且是重大的品位。那是什么呢？他自称为\Quote{仆人}，而非宗徒。因为这个品位确实也是伟大的，且是一切美事的总纲，就是作基督的仆人，而不仅是被称为仆人。\Quote{基督的仆人}，这实在是自由的人，关乎罪；且作为真正的仆人，他不作任何别人的仆人，否则就不是完全作基督的仆人了。同样，在写给罗马人时，他也说：\BibleQuote{保禄，耶稣基督的仆人。}\BibleRef{罗 1:1}。但写给格林多人及弟茂德时，他自称\Quote{宗徒}。那么这是为甚么呢？不是因为他们胜于弟茂德。远非如此。而是他尊敬他们，格外重视他们，胜过其他他所写信的人。因为他也在他们身上见证了许多美德。此外，在那里他即将安排许多事，因此以宗徒的身份自居。但在此时，他并未给他们任何他们自己不能明了的命令。\par

\Quote{致书给在斐理伯、在基督耶稣内的众位圣人。}由于犹太人很可能从最初的谕言就自称\Quote{圣人}，当他们被称为\BibleQuote{圣洁的子民，天主特选的子民}\BibleRef{出 19:6}时；因此他补充说：\Quote{在基督耶稣内的众位圣人。}唯独这些人是圣的，那些人是世俗的。\Quote{同為主教者们和执事们。}这是甚么意思？一座城市有几位主教吗？当然不是；但他称长老们如此。因为那时他们还互换称呼，主教也被称为执事。为此，在写给弟茂德时，他说：\BibleQuote{尽你的职务}，那时他是主教。因为他是主教，可见于他对他说：\BibleQuote{不可轻易给人覆手}\BibleRef{弟前 5:22}。又说：\BibleQuote{藉着预言，并藉着长老团的覆手，所赐给你的}\BibleRef{弟前 4:14}。然而长老们不会给主教覆手。同样，在写给弟铎时，他说：\BibleQuote{我从前留你在克里特，是要你将所缺的事办妥当，并在各城设立长老，正如我所吩咐你的。若有无可指摘的人，只作一个妻子的丈夫}\BibleRef{铎 1:5}；这是对主教说的。说了这话后，他立刻加上：\BibleQuote{因为监督必须无可指摘，如同天主的管家，不任性}\BibleRef{铎 1:7}。因此，如我所说，古代的长老们都称为基督的主教和执事，主教们则称为长老；因而直到如今，许多主教写信时写\Quote{致我同司铎者}和\Quote{致我同执事者}。但另一方面，专有名称清楚地归属于各人，即主教与长老。\Quote{同為主教者们}，他说，和执事们。\par

第二节。\BibleQuote{愿恩宠与平安，由我们的父天主和主耶稣基督赐与你们。}\BibleRef{斐 1:2}\par

为何他别处从不致书给圣职人员，不论在罗马、格林多、厄弗所或任何地方，而是普遍致书给\Quote{众圣人、信者、蒙爱者}，这里却致书给圣职人员呢？因为他们曾派遣、结出果子，也曾打发厄帕夫洛狄托到他那里去。\par

第三节。\BibleQuote{我每想起你们，就感谢我的天主。}\BibleRef{斐 1:3}\par

他在另一著作中说：\BibleQuote{应服从那些领导你们的人，且要顺从，因为他们为了你们的灵魂常警醒，好像要交账的人；使他们能喜乐地做这事，而不忧愁。}\BibleRef{希 13:17}如果那\Quote{忧愁}是由于门徒的邪恶，那么\Quote{喜乐}地做这事是由于他们的进步。每当我记起你们，我就光荣天主。但他做这事，是由于知道他们有许多善事。他说：我既光荣，又祈祷。我不因你们在德行上进步而停止为你们祈祷。而是\Quote{我每想起你们，就感谢我的天主，}\par

第四节。\BibleQuote{我每次祈祷，为你们众人祈祷时，常是喜乐地祈祷。}\BibleRef{斐 1:4}\par

\Quote{常是}，不仅在我祈祷时。\Quote{喜乐地}。因为也有可能忧愁地这样做，如他别处所说：\BibleQuote{我先前从心里痛苦、忧愁，流了许多泪，给你们写了信}\BibleRef{格后 2:4}。\par

第五节。\BibleQuote{因为你们从第一天直到现在，在传扬福音的事上，有份于我的恩宠。}\BibleRef{斐 1:5}\par

他对他们的作证是伟大的，而且非常伟大，是人对宗徒和福传者所能作的证词。他说：\Quote{你们没有因为被托付了一座城，就只关心那座城，反而毫不遗漏地分享我的劳苦，随时随处与我同工，参与我的宣讲。不是一次、两次或三次，而是从你们信主那时起直到如今，你们时常显出宗徒般的热忱。}看，那些在罗马的人怎样离弃了他；因为听他说：\Quote{你知道，所有在亚细亚的人都离弃了我。}\BibleRef{弟后 1:15} 又说：\Quote{德玛斯离弃了我}，和\Quote{在我初次辩护时，没有人支持我。}\BibleRef{弟后 4:10} 但这些（斐理伯人）虽然身在远处，却分担了他的苦难，既派人到他那里，又按自己的能力服事他，毫无遗漏。他说：\Quote{你们不但在现在，而且常常这样，在各方面帮助我。}所以，这是一种\OriginalTerm{共融}{fellowship}，是在福音事工上的共融。因为当一人宣讲，你们服事那宣讲者，你们就分享他的冠冕。即使在世俗的竞赛中，冠冕不仅归于那努力的人，也归于教练、陪同，以及所有帮助运动员备战的人。因为那些扶持他、使他恢复的人，完全有理由分享他的胜利。在战争中也是如此：不仅那赢得勇武之奖的人，所有随从他的人也完全有权利要求分享战利品并分享荣耀，因为他们通过随从参与了战斗。因为服事圣人不无小补，而是大有裨益。它使我们分沾为他们所预备的赏报。举例来说：假设某人为了天主放弃了大批财产，不断将自己奉献给天主，修习伟大德行，甚至在言行、思想、每件事上都极其严格。即使你没有表现出这样的严格，你也可以分享他为这些事所珍藏的赏报。如何做到？如果你在言语和行为上帮助他。如果你既满足他的需要，又尽一切可能服事他来鼓励他。因为那时，那条崎岖道路的平坦者就是你自己。因此，如果你景仰那些在旷野中实践天使般生活的人，以及教会中同样践行这些德行的人；如果你景仰他们，却为自己远远落后而悲伤，那么你可以通过另一种方式——服事他们、帮助他们——与他们共享。因为这正是天主的慈爱：将那些不太热心、不能承担那艰苦、崎岖而严格生活的人，我再说，也通过另一种方式，将他们带到与其他人相同的行列中。保禄所说的\Quote{共融}正是指此。他说：\Quote{他们在属血肉的事上与我们分享，我们在属神的事上与他们分享。}因为如果天主因微小而无价值的事物就赐予天国，那么祂的仆人也因微小而物质的事物，就在属神的事上给人分享；或者不如说，是祂藉着他们同时赐予两者。你不能守斋、不能独处、不能卧在地上、不能彻夜不眠？但如果你以另一种方式行此事，即服事那些在这些事上劳苦的人，不断地使他们得到休养和膏抹，减轻这些工作的痛苦，你仍可获得这一切的赏报。他站在那处战斗、挨打。你当在战斗结束时服事他，用双臂接住他，擦去汗水，使他舒畅；安慰、抚慰、恢复他疲惫的灵魂。如果我们愿意如此热心地服事圣人，我们必分沾他们的赏报。基督也这样告诉我们：\Quote{要藉着那\OriginalTerm{不义的钱财}{mammon of unrighteousness}结交朋友，使他们可以在永久的帐幕里接待你们。}\BibleRef{路 16:9} 你们看到他们已成为共享者吗？他说：\Quote{从第一天直到如今。}又说：\Quote{我欢喜}，不仅因过去的事，也因将来的事；因为从过去我就猜测那将来的事。\par

第六节。\BibleQuote{我深信那在你们中间开始了善工的，必予以完成，直到耶稣基督的日子。}\BibleRef{斐 1:6}\par

你看，他也教导他们要谦逊。因为既然他对他们作了一件伟大的见证，免得他们像人常有的那样自高，他就立刻教导他们要将过去和将来都归于基督。如何？他没有说：\Quote{我深信你们怎样开始，也必怎样完成}，而是说什么？\Quote{那在你们中间开始了善工的，必予完成。}他没有夺去他们的功劳（因为他说：\Quote{我因你们的共融而欢喜}，显然把这当作他们的行为），也没有称他们的善行全然是他们自己的，而主要是出于天主。他说：\Quote{因为我深信，那在你们中间开始了善工的，必予以完成，直到耶稣基督的日子。}那就是，天主必然完成。他暗示：我这样想，不是关乎你们自己，而是关乎你们的后裔。确实，天主在某人身上工作，这是不小的赞美。因为如果祂是\Quote{\OriginalTerm{不看人的情面}{no respecter of persons}}——事实上祂确实不看情面——而是在我们行善时，祂看我们的意向而帮助我们，那么显然我们是主动者，吸引祂到我们这里；所以即使从这个观点，他也没有夺去他们的赞美。因为如果祂的工作是不加分辨的，那么连外教人和所有人都有可能让祂在自己身上工作——即，如果祂像推动木头和石头那样推动我们，并不要求我们的参与。因此，当他说\Quote{天主必予完成}时，这又成了他们的赞美，因为他们吸引了天主的恩宠到自己身上，使祂帮助他们超越人性。在另一种意义上也是赞美，即\Quote{你们的善行是如此出色，不能出自人，而需要神圣的推动。}但如果天主必然完成，那么劳动也不会太多，他们应当有勇气，因为藉着祂的协助，他们必能轻易完成一切。\par

第七节。\BibleQuote{我这样思念你们众人是应当的，因为我把你们放在心上，无论我在锁链中，或在辩护与确定福音时，你们众人都与我同享恩宠。} \BibleRef{斐 1:7}\par

他在这里仍大大显示了他的渴望，因为他把你们放在心上；就是在狱中，虽被锁链捆绑，他仍记念斐理伯人。这对这些人来说，也不是小事，因为这位圣人并非因偏见而怀有爱，而是出于判断和正确的理由。因此，被保禄如此热切地去爱，是证明一个人是伟大和可钦佩的。\Quote{在辩护}，他说，\Quote{与确定福音时。}如果他即使在狱中都记念他们，那又有什么奇怪呢？他说，甚至在我前往法庭辩护的那一刻，你们也没有从我记忆中滑落。因为属神的爱是如此有王权，它不向任何时机让步，而是永远抓住爱者的灵魂，不容任何烦恼或痛苦压倒那灵魂。因为在巴比伦的火窑中，当那样大的火焰升起时，对那几位有福的儿童来说，它却是甘露。同样，占据爱主而又蒙主喜悦者灵魂的友谊，也能抖落每一火焰，产生奇妙的甘露。\par

\Quote{在确定福音时}，他说。这样，他的锁链就是福音的确定和辩护。确实如此。如何呢？因为如果他逃避锁链，他可能被视为骗子；但他忍受一切，包括锁链和苦难，显示他并非为了人的缘故而受苦，而是为了那赏报的天主。因为没有人愿意死，没有人愿意冒这样大的风险，没有人愿意与这样一位君王（我是指尼禄）冲突，除非他仰望另一位更大的君王。他的锁链实在是福音的\Quote{确定}。看他是如何超乎成功地将一切变为反面。因为他们以为软弱和贬损的，他却称为确定；若没有发生这事，那才是软弱。然后他表明他的爱并非出于偏见，而是出于判断。为什么呢？他说，我把你们放在心上，在我的锁链中，在我的辩护中，因为你们是\Quote{我恩宠的同分享者}。这是什么？使徒的\Quote{恩宠}就是被锁链捆绑，被驱逐，受千万苦难吗？是的。因为他说：\BibleQuote{我的恩宠够你用，因为我的德能在软弱中才全显出来。} \BibleRef{格后 12:9} 所以他说：\BibleQuote{我甘心情愿喜欢我的软弱、凌辱。} \BibleRef{格后 12:9} 既然我在你们的行动中看到你们德性的证明，并且你们同享这恩宠，而且欣然如此，我就合理地这样想。因为我曾考验过你们，极其深知你们和你们的善工：即使你们离我这样远，你们仍努力不在我们患难中缺席，而是为了福音的缘故，与我同受考验，分受我的劳苦，如同我亲临战场一般——你们既这样，我作见证，正是理所当然的。\par

为什么他不说\Quote{同享者}，而说\Quote{与我同享者}呢？他的意思是：我自己也与另一位同享，使我成为福音的同享者；就是说，我分享那为福音所存留的美好事物。而令人惊奇的正是他们都有这样的心志；因为他说\Quote{你们众人都同享恩宠}。从这些开端，我深信你们至终也都如此。因为如此光辉的开始，不可能熄灭而失败，它指向伟大的成果。\par

既然也有可能以其他方式同享恩宠、考验和磨难，我恳求你们，也让我们成为同享者。这里站着的众人，不，应该说所有的人，都愿意与保禄同享将来的美好事物！如果你们愿意，你们有能力这样做：为那些继承了他职务的人，当他们为基督的缘故受苦时，你们帮助他们，支援他们。你看见你的弟兄在考验中？伸出援手！你看见你的老师在争战中？站在他身边！但是，有人说，没有人像保禄！现在这是轻蔑！现在这是批评！那么没有人像保禄吗？好吧，我同意。但是，\BibleQuote{谁接纳一位先知，因他是先知，将领受先知的赏报。} \BibleRef{玛 10:41} 难道这些人得荣耀，是因为他们与\Emph{保禄}合作吗？不，是因为他们与承担了宣讲的人合作。保禄之所以可敬，是因为他为基督的缘故受了这些苦。\par

确实没有人像保禄。不，甚至没有一个人稍微接近那位有福的人。但宣讲与那时是一样的。\par

他们不仅在保禄的锁链上与他相通，更从起初就如此。请听他所说：\Quote{你们斐理伯人自己也知道，在福音初期，没有教会与我在给予和领受的事上互通，唯有你们。} \BibleRef{斐 4:15} 即使没有考验，老师也有许多劳苦：警醒、在言语上辛勤、教导、抱怨、控告、责难、嫉妒。忍受千万人的口舌，而自己只能独自忧虑，这岂是小事？唉！我该怎么办？因为我进退两难。我切愿激发并鼓励你们参与天主的圣徒们的联盟和援助；但我怕有人怀疑别的，以为我说这些不是为了你们，而是为了他们。但你们要知道，我说这些不是为了他们，而是为了你们自己。如果你们愿意听，我以我的话使你们信服：你们的收益与他们的收益并不相等。因为你们若施予，所施予的这些东西，无论愿意与否，不久之后必将离别，让位给别人；但你们所领受的却是伟大且远为丰盛的。或者，你们不这样认为，即施予时你们将领受？因为如果你们不这样认为，我甚至不希望你们施予。我绝非为他们做演讲！除非一个人首先这样认为：他宁可领受而非施予，他获得万倍，他受益而非施益，那就让他不要施予。如果像施惠于领受者那样，就让他不要施予。因为我更关心的不是圣徒们得到支持。因为即使你们不施予，别人也会施予。所以，我所愿的是你们能从自己的罪过中得到缓解。但不这样施予的人将得不到缓解。因为施舍不在于给予，而在于乐意地给予；在于欢喜，在于对领受者心存感激。因为他说：\Quote{不要勉强，也不要出于被迫，因为天主喜爱乐意捐献的人。} \BibleRef{格后 9:7} 既然如此，除非一个人这样施予，就让他不要施予：因为那是损失，不是施舍。如果你们知道你们将获益，而不是他们，要知道你们的获益会更大。因为至于他们，身体得到饱饫，但你们的灵魂得到赞许；至于他们，当他们领受时，他们的罪过没有一个获得赦免，但至于你们，你们过犯的绝大部分被除掉。那么，让我们在他们的巨大奖赏中与他们有份。当人们立君王时，他们不认为他们给予的多于领受的。你立基督为王，你将获得极大的保障。你也要与保禄有份吗？我为什么说保禄，而领受的却是基督呢？\par

为使你们知道，我所说所行的一切都是为了你们，并非为了顾惜他人的安逸，若教会中有任何掌权者生活富裕，一无所缺，即使他是位圣人，也不要施与，却要优先周济那贫困的，即使他并不可如此称赞。为何如此？因为基督也这样愿意，就如祂所说：\BibleQuote{你设午宴或晚宴时，不要请你的朋友，也不要请你的亲戚，而应请那些瘸腿的、残废的、瞎眼的，因为他们无力回报你。}\BibleRef{路 14:12}因为人不应不加区别地施予这样的关注，而是给饥饿的、口渴的、需要衣服的、外乡人，以及那些由富足沦为贫穷的人。因为祂并非简单地说\BibleQuote{我得了喂养}，而是说\BibleQuote{我饿了}，因为祂说：\BibleQuote{你们看见我饿了，便给我吃。}\BibleRef{玛 25:35}这要求是双重的：他既是圣人，又是饥饿的。因为如果仅仅饥饿的人就应得喂养，那么一个饥饿的圣人更应如此。所以如果他是圣人，但并不缺乏，就不要施与；因为这样没有益处。因为基督并没有这样命令；或者更确切地说，那个富有且接受施与的人并不是圣人。你们可见我说这些并非为贪图不义之财，而是为你们的益处？喂养饥饿的人，免得你喂养地狱之火。他吃你的东西，也圣化所剩余的。\BibleRef{路 11:41}试想那寡妇如何供养了厄里亚；她喂养别人并未超过自己被喂养：她付出并未超过所得。如今这事在更伟大的事上同样发生。因为并非\BibleQuote{一桶面粉,}，也不是\BibleQuote{一瓶油}\BibleRef{列上 17:14}，而是什么？\BibleQuote{百倍，并永生}\BibleRef{玛 19:21}，是这事的报偿——你成为天主的慈悲，属神的食物，纯洁的酵母。她是个寡妇，饥荒逼人，但这些事没有一样阻止她。她还有孩子，即便如此也没有被拦阻。\BibleRef{列上 17:12}这妇女成了与那投了两个小钱的妇女同等的人。她并没有对自己说：\Quote{我从这个人能得到什么？他需要我。如果他有什么能力，就不会挨饿，就不会使干旱停止，就不会遭受同样的苦难。或许他也得罪了天主。}这些事她一样也没有想。你们可见，以纯朴行善，不过分查究受惠者，是多么大的善事吗？如果她选择查究，她就会怀疑；她就不会相信。同样，亚巴郎如果选择查究，也不会接待天使。因为一个在这些事上过于挑剔的人，不可能，实在不可能遇到他们。不，这样的人通常只会撞上骗子；这是怎么回事，我来告诉你们。虔诚的人并不想显得虔诚，也不披上虚伪的外衣，因此可能被拒绝。但骗子以此为业，他装出大量难以看穿的虔诚。所以，那即使向看似不虔诚的人行善的，会遇到真正虔诚的人；而那专门寻找那些被认为是虔诚的人的，却常常遇到并非如此的人。因此，我恳求你们，让我们在一切事上都纯朴行事。即使我们假设来者是个骗子，你也不奉命查究这事。因为祂说：\BibleQuote{求你的，就给他}\BibleRef{路 6:30}；以及\BibleQuote{不可不赎回那将被杀的人}\BibleRef{箴 24:11}。然而大多数被杀的人是因为被定罪为邪恶而受此苦；祂仍然说：\BibleQuote{不可不}。因为在这事上我们将肖似天主，这样我们将受称赞，并获得那些不朽的祝福，愿我们所有人都被堪当承受，因我们的主耶稣基督的恩宠和慈爱，祂与圣父及圣神，共享光荣、权能、尊威，从今世直到永远，世世无尽。阿们。\par

\SourceRecord{Translated by John A. Broadus.；Nicene and Post-Nicene Fathers, First Series；Vol. 13.；Edited by Philip Schaff.；Buffalo, NY: Christian Literature Publishing Co.,；1889.；<http://www.newadvent.org/fathers/230201.htm>.}

% structure: p0001 source=h1 target=section reason=page-title

\section{斐理伯书讲道词 第二篇}

\BibleRef{斐 1:8-11}\par

\BibleQuote{天主为我作证：我是怎样以耶稣基督的慈爱，切望你们众人。我祈祷的，就是要你们的爱德，在知识和各种明辨上，不断增加，使你们能辨别最优良的事，好使你们在基督的日子，是纯洁无瑕的，并藉耶稣基督，充满义德的果实，为光荣赞美天主。}\par

他并非如同那该被怀疑的人一样呼求天主作证，而是出于他的深情厚爱，以及他极度的说服和信心；因为说过他们与他共融之后，他又加上这句：\Quote{以基督的慈爱}，免得他们认为他思念他们是为了这个原因，而非单纯为了他们本身的缘故。这话是什么意思：\Quote{以基督的慈爱}？它相当于\Quote{按照基督}。因为你们是信徒，因为你们爱基督，因为那按照基督的爱。他没有说\Quote{爱}，而是用了更温暖的说法：\Quote{基督的慈爱}，好像是说：\Quote{藉着在基督内的关系，成了你们的父亲}。因为这赋予我们温暖而炽热的心肠。祂赐予祂忠仆这样的心肠。\Quote{以这样的心肠}，祂说，好像有人说：\Quote{我用不是常人的心肠爱你们，而是用更温暖的心肠，就是基督的心肠}。\Quote{我是多么思念你们众人}。我思念所有的人，因为你们都有这样的本性；我无法用言语向你们描述我的思念；所以是无可形容的。因此，我把它留给那洞察人心的天主去知道。如果他是在奉承他们，就不会呼求天主作证，因为这样做不可能没有危险。\par

第九节\BibleRef{斐 1:9}。他说：\Quote{我为此祈祷，愿你们的爱德不断增加。}因为这是一种不会满足的美善；你看，他这样被爱，还渴望被爱得更多，因为爱其所爱之人，不愿停留在任何爱的限度内，因为如此高贵的事物不可能有尺度。保禄希望爱的债常欠着；\Quote{除了彼此相爱，你们不要欠人什么。}\BibleRef{罗 13:8}爱的尺度就是永不止步；他说：\Quote{愿你们的爱德，}他说，\Quote{越来越多。}想想这话的特色：他说：\Quote{愿它越来越多，}他说，\Quote{在知识和各种明辨上。}他并非一味赞扬友谊或爱，而是赞扬那出自知识的爱；也就是说，你不可用同样的爱对待所有人：因为这不是出于爱，而是出于麻木。他说的\Quote{在知识上}是什么意思？意思是：有判断、有理性、有辨别力。有些人没有理性地爱，简单地、随便地爱，所以这种友谊是脆弱的。他说：\Quote{在知识和各种明辨上，使你们能辨别最优良的事，}即有利的事。我说这话不是为我自己的缘故，他说，而是为你们，因为恐怕有人因异端者的爱而败坏；他暗示了这一切，看看他是怎么引入的。不是为我自己的缘故，他说，我说这话，而是为使你们纯洁无瑕，就是说，你们不可在爱的幌子下接受虚假的教义。那么，他说，\Quote{如果可能，尽力与众人和睦相处}？\Quote{和睦相处}\BibleRef{罗 12:18}，他说，不是爱到被那友谊所伤害；因为他说：\Quote{如果你的右眼使你跌倒，剜出它来，扔掉它；为使你们纯洁无瑕}\BibleRef{玛 5:29}，即在天主面前，\Quote{并无可指摘}，即在人面前，因为许多人的友谊常常成为他们的伤害。即使它不伤害你，他说，但别人仍可能因而跌倒。\Quote{直到基督的日子}；也就是说，那时你们被发现是纯洁的，没有使任何人跌倒。\par

第十一节\BibleRef{斐 1:11}。\Quote{藉耶稣基督，充满义德的果实，为光荣赞美天主；}即持有纯正教义，并度正直的生活。\par

不仅是正直，而且是\Quote{充满义德的果实}。因为确实有一种不按照基督的义，例如道德生活。\Quote{这些是藉耶稣基督，为光荣赞美天主。}你看见我不是讲我自己的光荣，而是天主的义；并且他常常称怜悯本身为义；他说，不要让你们的爱德间接伤害你们，阻碍你们对有益之事的知觉，并且谨慎，免得因对某人的爱而跌倒。因为我确实希望你们的爱德增加，但不要因此受到伤害。我也不希望它是出于偏见，而是出于对我所说的好坏之验证。他不是说你们要接受我的意见，而是你们要\Quote{验证}它。他没有直说，不要与这个人或那个人结交，而是说，我希望你们的爱德要顾及有益之事，而不是你们没有见识。因为如果你们不为了基督并藉着基督行义，那是愚蠢的。注意这句话：\Quote{藉着祂。}那么他把天主当作一个工具吗？绝无此念。他说，不是为要我得到赞美，而是为天主受光荣。\par

第十二、十三节。\BibleQuote{弟兄们，我愿意你们知道，我所遭遇的事，反而使福音更蓬勃，以致我在基督内的锁链，在全御营\OriginalTerm{御营}{prætorian guard}和其余众人中，都彰显出来。}\par

他们很可能在听说他被捆绑时感到悲伤，并以为宣讲工作停顿了。那么如何？他立刻消除了这个怀疑。这也显示了他的关爱，因为他宣告了发生在他身上的事，因为他们焦虑了。你怎么说？你被捆绑！你受阻碍！那么福音如何进展？他回答说：\BibleQuote{以致我的锁链在总督府中显露出来，都是为基督的缘故。}这件事非但没有使其他人沉默，也没有恐吓他们，反而鼓励了他们。如果那些接近危险的人不仅没有受害，反而获得了更大的勇气，那么你们更应该如此。假如他被捆绑时感到难过并保持沉默，他们很可能也会同样受影响。但既然他在被捆绑时更加勇敢地说话，他就给了他们比没有被捆绑时更大的信心。他的锁链如何\Quote{反而促进了福音的进展}？意思是，天主在他的安排中如此安排，我的锁链并非隐藏，这些锁链是\BibleQuote{在}基督内的，是\BibleQuote{为}基督的。\par

\BibleQuote{在整个总督府。}因为直到那时他们是这样称呼宫殿的。他说：\Quote{也在全城。}\par

第14节。\BibleQuote{并且大多数弟兄在主内，因我的锁链而获得信心，更加勇敢地宣讲天主的圣言，无所畏惧。}\BibleRef{斐 1:14}\par

这表明他们以前就已勇敢并放胆宣讲，但现在更甚。他说：如果别人因我的锁链而获得勇气，那么我更应如此；如果我是他人信心的原因，那么对我自己更是如此。\BibleQuote{并且大多数弟兄在主内。}因为说\Quote{我的锁链给了他们信心}是一件大事，所以他预先加上\BibleQuote{在主内}。你看，即使当他被迫说大事时，他也没有偏离节制。他说：\BibleQuote{更加勇敢地宣讲天主的圣言，无所畏惧}；\BibleQuote{更加}一词表明他们已经开始了。\par

第15节。\BibleQuote{有些人出于嫉妒和纷争宣讲基督，另一些人则出于善意。}\BibleRef{斐 1:15}\par

这意思值得探究。因为保禄被拘禁时，许多不信者为了激起皇帝更猛烈的迫害，自己也宣讲基督，目的是使福音的传播加剧皇帝的愤怒，并使所有怒火落在保禄头上。因此，从我的锁链产生了两种行动。一些人因此大得勇气；另一些人则希望导致我的毁灭而开始宣讲基督；\BibleQuote{他们中有些出于嫉妒}，即嫉妒我的名声和坚毅，并渴望我的毁灭及争斗精神，与我一同工作；或者他们自己受人尊敬，并期望从我荣耀中分得一些。\BibleQuote{另一些人出于善意}，即没有虚伪，全心诚意。\par

第16节。\BibleQuote{这些人出于分党，不真诚地宣讲基督。}\BibleRef{斐 1:16}\par

也就是说，不是出于纯洁的动机，也不是出于对事情本身的尊重；而是为什么呢？\BibleQuote{以为会给我的锁链增添苦难。}因为他们以为我会因此而陷入更大的危险，他们就苦难上加苦难。啊，残忍！啊，魔鬼的煽动！他们看见他被捆绑并投入监狱，仍然嫉妒他。他们想增加他的灾难，使他遭受更大的愤怒；他说的好：\BibleQuote{以为}，因为结果并非如此。他们确实想以此使我忧伤；但我却因福音被推进而欢喜。\par

第17节。\BibleQuote{但那些人出于爱心，知道我是为辩护福音而设立的。}\BibleRef{斐 1:17}\par

\BibleQuote{为辩护福音而设立}是什么意思？就是他们在为我必须向天主交账做准备，并且协助我。\par

\BibleQuote{为辩护}是什么意思？我被委派去宣讲，我必须交账，并为我被委派的工作答辩；他们帮助我，使我的辩护容易；因为如果发现许多人已被训导并相信了，我的辩护就会容易。因此，可能出于不好的动机去做善事。不仅这样的行为没有赏报，反而有惩罚。因为他们出于使基督的宣讲者陷入更大危险的欲望而宣讲基督，这不仅不会得到赏报，反而要受惩罚。\BibleQuote{有些人出于爱心。}也就是说，他们知道我必须为福音交账。\par

第18节。\BibleQuote{那么如何？无论如何，无论出于假善还是出于真诚，基督总被宣讲。}\BibleRef{斐 1:18}\par

但请看这人的智慧。他没有严厉控告他们，而是提到了结果；他说：无论这样做还是那样做，对我有什么差别？无论如何，\BibleQuote{无论出于假善还是出于真诚，基督总被宣讲。}他没有说\BibleQuote{让祂被宣讲}，正如有些人认为他那样说，为异端开放了道路；而是\BibleQuote{祂被宣讲}。首先，他并非制定法律，好像制定法律一般说\BibleQuote{让祂被宣讲}，而是报告正在发生的事；其次，即使他好像制定法律一样说话，这样也不会为异端开放道路。\par

我们来考查这事。因为即使他允许他们按照他所宣讲的去宣讲，也并没有为异端开路。何以见得呢？因为他们宣讲的是健康的教义；虽然他们行事的动机和目的败坏了，但宣讲本身并未改变，而他们被迫如此宣讲。为什么呢？因为，如果他们宣讲与保禄不同的东西，教导与他不同的道理，就不会增加皇帝的忿怒。但现在，通过推进他的宣讲，以同样的方式教导，并像他那样收徒，他们能够激怒皇帝，当他看到门徒众多时。但后来某个邪恶而愚昧的人抓住这段经文说：如果他们想惹恼他，他们本当反其道而行之，即赶走已经相信的人，而不是让信徒增多。我们该如何回答？那就是他们只关注这一件事：如何使他陷入眼前的危险，不给他逃脱的机会；因此他们认为这样能使他伤心，并扑灭福音，而不是用另一种方式。\par

如果采取另一种方式，他们就会平息皇帝的怒气，让他自由再去宣讲；但通过这种方式，他们认为只要除掉他，一切都会因他而毁灭。然而，大多数人不会有这种意图，只有某些心怀恶意的人才有。\par

然后他说：\Quote{在此事上，}他说，\Quote{我欢喜，而且要更加欢喜。}\Quote{是的，我要欢喜}是什么意思？即使这件事更加如此，他也说。因为他们甚至违背自己的意愿与我合作；他们会因自己的劳苦而受惩罚，而我，丝毫没有贡献，却要得奖赏。魔鬼的恶毒还有比这更甚的吗？它竟使宣讲受惩罚，使劳苦遭报应？你们看到他用多少恶事刺透自己的手下！一个仇恨和敌视他们得救的人怎能安排这一切？你们看到那与真理作战的，如何毫无能力，反而伤害自己，如同用脚踢刺一样吗？\par

第19节。\BibleQuote{因为我知道，}他说，\BibleQuote{这事藉着你们的祈祷和耶稣基督之圣神的供应，终必使我得救。}\BibleRef{斐 1:19}\par

没有什么比魔鬼更恶毒了。它就这样处处使它自己的人陷入无益的劳苦，并撕裂他们。它不仅不让他们获得奖赏，甚至使他们受惩罚。\par

因为它不仅命令他们宣讲福音，还同样命令他们守斋和守童贞，方式不仅使他们失去赏报，而且使那些遵行此道的人遭受重祸。关于这些人，他在别处也说：\BibleQuote{良心被烙过了一般。}\BibleRef{弟前 4:2}\par

为此，我恳求你们，事事感谢天主，因为他既减轻了我们的劳苦，又增加了我们的赏报。因为那些在他们中间过童贞生活的人，得不到在我们中间在婚配中贞洁生活的人所得的赏报；而在异端者中过童贞生活的人，则与奸淫者同受谴责。这一切都源于他们没有以正确的意向行事，反而指责天主的受造物和他不可言喻的智慧。\par

那么，我们不要懈怠。天主在我们面前安排了适度的竞赛，没有劳苦。但我们不要因此而轻视它们。因为如果异端者把自己投入到无益的劳苦中，我们若不愿意忍受那些更轻且有更大赏报的劳苦，还有什么借口呢？基督的诫命中哪一条是沉重的？哪一条是令人难以忍受的？你不能过童贞生活吗？你被允许结婚。你不能舍弃一切所有吗？你被允许用你所有的来供应别人的需要。\BibleQuote{你们的富裕补足他们的缺乏。}\BibleRef{格后 8:14}这些事确实显得沉重。哪些事？我是指轻视金钱，克服肉身的欲望。但他的其他诫命不需要代价，不需要强制。因为告诉我，不说坏话，仅仅不诽谤，有什么强制？不嫉妒他人的财物，有什么强制？不被虚荣所支配，有什么强制？忍受折磨并忍耐，这是力量的表现。实践哲学是力量的表现。一生忍受贫困是力量的表现。与饥渴搏斗是力量的表现。当这些都没有，而你可以像一个基督徒那样享受自己的所有，不嫉妒别人，又有什么强制呢？\par

嫉妒由此而生；不，一切邪恶无不由此而生，即我们依恋现世的事物。因为如果你把金钱和今世的荣耀视为虚无，你就不会对拥有它们的人投以恶眼。但既然你渴慕这些，崇拜它们，受它们奉承，为此嫉妒困扰你，虚荣困扰你；这一切都源于崇拜现世生活的事物。你嫉妒别人富有吗？那样的人倒是可悲可泣的对象。但你却发笑，立刻答道：我是可悲的对象，不是他！你也是可悲的对象，不是因为你贫穷，而是因为你自以为不幸。因为我们哭泣那些明明没什么事却不满的人，不是因为他们有什么问题，而是因为他们没有问题却以为有问题。例如：如果一个退了烧的人仍然焦躁不安，在床上辗转反侧，虽然健康地躺着，他不是比发烧的人更值得哭泣吗？不是因为他发烧（他并不发烧），而是因为没有病却仍以为有病。你之所以可悲，正是因为你自以为不幸，而不是因为你的贫穷。为你的贫穷，你倒应该被认为是幸福的。\par

你为甚么嫉妒那富人？是因为他使自己陷于许多挂虑中吗？陷于更艰苦的奴役中吗？因为他像一条狗，被千万条锁链捆绑——那就是他的财富吗？黄昏降临他，黑夜降临他，但休息的时刻对他来说却是烦恼、痛苦、焦虑的时候。有响声：他立刻跳起来。他的邻人遭劫了吗？那什么也没失去的人比失主更关心这事。因为那人失去了一次，但忍受了痛苦后便放下了挂虑；而另一个却总是带着它。黑夜来了，这我们灾祸的港湾，我们忧伤的慰藉，我们伤口的良药。因为那些被过度的悲伤压垮的人，常常对朋友、亲戚、密友充耳不闻——甚至对一位想要安慰的父亲也是如此，他们反而曲解他的话；可是当睡眠命令他们休息时，无人能直视他的脸。因为任何燃烧都比不上悲伤的苦涩折磨我们的灵魂。如同身体被日光的暴力煎熬折磨，被带到一处有许多泉源和微风轻拂的客栈，同样，黑夜将我们的灵魂交托给睡眠。是的，我更应该说，这不是黑夜或睡眠所做的，而是那认识我们劳苦世代的\Person{天主}所做的，而我们对自己毫无怜悯，反而好像与自己为敌，设计了一种比自然的休息匮乏更强大的暴政——因财富而来的失眠。因为经上说：\Quote{财富的焦虑赶走睡眠。}\BibleRef{德 31:1}看，\Person{天主}的眷顾何其大！但他没有将休息交托给我们的意志，也没有将我们对睡眠的需要交托给选择，而是将它绑在自然的必需中，好使益处甚至违背我们的意愿而临到我们。因为睡眠是自然的。但我们作为伟大的自恨者，像别人的敌人和迫害者，设计了一种比这自然必需更强大的暴政，即因钱财而来的暴政。天亮了？那么某人害怕告发者。黑夜笼罩他？他害怕盗贼。死亡临近？他必须将财物留给别人的念头比死亡更折磨他。他有一个儿子吗？他的欲望增加了；然后他幻想自己贫穷。他没有儿子吗？他的痛苦更大。你以为那从任何地方都无法得到快乐的人有福吗？你能嫉妒这样颠簸的人，而你自己却安息在贫穷的平静港湾中吗？事实上，这正是人性的不完美：它不能高尚地承受自己的好处，反而侮辱自己的繁荣。\par

这一切都在世上；但当我们离开那里时，请听那拥有无数财物的富人——如你所说（因为就我而言，我不称这些为好东西，而是无关紧要的）——请听这位拥有无数财物的主说什么，他需要什么：\Quote{亚巴郎父啊！他喊道，请打发\Person{拉匝禄}，用他的指尖蘸点水来凉润我的舌头，因为我在这火焰中非常痛苦。}因为即使那富人没有忍受我所提到任何一件事，如果他的一生都在恐惧和挂虑之外度过——我为什么说他的一生？倒是那一瞬间（因为一生只是一瞬间，与那无终的永恒相比）——如果一切事情都按他的愿望成就，难道他不该为这些话，不，为这种状况而受怜悯吗？你的桌子不是曾经酒流成河吗？现在你连一滴水也掌管不了，而且是在你最大的需要中。你不是曾忽视那满身疮痍的穷人吗？现在你请求见他一面，却没有人允许。他曾躺在你的门口；但现在在\Person{亚巴郎}的怀抱中。你当时躺在你高耸的天花板下；但现在在地狱的火中。\par

让这些富人们听吧。是的，更确切地说，不是富人，而是无怜悯心的人。因为他受惩罚不是因他富有，而是因他没有显示怜悯；因为一个既富有又有怜悯心的人，是可能得到一切好处的。因此，那富人的目光不落在别人身上，只落在当时向他求乞的那人身上，使他可以从过去行为的记忆中知道，他的惩罚是公义的。难道没有千万个义人中的穷人吗？但那当时躺在他门口的穷人单独被他看见，为教训他和我们，信赖财富是多么大的善行。他的贫穷没有阻碍那穷人获得天国；他的财富没有帮助那富人逃避地狱。人在什么点上算贫穷？人在什么点上沦落到乞讨？他不是，不是贫穷的，那什么也没有的人，而是那贪求许多东西的人！他不是富有的，那拥有大财产的人，而是那什么也不缺少的人。因为拥有全世界，却比那什么也没有的人活得更沮丧，有什么益处呢？他们的性情使人富有或贫穷，而非钱财的多少或缺少。你，一个穷人，想要变得富有吗？你可随意，无人能拦阻你。轻视世界的财富，视之为无，因为它本就是无。弃绝财富的欲望，你立刻就成了富人。那不想变得富有的人就是富人；那不愿贫穷的人就是穷人。如同那生病的人，即使在健康中自叹自怜，而不是那比完全健康更轻松地忍受疾病的人，同样，那不能忍受贫穷的人，在财富中自认比穷人更穷，就是穷人；而那比富人忍受财富更轻松地忍受贫穷的人，就不是穷人，因为他更富有。\par

请告诉我，你为甚么害怕贫穷？你为甚么战栗？不是由于饥饿吗？不是由于干渴吗？不是由于寒冷吗？难道不是为了这些吗？没有，没有一个在这事上曾经匮乏！\BibleQuote{因为请查看古代的各世代，有谁信赖了主，而遭遗弃？或有谁盼望了他，而蒙受羞辱？}\BibleRef{德 2:11}\par

再者，\BibleQuote{你们看那天上的飞鸟，也不种，也不收，也不积蓄在仓里，你们的天父尚且养活它们。}\BibleRef{玛 6:26} 没有人能轻易指出有谁因饥饿寒冷而死亡。那么你为什么惧怕贫穷呢？你说不出来。因为如果你有足够的必需品，为什么还要惧怕它呢？是因为你没有众多的仆人吗？这实在是摆脱了主人，这是持续的快乐，这是无忧无虑。是因为你的器皿、床榻、家具不是银子做的吗？拥有这些东西的人，比你有更大的享受吗？一点也没有。无论是什么材料，用途都是一样的。是因为你不是众人畏惧的对象吗？愿你永远不要如此！别人战战兢兢敬畏你，这有什么乐趣呢？是因为你害怕别人吗？但你无须惊慌。因为\BibleQuote{你愿意不惧怕掌权者吗？你只管行善，就必从它得称赞。}\BibleRef{罗 13:3} 有人会说，这是因为我们受人轻视，容易受伤害吗？造成这情况的不是贫穷，而是邪恶；因为许多穷人平静地度过了一生，而统治者、富人、有权势者，其结局却比一切作恶者、强盗、盗墓者更为悲惨。因为在你身上贫穷所带来的，在他们身上财富也会带来。因为那些想虐待你的人，因你卑贱的地位而做的事，他们却因嫉妒和恶眼而对富人做，后者比前者更甚，因为这是更强烈的虐待他人的欲望。嫉妒者竭尽全力行事，而轻视者却往往甚至可怜被轻视的人；他的贫穷和完全无力，常常成为他得救的原因。\par

并且有时候，我们对他说：\Quote{如果你除掉这样一个人，那将是多大的功绩！如果你杀死一个穷人，你将得到何等巨大的好处！}我们这样或许能平息他的怒气。但嫉妒却针对富人，不达目的不罢休，直到倾泻它的毒液。你们看，贫穷和财富本身都不是好的，而是我们自己的心性。让我们把它调整到良好的状态，用真智慧来训练它。如果这心性端正，财富不能将我们逐出天国，贫穷也不会使我们欠缺。反之，我们将温顺地忍受贫穷，享受未来美好事物时不会受到任何损失，甚至连在地上也是如此。而且我们将既享受地上的美好事物，也获得天上的美好事物——愿我们众人能得到这一切，赖天主的恩宠和慈爱，等。\par

\SourceRecord{Translated by John A. Broadus.；Nicene and Post-Nicene Fathers, First Series；Vol. 13.；Edited by Philip Schaff.；Buffalo, NY: Christian Literature Publishing Co.,；1889.；<http://www.newadvent.org/fathers/230202.htm>.}

% structure: p0001 source=h1 target=section reason=page-title

\section{斐理伯书讲道辞 第三篇}

\BibleRef{斐 1:18}\par

\BibleQuote{并且我在此欢喜，而且还要欢喜；因为我知道，这事借着你们的祈祷，以及耶稣基督的圣神的供应，终必使我得救，照着我热切的期待和盼望，我总不会蒙羞，反而要凡事放胆，无论是生是死，基督总在我身上显大。}\par

今世中所有令人痛苦的事都不能咬住那真正 \OriginalTerm{哲学的}{philosophic} 的崇高灵魂：无论是仇恨、控告、毁谤、危险还是阴谋。它仿佛逃向一座坚固的堡垒，在那里安全地抵御来自这低级尘世的一切攻击。保禄的灵魂就是这样；它占据了一个比任何堡垒更高的位置，即灵性智慧的所在，也就是真正的 \OriginalTerm{哲学}{philosophy}。因为那些外邦人（即异教徒）的所谓智慧只不过是空谈和儿戏。但我们现在不谈这些，而是论及保禄的事。那位有福者既有皇帝为敌，此外还有许多敌人以多种方式苦待他，甚至用恶毒的毁谤。而他说什么呢？我不仅不因这些事忧伤或消沉，反而 \BibleQuote{我还要欢喜，并且还要欢喜}，不是一时，而是我要永远为这些事欢喜。\BibleQuote{因为我知道，这事终必使我得救}，即未来的得救，甚至他们对我怀有的敌意和嫉妒也促进了福音。\BibleQuote{借着你们的祈祷}，他补充说，\BibleQuote{以及耶稣基督的圣神的供应，照着我热切的期待和盼望}。请看这位有福者的谦卑；他正在竞技中奋力，他已接近桂冠，他成就了万千伟业，因为他是保禄——对此还能添加什么呢？但他仍写信给斐理伯人：我可能因 \BibleQuote{你们的祈祷} 而得救——我，一个藉着无数成就获得救恩的人。\BibleQuote{以及耶稣基督的圣神的供应}，他说。\OriginalTerm{供应}{supply} 一词的意思是这样的：如果他（保禄）被认为配得你们的祈祷，那么他也将被认为配得更丰富的恩宠。因为 \Quote{供应} 的含义就是：如果圣神被供应给我，即更大量地赐予我。或者他在谈论解脱，\BibleQuote{直到得救}；也就是说，我也会像逃脱之前的危险一样逃脱当前的危险。关于同一件事，他说：\BibleQuote{在我初次申辩时，没有人站到我身边，众人都离弃了我；愿这罪不归于他们。但主站在我旁边，加给了我力量。} \BibleRef{弟后 4:16} 他现在预言的就是这事：\BibleQuote{借着你们的祈祷和耶稣基督的圣神的供应，照着我热切的期待和盼望}，因为我就是这样盼望的。为了使我们不把一切全托付给为我们的祈祷，而自己毫不贡献，请看他是如何提出自己那一部分的，即盼望——这位先知说，这是所有好事的根源。\BibleQuote{上主，求祢的仁慈临于我们，根据我们向祢所怀的盼望。} \BibleRef{咏 32:22} 正如另一处所载：\BibleQuote{你们且观察万民自古以来的世代；有谁信赖了上主，而受了羞辱？} \BibleRef{德 2:10} 还有，同一位有福者说：\BibleQuote{盼望不会使人蒙羞。} \BibleRef{罗 5:5} 这就是保禄的盼望——他盼望自己无论如何都不会蒙羞。\par

\BibleQuote{照着我热切的期待和盼望}，他说，\BibleQuote{叫我总不会蒙羞。} 你看到盼望天主是一件何等伟大的事吗？他说，无论发生什么事，我决不会蒙羞，也就是说，他们不能胜过我，\BibleQuote{反而要凡事放胆，如同以往一样，现在也照样，基督在我身上显大。} 他们确实指望在这圈套中捉住保禄，扑灭福音的宣讲，好像他们的诡计有任何力量似的。他说，这事不会发生，我现在不会死，反而 \BibleQuote{如同以往一样，现在也照样，基督在我身上显大。} 怎么会这样呢？我常陷入危险，那时众人都放弃了我，而且连我自己也放弃了。因为 \BibleQuote{我们心中已断定是必死的} \BibleRef{格后 1:9}，但主从这一切中解救我，所以现在他也要在我身上显大。那么怎么样？免得有人猜想并说：如果你死了，他难道不还是显大吗？是的，他回答，我知道他仍会显大；为此我没有说只有我的生命会使他显大，而且我的死亡也会。目前他是指 \Quote{借着生命}；他们不会毁灭我；即使他们这样做了，基督也仍然这样显大。怎么显大呢？借着生命，因为他解救了我；借着死亡，因为连死亡本身也不能说服我否认他，因为他给了我这样的准备，使我比死亡更坚强。一方面因为他使我脱离危险；另一方面，因为他使我不惧怕死亡的暴政：这样，他就借着生命和死亡显大。他说这话，不是因为他将要死，而是为了在他死的时候，他们不致像常人那样受影响。\par

但为叫你们知道，他的这些话并非指即将来临的死亡（这正是他们最感痛苦的想法），请看他是如何几乎以这样的话来缓解这痛苦：我说这些话，并非作为即将死去的人；因此他随即补充说：\Quote{并且我深信，我必存留，且与你们众人一同存留。}\Quote{我决不致于蒙羞，}他说，\Quote{在任何事上。}意思是，死亡不会给我带来羞耻，反而带来极大的益处。为什么？因为我不是不朽的，但我会比不朽者更加闪耀；因为一个不朽的人和一个会死的人轻视死亡，并非同一回事；因此，即使立即死亡也不会使我蒙羞，而我并不会死；\Quote{在任何事上我决不蒙羞，}无论生死。因为无论生死，我都会高尚地承受。他说得好！这是基督徒灵魂的特质！但他又加上：\Quote{全然放胆。}你们看到我如何完全脱离了羞耻吗？因为如果对死亡的恐惧削减了我的胆量，那么死亡就值得羞耻；但如果死亡来临时没有使我恐惧，那就没有羞耻；无论通过生我不会蒙羞，因为我仍宣讲那宣讲；还是通过死我不会蒙羞，恐惧没有阻拦我，因为我仍显出同样的胆量。当我提到我的锁链时，不要以为这事情可耻；它给我带来了如此多的益处，甚至给了别人信心。因为为基督被捆绑，不是羞耻；但出于对捆绑的恐惧而背叛任何属于基督的事，这才是羞耻。当没有这种事时，捆绑甚至成为胆量的缘由。但我多次脱离危险，并以此向外邦人夸耀，你们不要立即以为我蒙羞了，如果现在结果相反的话。这两种情况同样给予你们胆量。注意他如何以自身为例提出这一点，这在多处出现，如在\BibleBook{罗马}{书}中：\Quote{我不以福音为耻。}\BibleRef{罗 1:16}又如在\BibleBook{格林多}{前书}中：\Quote{这些事我以比方转用于自己和阿颇罗。}\BibleRef{格前 4:6}——\Quote{或借生或借死}：他说这话并非出于无知（因为他知道自己那时不会死，而是稍后）；但现在他仍准备他们的灵魂。\par

第21节。\Quote{因为对我来说，}他说，\Quote{活着是基督，死亡是收益。}\BibleRef{斐 1:21}\par

因为即使在死亡中，他的意思是，我并没有死，因为我在自身内拥有我的生命；那么，他们若是有能力通过这种恐惧将信德从我灵魂中赶走，就真的杀死我了。但只要基督与我同在，即使死亡追上了我，我仍然活着，而在现世生活中，不是这个，而是基督是我的生命。既然如此，即使在现世中也不是这样，而是\BibleQuote{但我现今在肉身内所活的生命，是活在信德内}；所以我也说在那种状态中，\BibleQuote{我活着，却不再是我，而是基督在我内活着}。\BibleRef{迦 2:20}基督徒应该如此！他说，我不活普通的生活。那么，蒙福的\Person{保禄}啊，你怎样生活呢？难道你不看见太阳，不呼吸空气吗？你不是和别人吃同样的食物吗？你不像我们一样踏足大地吗？你不需要睡眠、衣服、鞋子吗？你说\Quote{我不活了}是什么意思？你怎么不死？为何你自夸？这里没有自夸。因为如果事实不给他作证，人还可以有几分称之为自夸；但如果事实作证，哪里还有自夸呢？那么让我们学习他如何不活了，因为他在另一处说：\BibleQuote{我已与世界同钉十字架，世界也与我同钉十字架}。\BibleRef{迦 6:14}再听他说：\BibleQuote{我不再活了}。以及他说：\BibleQuote{对我来说，活着就是基督}。\Quote{生命}这个词意义重大，亲爱的诸位，同样\Quote{死亡}这个词也是。有身体的生命，有罪的生命，正如他在别处说：\BibleQuote{我们这些已死于罪恶的人，如何还能在罪恶中生活呢？}\BibleRef{罗 6:2}由此可见，是可能活罪的生命。请你们留心留意，我恳求你们，免得我的劳力白费。有永生和不朽的生命；有永恒的天上生命；\BibleQuote{因为我们的家乡是在天上}，他说，\BibleRef{斐 3:20}还有他所说的身体的生命：\BibleQuote{因为我们生活、行动、存在，都在他内}。\BibleRef{宗 17:28}他并不否认自己过着自然的生活，而是过着罪的生活，那是所有人所过的。那些不贪求现世生活的人，如何过着现世的生活？那些急趋另一种生命的人，如何过着这生命？那些藐视死亡的人，如何活着？那些一无所求的人，如何活着？因为如同一个金刚石般的人，即使被击打一千次，也毫不在意，\Person{保禄}更是如此。他说：\BibleQuote{我活着，却不再是我}，即不再是我；如同在别处：\BibleQuote{我这个人真不幸啊！谁能救我脱离这该死的肉身呢？}\BibleRef{罗 7:24}那些不为食物、不为衣服、不为任何现世事物做任何事的人，又如何活着呢？这样的人甚至不活自然生命：那些不关心任何维持生命之事的人，就不活着。我们活着，我们的每一个行动都关乎现世。但他不活；他不操心任何这里的事。那么他如何活着呢？正如我们平常在世俗事务中说，某人不与我同在，当他没有做任何与我有关的事时。同样，某人不为我而活着。他在别处表明他并未拒绝自然生命：\BibleQuote{我现今在肉身内生活，是生活在对天主子的信仰内；他爱了我，且为我舍弃了自己}。\BibleRef{迦 2:20}即，我活一种新的生命，一种改变了的生命。事实上，他说这一切都是为了安慰斐理伯人。不要说，他说，我将被剥夺这生命，因为即使在活着时，我也没活这生命，而是活基督所愿意的生命。告诉我吧？那些藐视金钱、奢侈、饥饿、干渴、危险、健康、安全的人，他活这生命吗？那些在此一无所有，又常常愿意在必要时抛弃生命，并不执着于它的人，他活这生命吗？绝不。我必须用一个例子向你们说明这一点。让我们想象一个拥有巨大财富、许多仆人和大量金子的人，却一点也不使用这一切；这样的人富有吗？绝不。让他看到他的子女挥霍他的财产，闲散游荡；让他毫不关心他们；当他被打时，甚至不感到痛苦；我们能称他为富人吗？绝不；尽管财富是他的。\BibleQuote{对我来说，}他说，\BibleQuote{活着就是基督；}如果你询问我的生命，那就是祂。\BibleQuote{死亡就是利益。}为何？因为我会更清楚地与祂同在；所以我的死亡更像是走向生命；那些杀我的人不会给我带来任何可怕的事，他们只是把我送到我真正的生命去，使我脱离那不属于我的生命。那么，当你在这里的时候，你不是属于基督的吗？是的，而且达到很高的程度。\par

22节。\BibleQuote{但如果活在肉身内——如果这是我工作的果实，那么我该选择什么，我不知道。}\BibleRef{斐 1:22}\par

恐怕有人会说，如果你所说的是生命，为何基督把你留在此地？他说：\Quote{它是，}\Quote{我工作的果实；}这样就有可能善用现世生活，同时并不活它。恐怕你以为这是对生命的责备。因为如果我们在此没有获得任何益处，为何我们不自杀呢？绝不，他回答说。我们甚至在此也能获益，只要我们不活这生命，而是活另一种生命。但或许有人会说，这能为你结出果实吗？是的！他回答说。现在异端者在哪里？看哪；\Quote{活在肉身内，}这就是\Quote{他工作的果实。}\BibleQuote{我现今在肉身内所活的生命，是活在信德内；}因此它是\Quote{我工作的果实。}\par

\Quote{我不知该选择什么。}真是奇妙！他的智慧何其深邃！他既消除了对现世生命的渴望，却又不加以贬斥。因为他说：\Quote{死亡乃是获益}，借此便消除了渴望；但他说：\Quote{在肉身内生活，乃是我工作的果实}，这便表明，倘若我们按需使用现世生命，并结出果实，它也是必要的；因为若不结果实，它便不再是生命。我们鄙视那些不结果实的树，视其如枯木，将它们交付给火。生命本身属于中性的中间物类，而生活得好与坏则在于我们自己。因此我们并不憎恨生命，因为我们也能生活得好。即使我们不善用它，我们也不归咎于它。为何？因为应受指责的不是生命本身，而是那些不善用之人的自由选择。天主使你活着，是为叫你为祂而活。但你却因败坏而犯罪，使自己承担一切指责。你说什么？告诉我。你不知道该选择什么？在此他揭示了一个伟大的奥秘：他的离世乃出于他自己的权能；因为凡有选择之处，我们便有权能。他说：\Quote{我不知该选择什么。}这在你自己的权能之内吗？是的，他回答，倘若我向天主祈求这恩宠的话。\par

第23节。\Quote{我在这两者之间受困，具有这渴望。} \BibleRef{斐 1:23}\par

请看这位有福者的情感；他以此安慰他们，让他们看到他主宰自己的选择，此事并非出于人的罪，而是天主的安排。他说：你们为何为我的死亡而哀伤？其实早已离世要好得多。\Quote{因为离世与基督同在，实在是好得无比。}\par

第24节。\Quote{然而为你们的缘故，留在肉身内更为必要。} \BibleRef{斐 1:24}\par

这些话是为预备他们面对他的死亡，使他们能高尚地忍受；这是教授真正的智慧。\Quote{为我而言，离世与基督同在更好，}因为死亡本身是中性之事；死亡本身并非恶事，死后受罚才是恶。死亡也非善事，但离世后\Quote{与基督同在}乃是善。死亡之后的事或是善或是恶。\par

因此，我们不要单纯为死者悲伤，也不要单纯为生者欢喜。该如何呢？让我们为罪人悲伤，不仅在他们死时，也在他们活着时。让我们为义人欢喜，不仅在他们活着时，也在他们死后。因为那些虽然活着，却是死的；而这些虽死，却仍活着。那些即使在世时也当被众人怜悯，因为他们与天主为敌；另一些即使离世去了那边，也是有福的，因为他们到了基督那里。罪人无论在哪里，都远离君王。因此他们堪受眼泪；而义人，无论在今生或在来世，都与君王同在；并且在那里，以更崇高、更亲近的方式，不是透过媒介，或通过信德，而是\Quote{面对面。} \BibleRef{格前 13:12}\par

让我们不要只为死者哀哭，却要为那些死于罪中的人哀哭。他们理应哀哭，理应捶胸流泪。因为请告诉我，当我们的罪伴随我们到那里——那里无法脱去罪——时，还有什么希望呢？只要他们还在这里，也许还有极大的期望，以为他们会改变，会变得更好；但当他们去了\OriginalTerm{阴府}{Hades}，在那里不能从悔改中获得任何益处（因为经上记载：\BibleQuote{在\OriginalTerm{阴府}{Sheol}中，谁还能称谢你？}\BibleRef{咏 6:5}），难道他们不值得我们哀悼吗？让我们为这样离世的人哀哭；让我劝你们哀哭，我不阻止；但不要以不当的方式，不要撕裂头发、裸露手臂、抓伤面容、穿着黑衣，而只在灵魂中安静地流下痛苦的眼泪。因为我们可以在没有这一切外在表演的情况下痛苦哀哭。而且不要只当是戏弄。因为许多人的哀哭与戏弄无异。那些公开的哀悼并非出于同情，而是出于表演、攀比和虚荣。许多妇女以此为业。要痛苦哀哭；在家里无人看见时悲叹；这是真正同情的表现；你也因此得益。因为这样哀悼别人的人，会更加努力绝不陷入同样的罪。罪从此将成为你的恐惧。要为不信者哀哭；要为那些与他们毫无区别的人哀哭，那些没有领受\OriginalTerm{光照}{illumination}、没有领受\OriginalTerm{印记}{seal}就离世的人！他们确实值得我们哀哭，值得我们叹息；他们在王宫之外，与罪人、与被定罪者在一起：因为\BibleQuote{我实实在在告诉你们：人若不是由水和圣神生的，不能进入天国。}\BibleRef{若 3:5} 要为那些死于财富，却从未从财富中想到任何灵魂安慰的人哀哭，他们本有能力洗去罪却不愿。让我们都为这些人私下和公开地哀哭，但要合宜、庄重，不要像演戏一样；让我们为他们哀哭，不是一天两天，而是终身。这样的眼泪不是出自无情的激情，而是出自真挚的情感。另一种是出自无情的激情。因此它们很快熄灭，而如果出自对天主的敬畏，它们就常与我们同在。让我们为他们哀哭；让我们尽力帮助他们；让我们为他们想些帮助，哪怕很小，仍然要帮助他们。如何帮助？通过祈祷和恳求别人为他们祈祷，通过不断为他们施舍给穷人。这种行为有些安慰；因为请听天主自己的话，祂说：\BibleQuote{我要为自己的名誉，并为我的仆人达味的缘故，保护这城。}\BibleRef{列下 20:6} 如果只是一位义人的纪念，在为某人行善时有如此大的力量，那么它岂不会更有力量？宗徒们规定在可畏的奥迹中纪念亡者，并非徒然。他们知道这为他们带来巨大益处，巨大好处；因为当全体人民举手站立，司祭的集会，以及那可畏的祭献陈列出来时，我们怎能不为他们向天主恳求而获得垂允呢？我们为那些怀着信德离世的人这样做，而\OriginalTerm{望教者}{catechumens}甚至被认为不配得到这种安慰，被剥夺了一切帮助的手段，只除了一样。这是什么？我们可以为他们施舍给穷人。这种行为在某种程度上使他们得到安息。因为天主愿意我们互相帮助；否则祂为什么命令我们为世界的和平与福祉祈祷？为什么为所有的人祈祷？因为在这数字中包括强盗、盗墓者、小偷、身负无数罪的人；然而我们为所有人祈祷；也许他们会悔改。正如我们为那些活着的人祈祷，他们与死者无异，同样我们也可以为他们祈祷。约伯为他的儿女献祭，使他们脱离罪。他说：\BibleQuote{或许他们的心中已背弃了天主。}\BibleRef{约 1:5} 这就是一个人为自己的儿女所做的预备！他没有像现在许多人那样说：我要留给他们财产；他没有说：我要为他们赢得荣誉；他没有说：我要为他们买官职；他没有说：我要为他们买田地；而是说：\BibleQuote{或许他们的心中已背弃了天主。} 因为那些东西有什么益处呢？留在这世上的完全没有益处。我要使万有的君王对他们仁慈，然后他们就不再缺乏任何东西。经上说：\BibleQuote{上主是我的牧者，我必一无所缺。}\BibleRef{咏 22:4} 这就是巨大的财富，这就是宝藏。如果我们有对天主的敬畏，我们就一无所缺；如果没有，即使拥有王权本身，我们也是所有人中最贫穷的。没有什么比敬畏上主的人更伟大。因为经上说：\BibleQuote{敬畏上主，超越一切。}\BibleRef{德 25:11} 让我们获得这敬畏；让我们为它做一切。如果需要舍命，如果需要肢体被撕裂，我们不要吝惜；让我们做一切以得到这敬畏。因为这样我们将在所有人之上丰盛有余，并将在我们的主基督耶稣内获得那未来的美善，愿光荣归于祂，等等。\par

\SourceRecord{Translated by John A. Broadus.；Nicene and Post-Nicene Fathers, First Series；Vol. 13.；Edited by Philip Schaff.；Buffalo, NY: Christian Literature Publishing Co.,；1889.；<http://www.newadvent.org/fathers/230203.htm>.}

% structure: p0001 source=h1 target=section reason=page-title

\section{讲道第四篇 论斐理伯书}

\BibleRef{斐 1:22-26}\par

\BibleQuote{那么，我不知该选择什么。我夹在两者之间：我渴望离世与基督同在，这实在好得无比；然而，我留在肉身中，为你们更是必要的。我既然深信这点，就知道我要留下，且要与你们众人一同留下，为使你们在信德上长进和喜乐，叫你们因我再次与你们同在，在耶稣基督内以我为夸耀。}\par

再没有比保禄的精神更蒙福的了，因为再没有比它更高贵的。我们大家都害怕死亡——我常这样说——有些人因我们许多的罪过而害怕，我也是其中之一；另一些人因贪生和怯懦而害怕，但愿我永远不是其中之一；因为陷入这种恐惧的人只是动物罢了。那么，连我们都害怕的这件事，他竟祈求，并向它奔赴，说：\BibleQuote{离世实在好得无比。}你说什么？当你即将从大地转向天堂，并与基督同在时，你竟不知道选择什么？不，这与保禄的精神相去甚远；因为若有人能确凿地得到这样的提议，难道他不会立刻抓住它吗？是的，因为既然\BibleQuote{离世与基督同在}不属我们，同样，如果我们能达致此事，留在这里也不属我们。两者都属于保禄和他的精神。他充满确信。什么？你即将与基督同在，却说\BibleQuote{我不知该选择什么}？不仅如此，你还选择这里的，\BibleQuote{留在肉身中}？这世上，难道你没有过极其痛苦的生活，在\BibleQuote{守夜}、船难、\BibleQuote{饥饿和干渴}、\BibleQuote{赤身露体}、忧虑和焦虑中？\BibleQuote{软弱的人}，你就\BibleQuote{软弱}；对那些\BibleQuote{被绊倒的人}，你就\BibleQuote{焦急}。\BibleRef{格后 11:23}\BibleQuote{以极大的坚忍、患难、困苦、艰难、鞭打、监禁、暴乱、禁食、纯洁。}\BibleRef{格后 6:5}\BibleQuote{五次}你\BibleQuote{受了四十下除去一下的鞭打}，\BibleQuote{三次}你\BibleQuote{被棍击打}，\BibleQuote{一次}你\BibleQuote{被石头砸}，\BibleQuote{一昼一夜}你\BibleQuote{在深海里，在江河的危险中，在盗贼的危险中，在城中的危险中，在旷野的危险中，在伪弟兄的危险中。}\BibleRef{格后 11:24}当整个迦拉达民族重新遵守法律时，难道你没有大声呼喊说：\BibleQuote{凡你们想要靠法律成义的，就都从恩宠上跌落了。}\BibleRef{迦 5:4}那时你的悲伤何等巨大，而现在你仍留恋这必朽的生命？即使这一切从未临到你，而是你的成功——无论何处成功伴随着你——都无恐惧、充满喜乐，你难道不该因惧怕不确定的未来而赶紧奔向某个港湾吗？请告诉我：哪个商人，他的船满载数不尽的财富，当他可以进港休息，却宁愿还留在海上？哪个摔跤手，当他可以加冕时，却宁愿继续比赛？哪个拳击手，当他可以戴上冠冕时，却选择重新进入竞赛，把自己的头献上挨打？哪个将军，当他可以带着美名和战利品脱离战争，与君王在宫殿里休息，却选择继续劳苦、列阵作战？那么，你这个过得如此痛苦的生命，怎么还想留在这里？你难道没有说过：我害怕\BibleQuote{免得我给别人传了福音，自己反而被淘汰。}\BibleRef{格前 9:27}即使没有别的缘故，单为这个，你也当渴望解脱；即使现在有无数好处，仍要为你的渴望——基督。\par

哦，保禄的精神！再没有像它一样的，也永远不会有了！你害怕未来，被无数可怕的事包围，你还不愿与基督同在？不，他回答，这全是为了基督，好让我使那些我使之成为祂仆人的人对祂更忠诚，好让我所栽种的园地结出丰硕的果实。\BibleRef{格前 3:9}你难道没有听见我声明：我不寻求\BibleQuote{自己的利益}\BibleRef{格前 10:33}，而是邻人的吗？你难道没有听见这些话：\BibleQuote{我宁愿自己从基督被诅咒}\BibleRef{罗 9:3}，好使许多人归向祂？我，选择了那部分，难道不更该选择这个吗？我不乐意因这拖延和延迟而损害自己，使他们得救吗？\par

\BibleQuote{谁能宣扬你的大能，上主？}\BibleRef{咏 6:2}因为你没有让保禄被隐藏，因为你向世界彰显了这样一个人。当祢创造星辰时，天主的众天使都齐声赞美祢\BibleRef{约 38:7}，同样，当祢创造太阳时，也如此，但都不及祢向全世界彰显保禄时那样。藉此，大地变得比天空更光辉，因他比太阳光更明亮，放射出更灿烂的光芒，散发出更欢悦的光线。这个人为我们结出了何等的果实！不是使我们的五谷丰登，不是滋养我们的石榴，而是产生并成全了圣德的果实，并且当它们衰败时，不断恢复它们。因为太阳本身对已经枯朽的果实无能为力，但保禄从罪恶中唤回了那些有多种败坏的人。太阳让位于黑夜，但他却制服了魔鬼。没有什么制服过他，没有什么征服过他。太阳升上天空，射下光芒，但他从下面升起，不是以光充满天地之间，而是当他开口时，使天使充满极大的喜乐。因为如果\BibleQuote{对于一个罪人悔改，天上也有喜乐}\BibleRef{路 15:7}，而他第一次讲道就抓住了众人，难道他不使天上的众天使充满喜乐吗？我说什么？只需保禄的名字被提及，诸天就欢跃。因为如果当以色列人\BibleQuote{出离埃及时，诸山跳跃如公羊}\BibleRef{咏 113:4}，那么当人从地上升到天上时，喜乐该是何等大呢！\par

第二十四节。因此，\BibleQuote{在肉身内存留，为你们是更有必要的} \BibleRef{斐 1:24}。\par

那么，我们还有什么借口呢？常有这样的事：一个人拥有一个小小的穷城，却不愿离开到别处去，宁愿安于自身。\Person{保禄}本可以到\Person{基督}那里去，却不愿意（他渴望\Person{基督}到甚至为了祂的缘故宁愿选择地狱），但仍留在尘世为人类争战。我们还有什么借口呢？那么，我们甚至敢提到\Person{保禄}吗？请看他的行为。他表明离去更好，却说服自己不要悲伤；他向他们表明，如果他留下，是为了他们的缘故才留下，而不是因为那些图谋害他的人邪恶。他又加上理由，要使他们确信。因为如果这是必要的，也就是说，我无论如何都会留下，而且我不只是简单\Quote{留下}，而是\Quote{将留在你们那里}。因为这话的意思就是：\Quote{我将与你们同在}，也就是我要见到你们。为什么呢？\Quote{为你们的进步和信德中的喜乐}。这里他也激励他们，要他们谨慎自己。他说，如果我为你们的缘故留下，你们要留意不要羞辱我的留下。\Quote{为你们的进步}，我选择了留下，尽管我本要见到\Person{基督}。我选择了留下，因为我的同在促进你们的信德和喜乐。那么怎样呢？他只为了\Person{斐理伯}人留下吗？他留下不单是为了他们，但他说这话是为了表示对他们的关心。他们如何\Quote{进步}于\Quote{信德}呢？为使你们更坚固，如同幼鸟，在羽毛长齐之前需要母亲一样。这是他大爱的明证。同样，我们也激励你们中的一些人，对你们说：为你们的缘故我留下了，好使我把你们造就成好人。\par

第二十六节。\BibleQuote{叫你们在基督耶稣内的夸耀，因我再到你们那里去，在你们身上越发增加} \BibleRef{斐 1:26}。\par

你看，这解释\Quote{与你们同在}一词。请看他的谦卑。他说了\Quote{为你们的进步}，又表明这也是为他自己的益处。当他写给罗马人时也这样行，说：\Quote{就是使我们因你们彼此得到安慰} \BibleRef{罗 1:11}。他先前曾说：\Quote{把一些神恩分给你们}。而\Quote{叫你们的夸耀越发增加}是什么意思？这夸耀就是他们在信德上的坚固。正直的生活就是在\Person{基督}内的夸耀。你却说\Quote{你们的夸耀因我再到你们那里去}？是的，他回答：\Quote{什么是我们的希望、喜乐或所夸耀的冠冕呢？不正是你们吗？} \BibleRef{得前 2:19} 因为\Quote{你们是我们的夸耀，正如我们也是你们的夸耀} \BibleRef{格后 1:14}，也就是说，我能够在你们身上大大喜乐。你怎么说\Quote{叫你们的夸耀越发增加}？当你们进步时，我更能夸耀。\par

\Quote{因我再到你们那里去}。那么！他来到他们那里了吗？你们查考一下他是否来了。\par

第二十七节。\BibleQuote{惟愿你们的生活相称于基督的福音} \BibleRef{斐 1:27}。\par

你看见吗？他所有的话都倾向于引导他们达到这一件事，就是品德的进步。\Quote{惟愿你们的生活相称于基督的福音}。\Quote{惟愿}这个词是什么意思？无非是指，只有这个，而非其他，才是我们所当追求的。如果我们有它，便没有什么痛苦会临到我们。\Quote{这样，我或来看你们，或不在你们那里，都可以听到你们的情况}。他说这话，并非好像他改变了主意，不再打算访问他们。而是说，如果这事发生，即使我不在，我也能喜乐。也就是说，\Quote{如果}我\Quote{听说你们在同一圣神内，以同一灵魂站稳}。这正是最能使信徒合一、保持爱德不断的事：\Quote{使他们合而为一} \BibleRef{若 17:11}。因为\Quote{一家若自相纷争，必站立不住} \BibleRef{谷 3:24}。为此，他处处教训他的门徒要同心合意。\Person{基督}说：\Quote{如果你们彼此相爱，众人因此就认出你们是我的门徒} \BibleRef{若 13:35}。也就是说，不要期盼我而沉睡，等候我的到来，然后当你们看不到我来时就灰心。因为我从报告也能同样得到喜乐。\par

\Quote{在同一圣神内}是什么意思？是藉着同一恩宠的恩赐，即同心同德与热忱；因为圣神是唯一的，他也表明了这一点；因为当我们都拥有\Quote{同一圣神}时，我们就能够站在\Quote{同一灵魂}上。请看\Quote{一}这个词如何用于表示同心。看他们的灵魂虽多却被称为一。古时就是这样。经上记载：\Quote{众信徒都是一心一意，为福音的信德一同奋斗} \BibleRef{宗 4:32}。他说互相奋斗，难道信德在奋斗吗？难道他们彼此摔跤吗？而是说，你们在为了福音的信德奋斗时互相帮助。\par

第二十八节。\BibleQuote{什么事都不怕敌人的惊吓；这为他们是沉沦的明显证据，但为你们却是得救的} \BibleRef{斐 1:28}。\par

他说得好：\Quote{惊慌}，这就是我们从敌人那里遭遇的，他们只会恐吓。\Quote{一无所惧}，因此，他说，无论发生什么事，无论是危险还是阴谋。因为这是那些站立得稳的人的本分；敌人除了恐吓之外什么也做不了。既然当\Person{保禄}遭受如此无数的苦难时，他们很可能大为困扰，他说，我劝你们不仅不要动摇，而且不要惊慌，反而要真心鄙视他们；因为如果你们这样受影响，你们将立即借此同时显明他们的毁灭和你们的得救。因为当他们看到，他们无数的阴谋无法恐吓你们时，他们会将此视为自己毁灭的证据。因为当迫害者不能胜过被迫害者，阴谋者不能胜过他们阴谋的对象，有权势者不能胜过那些在他们权下的人时，这不是不言而喻的吗：他们的灭亡近了，他们的权势是虚无，他们的份是虚假，他们的份是软弱的？\Quote{而这，}他说，\Quote{来自天主。}\par

第29节。\BibleQuote{因为为了基督，赐给你们的不但是信祂，而且是为祂受苦。}\BibleRef{斐 1:29}\par

他又藉着将一切归因于天主，教导他们谦卑的精神，并说为基督受苦是出于恩宠，是恩宠的赐予，是白白的恩赐。因此，不要以这恩宠的赐予为耻，因为它比复活死者或行奇迹的能力更奇妙；因为在那里我是欠债者，但在这里我却有基督作我的欠债者。所以我们不仅不该羞耻，反而应当喜乐，因为我们拥有这恩赐。他称德行为恩赐，但不像其他事物那样；因为那些完全来自天主，而在这些德行中我们也有份。但既然即使在此，最大部分也属于天主，他就完全归功于祂，并非要推翻我们的自由意志，而是要使我们谦卑和端正。\par

第30节。\BibleQuote{你们有过同样的斗争，如同你们在我身上所看见的那样。}\BibleRef{斐 1:30}；也就是说，你们也有一个榜样。他再次提升他们，藉此表明他们的斗争处处与他的相同，他们的挣扎与他的相同，无论是各别的，还是他们与他联合承担考验。他没有说'你们听说了'，而是说\Quote{你们看见了}，因为他在斐理伯也曾奋斗。这确实是一种卓越的德行。因此，在写给迦拉达人的信中也说：\BibleQuote{你们受了这么多的苦，难道都是白费的吗？如果真是白费的话。}\BibleRef{迦 3:4} 再次，在写给希伯来人的信中说：\BibleQuote{请你们回想从前那些日子，你们受了光照以后，忍受了极大的痛苦斗争；一方面，你们因讥讽和磨难成了观瞻的对象。}\BibleRef{希 10:32} 再次写给马其顿人，即得撒洛尼人，他说：\BibleQuote{因为他们自己传报我们的事，说我们怎样进入了你们中间。}\BibleRef{得前 1:9} 又说：\BibleQuote{兄弟们，你们自己知道我们进入你们中间，并不是徒然的。}\BibleRef{得前 2:1} 同样，他为众人作证同样的事：劳苦和奋斗。但如今在我们中间你们找不到这样的事；如今若有人只在财物上受一点损失，就算很多了。而且在财物方面，他也为他们作了大见证。因为对一些人他说：\BibleQuote{因为你们欣然接受了你们的财物被抢掠。}\BibleRef{希 10:34}；对另一些人他说：\BibleQuote{因为马其顿和阿哈雅乐意作出某种捐输，为资助穷人。}\BibleRef{罗 15:26} 以及\BibleQuote{你们的热情激励了他们中的许多人。}\BibleRef{格后 9:2}\par

你可看见当时之人的赞美？但我们连殴打和击打也忍受不了，既不能忍受侮辱，也不能忍受财产的损失；他们立即热心起来，全都奋力如同殉道者，而我们却对基督的爱冷淡了。我又不得不指责现今的事；我该当如何呢？这违反我的意愿，但我不得不如此。倘若我能以缄默对待所行之事，闭口不言，不提任何事，便能消除它们，我应当缄默。但若适得其反；沉默不但不能消除这些事，反而使它们变得更糟，我们便被迫发言。因为责备罪人的人，即使别无作为，也阻止他们更进一步。因为没有哪个如此无耻和卤莽的灵魂，在听见有人不断责备它时，不会回转并收敛其恶行的过度放荡。确实，确实，即使在无耻之人中，也仍有一小部分羞耻。因为天主在我们本性中撒下了羞耻的种子；既然恐惧不足以使我们调整到正确的心态，衪也预备了许多其他避免犯罪的方式。例如，一个人被责备，对成文法律的恐惧，对名誉的爱惜，结交友谊的愿望；所有这些都是避免犯罪的道路。常有不因天主而行的事，却因羞耻而行；不因天主而行的事，却因怕人而行。我们所寻求的，首先是不犯罪，然后我们便会为天主而成功地做到这点。不然，\Person{保禄}为何劝告那些将要战胜敌人的人，不是出于对天主的敬畏，而是基于等候报复？\Quote{因为这样做，}他说，\Quote{你便是将炭火堆在他头上。}\BibleRef{罗 12:20}因为他的首要愿望，是我们的德行得以建立。正如我那时所说，我们心中有羞耻之感。我们有许多自然的善情，引领我们走向德行；例如，我们所有人天生容易被同情所感动，没有其他善事如此内在于我们的本性，唯独此一善事。由此任何人可能合理探问，为何这些种子超越其他一切而被撒在我们本性中，使我们因眼泪而融化，使我们转向怜悯，并预备好同情。没有人天生懒惰，没有人天生不顾名誉，没有人天生不争强好胜，但怜悯深植于每个人的本性中，无论他多么凶猛和不温和。这有什么奇怪呢？我们怜悯禽兽，如此丰盛的怜悯深植于我们内。若我们看见一只幼狮，我们会有些触动；对于我们的同类，更是如此。看，有多少残疾的人！这足以引导我们怜悯。没有比仁慈更悦乐天主的事。因此，司铎、君王和先知都被傅油，因为他们藉油获得了天主对人的爱的预象；他们还进一步学到，统治者应当更富有仁慈。这显示圣神藉仁慈降临于人，因为天主怜悯人并善待人类。因为经上记载：\Quote{你怜悯众生，}\Quote{因为你是全能的。}\BibleRef{智 11:23}为此，他们被傅油；事实上，正是出于仁慈，衪设立了司祭职。君王也被傅油；若有人赞美一位统治者，他能提及的再没有比仁慈更合适他的了。因为怜悯是权力所特有的。要想到世界是因怜悯而建立的，然后效法你的主。\Quote{人的仁慈仅限于近人，但主的仁慈却普及一切血肉。}\BibleRef{德 18:13}如何理解\Quote{普及一切血肉}？无论你指罪人，还是义人，我们都需要天主的仁慈；我们都享有它，无论是\Person{保禄}，还是\Person{伯多禄}，或是\Person{若望}。请听他们自己的话，无须我来言说。这位蒙福者说了什么？\Quote{但是我蒙受了怜悯，因为我当时是在不信中行事。}\BibleRef{弟前 1:13}那么，此后不再需要仁慈了吗？听他怎么说：\Quote{然而，我比他们众人更劳苦；不是我，而是天主的恩宠与我同在。}\BibleRef{格前 15:10}关于\Person{厄帕夫洛狄托}，他说：\Quote{他实在病了，几乎要死；但天主怜悯了他，不但怜悯了他，也怜悯了我，免得我忧上加忧。}\BibleRef{斐 2:27}他又说：\Quote{我们被压得太过分，甚至超出了力量，以致我们连活命的指望都绝了。我们自己心里也断定是必死的，叫我们不靠自己，只靠那使死人复活的天主；衪曾救我们脱离那极大的死亡，并且仍要救我们。}\BibleRef{格后 1:8}又说：\Quote{我从狮子口中被救出来；主必救我。}\BibleRef{弟后 4:17}我们到处都会发现他以此夸耀，即他是因仁慈而得救。\Person{伯多禄}也是如此伟大，因他蒙受了仁慈。请听基督对他说：\Quote{撒旦想要筛你们像筛麦子一样；但我已为你祈求，使你的信德不至失落。}\BibleRef{路 22:31}\Person{若望}也同样因仁慈而变得如此伟大，总之，他们所有人都是。请听基督所说：\Quote{不是你们拣选了我，是我拣选了你们。}\BibleRef{若 15:16}因为我们都需天主的仁慈，正如经上所记：\Quote{天主的仁慈普及一切血肉。}但如果这些人需要天主的仁慈，对其他人该说什么呢？因为，请告诉我，为何衪\Quote{叫太阳上升，光照恶人，也光照善人}？若衪一年不降雨，岂不毁灭一切？若衪降下倾盆大雨呢？若衪降下火呢？若衪派遣苍蝇呢？但我说什么呢？若衪像从前那样行动，岂不是所有人都要灭亡？若衪震动大地，岂不是所有人都要灭亡？现在正合时宜地说：\Quote{世人算什么，你竟眷顾他？}\BibleRef{咏 8:4}若衪只是威胁大地，所有人都会成为一个坟墓。\Quote{看哪，万民像桶中的一滴水，}经上记着，\Quote{在衪眼中如同微尘，又如天平上的细沙。}\BibleRef{依 40:15}衪毁灭万有再创造，对我们来说，就如我们转动天平一样容易。衪对我们有这样的大能，看见我们天天犯罪，却不惩罚我们，若不是因仁慈容忍我们，又是为何？就连野兽也是因仁慈而存在：\Quote{主啊，你必保护人和牲畜。}\BibleRef{咏 35:7}衪眷顾大地，使之充满活物。这是为什么？是为了你！衪为何造了你？因衪的良善。\par

没有什么比油更好的了。它是光的成因，在那里也是光的成因。\BibleQuote{那时，你的光明将如晨光般射出}\BibleRef{依 58:8}，先知说，如果你向邻人表示怜悯。正如天然的油含有光，同样，怜悯（施舍）也赐予我们伟大而奇妙的光。保禄也多次提及这种怜悯。在一处，听他说：\BibleQuote{只是我们要记念穷人}\BibleRef{迦 2:10}。在另一处，\BibleQuote{若我也该去}\BibleRef{格前 16:4}。无论你转向何处，你都会看到他对此事的挂虑。又说：\BibleQuote{我们的人也应学习行善}\BibleRef{铎 3:14}。又说：\BibleQuote{这些事是美好而有益于人的}\BibleRef{铎 3:8}。请听另一位所说：\BibleQuote{施舍能救你脱离死亡}\BibleRef{多 12:9}；如果你除去怜悯，\BibleQuote{主，主，谁能站立得住}\BibleRef{咏 129:3}；经上又说，你若\BibleQuote{与仆人审判}\BibleRef{咏 142:2}；\BibleQuote{人是伟大的}；为什么？\BibleQuote{而怜悯的人是可敬的。}\BibleRef{箴 20:6}因为这才是人的真正特性，即怜悯；更确切地说，这是天主的特性，即施以怜悯。你看见了吗？天主的怜悯是何等强大？祂藉此创造了万物，形成了世界，造了天使，全是出于纯粹的良善。为此，祂也以地狱相威胁，好使我们能获得天国，而藉着怜悯我们确实获得天国。因为天主独一，为何创造如此众多？难道不是出于良善？难道不是出于对人的爱？你若问为何如此这般，你总会在良善中找到答案。让我们对邻人施以怜悯，好使怜悯也施予我们。这些怜悯的行为，与其说是施予他们，不如说是为我们自己积攒起来，以待那一日。当火焰炽烈时，这油（怜悯）能熄灭火焰，并给我们带来光明。这样，我们将从地狱之火中获得释放。因为祂何时会施以怜悯和慈爱？怜悯源自爱！没有什么比没有怜悯更惹天主发怒的了。\Quote{有人带了一个欠他一万塔冷通的来，他动了慈心，就豁免了他。而那个人的同伴欠他一百德纳，他却扼住他的喉咙。因此主人把他交给刑役，直到他还清所欠的债务。}让我们听了这些，对那些欠我们钱财或得罪我们的人心怀怜悯。愿没有人记仇，至少如果他不想伤害自己的话；因为他并不如此伤害别人（正如他伤害自己）。因为他要么会报复，要么还未报复；而你自己，若不宽恕邻人的罪，却寻求天国？为避免这事发生在我们身上，让我们宽恕所有的人（因为我们正是在宽恕自己），好使天主宽恕我们的罪，使我们获得那预备的福乐，藉着恩宠和慈爱等。\par

\SourceRecord{Translated by John A. Broadus.；Nicene and Post-Nicene Fathers, First Series；Vol. 13.；Edited by Philip Schaff.；Buffalo, NY: Christian Literature Publishing Co.,；1889.；<http://www.newadvent.org/fathers/230204.htm>.}

% structure: p0001 source=h1 target=section reason=page-title

\section{论\WorkTitle{斐理伯书}第五篇讲道}

斐理伯书 2:1-4 \BibleRef{斐 2:1-4}\par

\Quote{所以，若在基督内有任何鼓励，若有爱的任何安慰，若有圣神的任何共融，若有任何慈悲与怜悯，你们就要使我喜乐圆满，好使你们同心合意，有同样的爱，同心合意，同心一致；不论做什么，都不出于结党或虚荣，却要谦卑，各人看别人比自己强；各人不要只求自己的利益，也要顾及别人的利益。}\par

没有什么比一位神修导师更美好、更亲切的了；他超越任何生身父亲的慈爱。试想，这位有福之人是如何恳求斐理伯人有关他们自身利益的事。他在劝勉他们同心合意——那一切美善之源——时说了什么？请看他是多么恳切、多么热切、多么同情地说：\Quote{所以，若在基督内有任何鼓励，}就好像他说：若你们对我有任何顾念，若你们对我有任何关心，若你们曾从我手中得到过任何好处，就请这样做。当我们主张一件我们看为比一切都重要的事时，我们就会采用这种恳切的方式。因为若不是我们把它置于一切之上，我们就不会愿意在它里面得到我们一切的报答，也不会说通过它一切都被代表了。我们确实会提醒人们我们属肉体的要求；例如，如果一位父亲说：若你对父亲有任何敬畏，若你对我养育你的任何记念，若你对我有任何爱，若你对我曾给你的尊荣有任何记忆，若你对我的任何恩惠，就不要与你的兄弟为敌；也就是说，我要求用这一切来回报这件事。\par

但保禄并非如此；他并没有让我们想起任何属肉体的事，而是全都属灵的益处。也就是说，若你们愿意在我的试探中给我任何鼓励，以及在基督内的鼓励，若有任何爱的安慰，若你们愿意显示任何在圣神内的共融，若你们有任何慈悲与怜悯，就使我喜乐圆满。\Quote{若有任何慈悲与怜悯。}保禄把门徒们的同心合意说成是对自己的怜悯，从而显示出危险是极大的，如果他们不一心一意的话。若我能从你们那里得到安慰，若我能从我们的爱中得到任何安慰，若我能在圣神内与你们相通，若我能在主内与你们共融，若我能从你们手中得到怜悯与慈悲，请用你们的爱显示这一切的回报。如果你们彼此相爱，这一切我就都得到了。\par

\Quote{你们就要使我喜乐圆满。} \BibleRef{斐 2:2}\par

为了不让劝勉看起来是针对仍有欠缺的人，请看他是如何说的：他没有说\Quote{使我喜乐}，而是说\Quote{使我喜乐圆满}；也就是说，你们已开始在我内播种它，你们已给了我一些平安的份量，但我希望达到它的圆满？你说，你愿意什么？要我们把你从危险中解救出来？要我们供应你的需要？不，而是\Quote{使你们同心合意，有同样的爱，}你们已经开始这样做了，\Quote{同心合意，同心一致。}请看，他因着极大的爱重复同样的事多少次！\Quote{使你们同心合意，}或者更确切地说，\Quote{使你们一心一意。}因为这一句比\Quote{同心合意}更进一步。\par

\Quote{有同样的爱。}也就是说，不要仅仅只在信仰上，也要在一切事上；因为有一种情况是同心合意却无爱。\Quote{有同样的爱，}也就是说，要同样地爱和被爱；不要你享受许多爱，却付出较少爱，以至于在这事上也贪婪；但不要容忍自己这样。他补充说\Quote{同心合意，}也就是说，以同一个灵魂占据众人的身体，不是在实体上（因为那是不可能的），而是在意向和心意上。让一切事都好像从一个灵魂发出。\Quote{同心合意}是什么意思？他在说\Quote{同心一致}时表明了这一点。愿你们的心思成为一体，好像出自同一个灵魂。\par

\Quote{不论做什么，都不出于结党或虚荣。} \BibleRef{斐 2:3}\par

他最终向他们要求这一点，并告诉他们如何能做到。\Quote{不论做什么，都不出于结党或虚荣。}正如我常说的，这是一切邪恶的原因。由此产生争斗和冲突。由此产生嫉妒和纷争。由此爱变冷淡，当我们爱人的称赞，当我们做多数人所施予的荣誉的奴隶时，因为人不可能既做赞誉的奴隶，又做天主的真仆人。那么，我们如何逃避虚荣呢？因为你还没有告诉我们方法。那么请听下文。\par

\Quote{却要谦卑，各人看别人比自己强。}哦！他所提出的教训是多么充满真智慧，多么普遍地凝聚了我们的救恩！他的意思是，你若认为别人比你大，并且说服自己，更甚者，你不但这样说，而且完全确信，那么你就将尊荣归给他；你若将尊荣归给他，看见他被别人尊荣时，就不会不悦。所以，不要只认为他比你大，而是\Quote{比自己强，}这是一个非常大的优越性，而你并不觉得奇怪，也不因此痛苦，如果你看见他受尊荣。即使他藐视你，你也能高尚地忍受，因为你已认为他比你大。即使他辱骂你，你也顺服。即使他恶待你，你也默默忍受。因为一旦灵魂完全确信他更大，当受到他的恶待时，就不会陷入愤怒，也不陷入嫉妒，因为没有人会嫉妒远远超越自己的人，因为一切都属于他的优越性。\par

因此他劝诫这一方应有这样的心意。但是当那受你们如此尊敬的人也这样对待你们时，想想这是何等双重谦让之墙（参\BibleRef{斐 4:5}）；因为当你这样敬重他，他也同样敬重你时，任何痛苦的事都不会发生；因为如果单方这样做就足以消除一切纷争，双方都这样做时，谁能攻破这防线呢？连魔鬼自己也不能。这防御是三倍、四倍，甚至多倍的，因为谦卑是万善之源；为使你明白这点，请听先知的话：\BibleQuote{你若要祭献，我必献上；全燔祭你必不喜爱。天主，我的祭献就是痛悔的精神，痛悔和谦卑的心，天主必不轻看。}\BibleRef{咏 51:16}不单是谦卑，而是深切的谦卑。犹如在物质上，\Quote{碎}的不会起来攻击\Quote{硬}的，无论遭受多少苦，宁可自身消亡也不攻击对方；同样，灵魂即使不断受苦，也宁愿死，不愿以攻击来报复。\par

我们要这样可笑地骄傲到几时呢？正如我们看见小孩子挺胸昂首、自高自大，或看见他们捡起石头乱扔就会发笑，同样，人的骄傲属于幼稚的理智和未成熟的心智。\BibleQuote{为什么尘土和灰烬骄傲呢？}\BibleRef{德 10:9}人啊，你心高气傲吗？为什么？告诉我有什么好处？你为什么对同类心高气傲？你们不都共享同样的本性吗？同样的生命吗？你们不是从天主领受了同样的尊荣吗？但你聪明吗？你应当感恩，而不是骄傲。傲慢是忘恩负义的首要行为，因为它否认恩宠的恩赐。骄傲的人骄傲，仿佛他凭自己的力量胜过别人；凡以为自己这样胜过别人的，就是对赐予尊荣的天主忘恩负义。你有什么好的？要感谢那赐予者。请听\Person{若瑟}和\Person{达尼尔}所说的话。当埃及王派人召见他，在全军面前问他关于那件事，那些精通此道的埃及人都束手无策时，他几乎要从他们那里夺走一切，显得比占星家、术士、巫师和当时所有的智者更聪明，而且他出身俘虏和奴仆，还只是个青年（他的荣耀因此更大，因为一个人出名与出人意料不同，出乎意料使他更令人钦佩）；那么，当他来到\Person{法郎}面前时，他说\Quote{是，我知道}吗？不，是什么？当没有人催逼他时，他出于自己卓越的精神说：\Quote{解梦不是出于天主吗？}你看，他立刻光荣了他的主，因此他也得了光荣。这也不是小事。因为天主启示给他，比他本人卓越要伟大得多。因为他证明了他的话可信，这也是他与天主亲密无间的极大明证。没有什么比作天主的挚友更好的了。圣经说：\BibleQuote{因为如果\Person{亚巴郎}因行为成义，他就有可夸的，但不是在天主前。}\BibleRef{罗 4:2}因为蒙受恩宠的人夸耀天主，因他爱他、他的罪得赦免，那凭行为的人也有可夸的，但不在天主面前，像前者那样（因为这证明我们极度的软弱）；那从天主领受智慧的人，岂不更加令人钦佩？他光荣天主，也受天主光荣，因为他说：\Quote{尊敬我的，我必尊重他。}\BibleRef{撒上 2:30}\par

再听那从\Person{若瑟}的后裔降生者，无人比他更有智慧。他说：\Quote{你比\Person{达尼尔}更有智慧吗？}\BibleRef{则 28:3}这位\Person{达尼尔}，当\Place{巴比伦}所有的智者、占星家、先知、术士、巫师，甚至他们的智慧不仅被揭露，而且完全被摧毁时（他们的被摧毁清楚地证明他们以前是在欺骗），\Person{达尼尔}上前准备解答国王的问题，却不将尊荣归于自己，而首先将一切归于天主，说：\BibleQuote{王啊，我绝不是凭自己有超过众人的智慧而蒙启示。}\BibleRef{达 2:30}\BibleQuote{于是国王崇拜他，下令向他献祭。}\BibleRef{达 2:46}你看见他的谦卑吗？你看见他卓越的精神吗？你看见这谦逊的习性吗？也听宗徒们说：\Quote{为什么注视我们，好像我们凭自己的能力或虔诚使这人行走呢？}\BibleRef{宗 3:12}又说：\Quote{我们和你们有同样的性情。}\BibleRef{宗 14:15}既然他们这样拒绝所受的尊荣，这些人因基督的谦卑和能力行了比基督更大的事（因为他说：\BibleQuote{信我的，要做比这些更大的事}\BibleRef{若 14:12}），难道我们这些可怜、悲惨的人，连蚊子都赶不走，更不用说魔鬼了，没有能力使一个人受益，更不用说全世界，却如此看重自己，连魔鬼都比不上我们吗？\par

没有什么比骄傲更与基督徒的灵魂格格不入的了。我说的是骄傲，不是勇敢或勇气，因为这些是与灵魂相宜的。但这是一回事，那是另一回事；谦卑是一回事，卑屈、奉承和谄媚是另一回事。\par

现在，若你愿意，我将为你举出所有这些品质的例子。因为这些似乎相反的事物，在某些方面彼此相近，如同莠子与麦子，荆棘与玫瑰。然而婴孩容易被欺骗，那些真正成人的、精通属灵农事的人，却懂得如何从恶中分辨真正的善。让我从圣经中为你举出这些品质的例子。什么是谄媚、卑屈和奉承？齐巴不合时宜地谄媚达味，并诬告他的主人。\BibleRef{撒下 16:1}阿希托费耳更是谄媚阿贝沙隆。\BibleRef{撒下 17:1}但达味却不如此，他是谦卑的。因为诡诈的人是谄媚者，例如他们说：\Quote{大王万岁！}\BibleRef{达 2:4}再看那些术士是何等谄媚。\par

我们会在宗徒大事录中保禄身上找到许多这样的例子。当他和犹太人辩论时，并不谄媚他们，却心怀谦卑——因为他知道如何放胆直言——如他说：\Quote{诸位弟兄！我虽然没有做什么反对人民或祖先规矩的事，却从耶路撒冷被囚禁交出来。}\BibleRef{宗 28:17}\par

这些话是出于谦卑。请听他在后面如何斥责他们：\Quote{圣神说得对啊，他说：你们听是听得见，却总不明白；看是看得见，却总看不透。}\BibleRef{宗 28:25}\par

你看见他的勇气了吗？再看看洗者若翰在黑落德面前所用的勇气：他说：\Quote{你占有你兄弟斐理伯的妻子是不合法的。}\BibleRef{谷 6:18}这是大胆，这是勇气。但史米的言语却不然，他说：\Quote{滚，你这流人血的人！}\BibleRef{撒下 16:7}然而他也放胆直言，但这不是勇气，而是狂妄、傲慢和放肆的舌头。依则贝耳也辱骂耶胡，她说：\Quote{杀了你主人的人}\BibleRef{列下 9:31}，但这是狂妄，不是大胆。厄里亚也责骂，但那是大胆和勇气：\Quote{不是我叫以色列遭难，而是你和你的父家。}\BibleRef{列上 18:18}厄里亚又向全体民众放胆直言：\Quote{你们要跛着双腿到几时呢？}\BibleRef{列上 18:21}这样的责斥是大胆和勇气。先知们也是如此行，但另一种则是狂妄。\par

你想看到既谦卑又不谄媚的言语吗？听保禄说：\Quote{但对我来说，受你们审断，或受人间法庭审断，都是极小的事；我也不审断自己。因为我虽不觉得自己有错，却不能因此就成义。}\BibleRef{格前 4:3}这是基督徒应有的精神；他又说：\Quote{你们中间有彼此相争的事，竟敢在不义的人面前诉讼，而不在圣徒面前吗？}\BibleRef{格前 6:1}\par

你想看到愚昧犹太人的谄媚吗？听他们说：\Quote{除了凯撒，我们没有君王。}\BibleRef{若 19:15}你想看到谦卑吗？再听保禄的话：\Quote{因为我们不是宣传我们自己，而是宣传耶稣基督为主，并因耶稣的缘故，做你们的仆人。}\BibleRef{格后 4:5}你想看到既谄媚又狂妄吗？纳巴耳身上的\Quote{狂妄}\BibleRef{撒上 25:10}和齐弗人的\Quote{谄媚}\BibleRef{撒上 23:20}；因为他们在谋划中背叛了达味。你想看到\Quote{智慧}\BibleRef{撒上 26:5}而非谄媚吗？达味的智慧，他如何将撒乌耳置于掌握之下，却饶恕了他。你想看到那些杀害默黎巴耳之人的谄媚吗？达味也杀了他们。最后，仿佛是纲要，总括一切：狂妄是在无缘无故发怒、侮辱他人、或为自己报复、或以不义方式放肆时显出的；而大胆和勇气则是当我们为了讨天主喜悦，而敢于面对危险和死亡，轻看友情和敌意。反之，谄媚和卑屈是当人不是为了任何正当目的而奉承别人，而是贪求今世的事物时；谦卑则是当人为了讨天主喜悦而这样做，并且从自己的高位降下，以便成就一件伟大而可赞叹的事。如果我们知道这些，又实行这些，便是有福的。因为仅仅知道还不够，圣经说：\Quote{不是听法律的，而是行法律的才成义。}\BibleRef{罗 2:13}是的，知识本身若不伴随行动和美德，反倒定我们的罪。因此，为逃避定罪，让我们追随实践，以获取那应许给我们的福乐，凭我们的主耶稣基督的恩宠和爱。\par

\SourceRecord{Translated by John A. Broadus.；Nicene and Post-Nicene Fathers, First Series；Vol. 13.；Edited by Philip Schaff.；Buffalo, NY: Christian Literature Publishing Co.,；1889.；<http://www.newadvent.org/fathers/230205.htm>.}

% structure: p0001 source=h1 target=section reason=page-title

\section{论斐理伯书第六讲}

\BibleRef{斐 2:5-8}\par

\BibleQuote{你们该怀有基督耶稣所怀有的同一心情：他虽具有天主的形体，并没有以自己与天主同等为应当把持不舍的，却使自己空虚，取了奴仆的形体，与人相似，形状也一见如人；他贬抑自己，听命至死，且死在十字架上。}\par

我们的主耶稣基督在劝勉门徒们行伟大的事时，常将他自己、圣父和先知们作为榜样放在他们面前；例如当他说：\BibleQuote{因为他们也这样迫害了在你们以前的先知。}\BibleRef{玛 5:12}；又说：\BibleQuote{如果他们迫害了我，也要迫害你们。}\BibleRef{若 15:20}；以及：\BibleQuote{你们跟我学吧，因为我是良善心谦的。}\BibleRef{玛 11:29}；又说：\BibleQuote{你们应当慈悲，就像你们的父那样慈悲。}\BibleRef{路 6:36}。有福的保禄也这样做了：在劝勉他们谦卑时，他引出了基督。他不只是在这里这样做，在论及对穷人的爱德时，他也这样说：\BibleQuote{因为你们知道我们的主耶稣基督的恩宠：他本是富有的，为了你们却成了贫穷的。}\BibleRef{格后 8:9}。没有什么比知道自己在善工上相似天主更能激发一个伟大而有哲学修养的灵魂去行善了。还有什么鼓励能与此相比？没有。保禄深知这一点，当他想要劝勉他们谦卑时，先是恳求和劝告他们，然后为了敬畏他们，他说：\BibleQuote{你们要一心一意站稳；}\BibleQuote{这对他们是灭亡的明显证据，但对你们却是得救的证据。}\BibleRef{斐 1:27}。最后他说了这句话：\BibleQuote{你们该怀有基督耶稣所怀有的同一心情：他虽具有天主的形体，并没有以自己与天主同等为应当把持不舍的，却使自己空虚，取了奴仆的形体。}\BibleRef{斐 2:5}。请你们留意，我恳求你们，并要警醒！因为正如一把双刃利剑，无论落在何处，即使身处千万方阵之中，也能轻易劈开并摧毁，因为它两面锋利，无物能挡其锋；圣神的话语也是如此。\BibleRef{希 4:12}。因为凭这些话，他击倒了亚历山大城的亚略、萨莫萨塔的保禄、加拉提亚的马尔塞鲁、利比亚的撒伯里乌、本都的玛尔西翁、瓦伦提努、摩尼、劳底西亚的亚波里拿留、福提诺、索夫罗尼乌的追随者，总而言之，所有异端。所以请你们警醒，以便目睹如此壮观的景象，如此众多的军队被一击而倒，免得你们失去目睹此情景的快乐。因为如果在赛马中，当一辆战车冲撞并颠覆整辆战车及其驭手，在击倒许多站立其上的战车手后，独自驶向终点和赛程的尽头，在四周响彻云霄的欢呼与呐喊中，仿佛被那喜悦与欢呼插上翅膀的战马驰过整个场地时，没有什么比这更令人愉悦的了；那么在这里，当藉着天主的恩宠，我们一下子整体推翻所有这些异端连同其驭手的联合与魔鬼般的阴谋时，快乐又该多么大呢？\par

如果你们同意，我们先将这些异端本身排列好。你们希望按他们不敬的程度排列，还是按时间顺序？按时间顺序，因为难以判断他们不敬的程度。那么首先让利比亚的撒伯里乌上前。他主张什么？他说：父、子和圣神只是一个位格的三个名字。本都的玛尔西翁说：创造万有的天主不是良善的，也不是良善的基督的父，而是另一位正义者，并且基督并没有为我们取了肉身。马尔塞鲁、福提诺和索夫罗尼乌声称：圣言是一种能力，正是这能力住在达味后裔中的人身上，而不是一个有位格的实体。\par

亚略确实承认圣子，但只是口头上的；他说圣子是受造物，远逊于父。还有人说圣子没有灵魂。你们看见战车站立了吗？现在看他们的跌倒，他如何一下子全部推翻他们。怎么推翻的？他说：\BibleQuote{你们该怀有同一心情，}\BibleQuote{这心情在基督耶稣内，他虽具有天主的形体，并没有以自己与天主同等为应当把持不舍的。}萨莫萨塔的保禄倒了，马尔塞鲁和撒伯里乌也倒了。因为他说：\BibleQuote{具有天主的形体。}如果\BibleQuote{具有天主的形体}，邪恶的人啊，你怎么能说他是从玛利亚取得起源，以前并不存在？你怎么能说他是能力？因为经上写着：\BibleQuote{天主的形体取了奴仆的形体。}\BibleQuote{奴仆的形体}是奴仆的能力，还是奴仆的本性？我想，当然是奴仆的本性。同样，天主的形体就是天主的本性，因此并不是能力。看，加拉提亚的马尔塞鲁、索夫罗尼乌和福提诺也倒了。\par

看，撒伯里乌也倒了。经上写着：\BibleQuote{他没有以自己与天主同等为应当把持不舍的。}平等不能指同一个位格，因为平等的事物需要与之平等的事物。你们没有看见两个位格的实体，而不是没有实体的空名吗？你们没有听见独生子的永恒先在吗？\par

最后，我们该怎样反驳\Person{亚略}，他断言圣子的本质不同呢？请告诉我，\Quote{他取了仆人的形象}是什么意思？意思是他成了人。因此，\Quote{既然具有天主的形象}，他就是天主。因为既然提到了一个\Quote{形象}和另一个\Quote{形象}；如果一个是真实的，另一个也是真实的。\Quote{仆人的形象}意思是人性，因此，\Quote{天主的形象}意思是天主性。他不仅为此作证，也证明了他的平等，正如\Person{若望}也作证，他决不低于圣父，因为他说：\Quote{他不认为与天主同等是应当强夺的}。现在他们那聪明的推理是什么呢？不，他们说，他证明的恰恰相反；因为他说，\Quote{既然具有天主的形象，他没有强夺与天主同等}。如果他本是天主，他怎能\Quote{强夺}呢？这不是毫无意义吗？谁会说出一个人没有强夺做人呢？因为人怎能强夺他本来就是的呢？不，他们说，他的意思是：作为一个较小的神，他没有强夺与那比他大的大神平等。有大小之分的神吗？你们是把希腊人的教义带进教会的教义里吗？在他们那里，有大小神之分。如果你们那里也如此，我不知道。因为在圣经里你绝找不到：你在那里从头到尾只会找到全能的天主，没有一个地方说是小的。如果他小，他怎能又是天主呢？如果人不是又小又大，而只有一个本性，凡不属于这个本性的就不是人，那么怎能有一个小的神和一个大的神呢？\par

不属于那本性的，就不是天主。因为圣经处处称祂为大；\BibleQuote{主是伟大的，应受赞美。} \BibleRef{咏 47:1}这话也是关于圣子说的，因为圣经常称祂为主。\BibleQuote{祢是伟大的，行奇妙的事。只有祢是天主。} \BibleRef{咏 85:10}又说：\BibleQuote{我们的主是伟大的，祂的能力伟大，祂的伟大不可测量。} \BibleRef{咏 144:3}\par

但是圣子，他说，是小的。但这是你说的，因为圣经说的相反：如同论圣父一样，也论圣子说；请听\Person{保禄}说：\BibleQuote{期待那蒙福的盼望，和我们的伟大天主的荣耀显现。} \BibleRef{铎 2:13} 但他能说\Quote{显现}是指圣父吗？不然，为更能说服你，他加了一句关于显现\Quote{这位伟大的天主}。那么，难道不是论圣父说的吗？绝不是。因为下文不容许这样，它说：\BibleQuote{我们的伟大天主和救主耶稣基督的显现}。你看，圣子也是伟大的。那么你怎么说大小呢？\par

也请听先知称祂为\BibleQuote{伟大的谋士}。\BibleRef{依 9:6} \BibleQuote{伟大的谋士}，祂自己不伟大吗？\BibleQuote{大能的天主}，祂是小的而不伟大吗？那么这些无耻狂妄的人说祂是小神时，是什么意思？我常常重复他们的话，好让你们更远离他们。祂作为较小的神，没有为自己强夺与那较大的神相似！请告诉我（但不要以为这是我的话），如果祂如他们所说的小，而且权力远次于圣父，祂怎能强夺与天主平等呢？因为一个较低的本性不能强夺进入那较高的本性；例如，一个人不能强夺成为天使的本性；一匹马，即使它愿意，也不能强夺成为人的本性。但除此之外，我还要说这一点。\Person{保禄}用这个例子想要建立什么？当然，是引导\Place{斐理伯}人走向谦卑。那么他提出这个例子是为了什么目的？因为没有一个劝人谦卑的人会这样说：\Quote{你要谦卑，在同等尊荣的人中看自己更卑微，因为奴隶没有起来反抗主人；你要效法他。} 任何人都会说，这不是谦卑，而是傲慢。学习什么是谦卑吧，你们这些有魔鬼般骄傲的人！那么谦卑是什么？是谦虚的心。谦虚的心是那自己谦卑的人，不是那因不得已而谦卑的人。解释一下我所说的，请你们注意：谦虚的心，当它有能力变得高傲时，却是谦卑的；但那因为不能高傲而谦卑的，不再算谦虚。例如，如果一个国王顺服于自己的官员，他是谦卑的，因为他从高位降下；但如果一个官员这样做，他就不是谦虚的心；因为怎么呢？他没有从任何高位上降下自己。除非我们能够做到相反，否则无法表现出谦虚的心。因为如果我们必须违背自己的意愿谦卑，那优点并非来自精神或意志，而是来自必然。这德行被称为谦虚的心，因为是心灵的谦卑。\par

如果有人没有能力夺取别人的财物，而继续占有自己的财物，你以为我们应当为他的正义而称赞他吗？我想不会。为什么呢？因为选择自由的可称赞之处被必然性剥夺了。如果某人没有能力篡权为王，而只是一个平民，我们应当为他的安静而称赞他吗？我想不会。同样的原则也适用在这里。因为，最无知的人哪，称赞不是因戒绝这些事而给予的，而是因行善；前者固然免于责备，但尚未有称赞，后者才配得颂扬。请看，基督正是本着这一原则而称赞的，祂说：\BibleQuote{我父所祝福的，你们来吧！承受自创世以来为你们预备的国度，因为我饿了，你们给了我吃的；我渴了，你们给了我喝的。}\BibleRef{玛 25:34}祂没有说：因为你们没有贪心，没有抢劫；这些都是小事；而是说：\Quote{你们看见我饿了，给了我吃的。}有谁曾这样称赞他的朋友或敌人呢？没有人曾这样称赞保禄。何必说保禄呢？没有人曾这样称赞一个普通人，如同你称赞基督一样，因为祂没有夺取那不属于祂的权柄。因这样的事而钦佩，正是许多邪恶的证据。为什么呢？因为在恶人当中，这乃是可称赞的事，如同一个贼，若不再偷窃，就被称赞；但在善人当中则不然。\BibleRef{弗 4:28}一个人没有夺取不属于他的权柄和尊荣，他就可称赞吗？这是何等愚蠢！\par

我恳求你们注意，因为这一推论很长。再者，谁会从这样的理由出发来劝人谦卑呢？榜样应当比我们所劝勉的主题大得多，因为没有人会被与主题无关的事所感动。例如，当基督引导我们善待仇敌时，祂带来了一个伟大的榜样，即祂的父的榜样，\BibleQuote{因为祂使太阳上升，光照恶人，也光照善人；降雨给义人，也给不义的人。}\BibleRef{玛 5:45}当祂引导人忍受冤屈时，祂带来了一个榜样，\BibleQuote{你们跟我学习，因为我是良善心谦的。}\BibleRef{玛 11:29}又说：\BibleQuote{我为主、为师傅的，尚且这样做，何况你们呢？}\BibleRef{若 13:14}你们看，这些榜样并非遥远，因为不需要如此遥远，因为我们确实也做这些事；尤其在此处，榜样甚至不是相近的。怎样呢？如果祂是仆人，祂就较为低微，且服从那更大的，但这并非谦卑。需要表明相反的，即那更大的使自己服从那较小的。但由于在天主身上他找不到这种\Quote{更大}与\Quote{较小}的区别，他至少建立了一种平等。如果圣子是低微的，这就不是足以引导我们谦卑的榜样。为什么呢？因为那不是谦卑：那较小的不反对那更大的，不夺取权柄，且\Quote{服从至死}，这并非谦卑。\par

再者，请看他如何在榜样之后说：\BibleQuote{存心谦卑，各人看别人比自己强。}\BibleRef{斐 2:3}他说\Quote{看}，因为你们在本质上是一样的，在来自天主的尊荣上也是一样的，所以这事关乎看法。在那些较大和较小的人之间，他不会说\Quote{看}，而是说尊敬那些比你们强的，正如他在另一处所说：\BibleQuote{你们要信服并顺从那些领导你们的人。}\BibleRef{希 13:17}在那例子中，服从是事情本质的结果；在此处，却是出于我们自己的判断。他说：\BibleQuote{存心谦卑，各人看别人比自己强}，正如基督所做的一样。\par

这样，他们的解释就被推翻了。现在，在首先简要地论说了他们的观点之后，我该谈谈我们自己的观点了。当劝人存心谦卑时，保禄绝不会提出一个较小的、服从更大的例子。如果他是劝仆人服从主人，他可以恰当地这样做；但当他劝自由的人服从自由的人时，提出仆人服从主人（即较小的服从较大的）的例子有什么用处呢？他不说：\Quote{让较小的服从较大的}，而是说你们这些彼此同样尊荣的人要彼此服从，\BibleQuote{各人看别人比自己强}。那么，他为什么连妻子的服从也没有提出，说：如同妻子服从丈夫，你们也该如此服从？既然在那种有平等和自由的情况中，服从是轻微的，他尚且没有提出，那么他更不会提出奴隶的服从了。我前面说过，没有人因戒恶而称赞一个人，甚至根本不提他；没有人想要称赞一个人的节制时，会说：他没有犯奸淫，而是说：他戒绝了自己的妻子；因为我们根本不把戒恶看作可称赞的事，那将是可笑的。\par

我说过，\OriginalTerm{奴仆的形体}{form of a servant}是真实的形体，不是虚假的。因此，\OriginalTerm{天主的形体}{form of God}也是完全的，不是虚假的。为什么他不说\BibleQuote{成为天主的形体}，而说\BibleQuote{有天主的形体}呢？这与\BibleQuote{\OriginalTerm{我是自有者}{I am that I am}}这句话相同。\BibleRef{出 3:14}\Quote{形体}意味着不变，就其是形体而言。同实体的事物不可能有另一事物的形体，如同人没有天使的形体，兽也没有人的形体。那么，圣子又怎能呢？\par

在我们自身，因为我们人是复合的本性，形体属于身体；但在那单纯且完全非复合的本性中，形体就是实体本身。但如果你争辩说，祂所说的不是指圣父，因为这个词没有冠词，那么在许多地方，尽管词没有冠词，却是指圣父。我何必说在许多地方呢？因为就在此处，祂说：\BibleQuote{祂不认为\OriginalTerm{与天主平等}{equality with God}是应当紧抓不放的}\BibleRef{斐 2:6}，用词时没有冠词，却是指天主圣父。\par

我本愿意加上我们自己的解释，但我怕会使你们的心思负担过重。同时，请记住以上所说的话，以驳斥他们；同时，让我们先根除荆棘，然后等荆棘被根除、土地稍得安歇之后，我们再撒上好种子；这样，土地摆脱了由此所沾染的一切恶，就能以全部力量接受神圣的种子。\par

让我们为所说的话感谢天主；让我们恳求祂赐予我们守护和保全这些话的恩宠，好使我们和你们都欢喜，并使异端者蒙羞。让我们祈求祂为接下来的事开启我们的口，使我们能以同样的热忱陈述我们自己的见解。让我们恳求祂赐予我们相称于信德的生活，使我们能为祂的光荣而生活，并使祂的名不致因我们而受亵渎。因为经上记载说：\Quote{祸哉，你们}，\Quote{因你们天主的名受亵渎。} \BibleRef{依 52:5} 因为如果我们有一个儿子（还有什么比儿子更属于我们自己的呢？），因此如果我们有一个儿子，并且因他而受亵渎，我们会公开弃绝他，远离他，不愿接纳他；那么何况天主，当祂有忘恩负义的仆人亵渎并侮辱祂时，祂岂不更要远离他们、憎恨他们吗？天主所憎恨并远离的人，除了魔鬼和邪魔，谁会接纳他呢？凡被邪魔所接纳的人，还有什么救恩的希望？还有什么生活的安慰？\par

只要我们是在天主手中，\Quote{没有人能从我手中把他们夺去} \BibleRef{若 10:28}，因为那手是强而有力的；但当我们从那只手和那帮助中坠落时，我们就丧亡了，就暴露无遗，轻易被夺去，有如\Quote{将倾的墙壁和摇摇欲坠的篱笆} \BibleRef{咏 61:3}；墙软弱时，任何人都容易跨越。不要以为我下面要说的只指耶路撒冷，而是指向所有人。那关于耶路撒冷所说的是什么？\Quote{现在我要为我的爱友唱一首爱歌，论及祂的葡萄园。我的爱友在肥沃的山丘上有一个葡萄园，我给它围上篱笆，挖了水池，栽种了索勒克的葡萄树，在园中筑了一座塔，也凿了一个酒榨，我期望它结出葡萄，它却长出了野荆棘。现在，耶路撒冷的居民和犹大人啊，请你们在我与我葡萄园之间判断。我为我葡萄园所做的，还有什么没有做到呢？我期望它结出葡萄，它为何长出野荆棘？现在我要告诉你们，我必向我的葡萄园怎样行：我必撤去它的篱笆，使它被吞噬；我必拆毁它的墙壁，使它被践踏。我必使我的葡萄园荒芜，不再修剪，不再锄挖，荆棘和蒺藜要长起来，如同在荒野之地。我也必命令云彩不再降雨在它上面。因为万军上主的葡萄园就是以色列家，犹大人是祂所喜爱的植物。我期望它结出正义，它却行不义，期望它发出公道，它却发出哀号。} \BibleRef{依 5:1} 这也指向每个灵魂。因为当天主爱了人，做了所有必要之事，而人却长出荆棘而非葡萄时，祂必撤去篱笆，拆毁墙壁，我们便成了被掠夺之物。因为请听另一位先知在哀歌中所说的：\Quote{你为何拆毁了它的篱笆，使所有过路的人都能摘取它？森林中的野猪践踏它，田野的野兽吞食它。} \BibleRef{咏 80:12} 前一处提到的是玛待人和巴比伦人，这里却未提他们，而是\Quote{野猪}和\Quote{独行的野兽}指魔鬼及其所有部下，因为他们性情凶残污秽。因为圣经为了显示他的贪婪，说：\Quote{如同咆哮的狮子巡行，寻找可吞食的人} \BibleRef{伯前 5:8}；为了显示他毒害、致命、毁灭的本性，称他为蛇和蝎子；\Quote{因为}祂说，\Quote{你们要践踏毒蛇和蝎子，并胜过仇敌的一切势力} \BibleRef{路 10:19}；为了显示他的力量和毒液，称他为龙；如经上说：这条龙，\Quote{你所造的在其中戏耍的} \BibleRef{咏 103:26}。圣经各处称他为龙、弯曲的蛇和毒蛇 \BibleRef{咏 74:13}；他是多褶的野兽，诡计多端，力量强大，他搅动一切，扰乱一切，颠覆一切。\BibleRef{依 27:1} 但不要害怕，也不要恐惧；只需警醒，他就如麻雀一般；\Quote{因为}祂说，\Quote{你们要践踏毒蛇和蝎子。} 只要我们愿意，祂就使他在我们脚下被践踏。\par

请看看，这是何等的鄙视，甚至何等的悲惨，看见他站在我们头上，而那本是赐给我们去践踏的。这从何而来？是出于我们自己。如果我们愿意，他就变得强大；如果我们愿意，他就变得弱小。如果我们谨慎自守，与作为我们君王的祂站在一起，他就蜷缩起来，在攻击我们时如同小孩一般。每当我们离开祂，他就大大膨胀，发出可怕的声音，咬牙切齿，因为他发现我们失去了最大的帮助。因为他不会接近我们，除非天主容许他；因为如果未经天主许可，他连进入猪群都不敢，更何况进入人的灵魂？但天主容许他，或是为了惩戒、惩罚我们，或是为了使我们更经得起考验，如同约伯的例子。你看见了吗？他没有到约伯那里，也不敢靠近他，而是战栗发抖？我何必提约伯呢？当他对付犹达斯时，他也不敢完全抓住他，进入他内，直到基督将他从神圣的团体中分离。他确实从外面攻击他，但不敢进入；但当他看见犹达斯从圣洁的羊群中被割离，他就以超过豺狼的猛烈扑向他，并且不放他，直至以双重的死亡杀死他。\par

这些事记载下来，是为了劝诫我们。知道十二人中有一个是叛徒，对我们有什么益处？有什么获益？有什么好处？很多。因为当我们知道他如何陷入这致命的阴谋时，我们就能警惕，免得自己也遭遇同样的事。他为何如此？由于贪爱金钱。他是个贼。为了三十块银钱，他背叛了主。他被这情欲冲昏了头脑，为了三十块银钱背叛了世界的\OriginalTerm{主}{Lord}。有什么比这疯狂更糟糕？对他而言，没有任何事物是等同的，没有任何事物是相称的，\Quote{万民在他面前好像虚无}\BibleRef{依 40:15}，他却为了三十块银钱背叛了他。贪爱金钱确实是一个残酷的暴君，可怕地使人神魂颠倒。一个人因醉酒而失去理智，不如因贪爱金钱而失去理智；不是因为疯狂和神智错乱，而是因为贪爱金钱。\par

请告诉我，你为什么背叛他？他呼召了你这个无名无闻的人。他使你成为十二人之一，让你分享他的教导，应许你无数美善，让你行奇迹，你与其他宗徒共享同一餐桌、同一行程、同一团体、同一交往。这些难道不足以约束你吗？你为何背叛他？你指控他什么，恶人？不如说，你从他手中没有领受什么美善？他知道你的心思，却没有停止尽他的本分。他多次说：\Quote{你们中有一人将要出卖我。}\BibleRef{玛 26:21} 他多次指出你，却仍然宽容你，尽管他知道你是这样的人，却没有把你赶出团体。他仍然容忍你，仍然尊重你，爱你如同真门徒，如同十二人之一，最后（噢，你的卑劣！）他拿了一条毛巾，用自己无玷的双手洗了你污秽的脚，即便如此，你仍不回头。你偷窃穷人的东西，为了不让你陷入更大的罪，他也容忍了。没有什么能说服你。即使你是一头野兽或一块石头，难道不会因这些恩待、这些奇迹、这些教导而改变吗？虽然你如此残忍，他仍呼召你，以奇事吸引你，你比石头更迟钝，他仍吸引你归向他。然而，你没有因任何一件事而变得更好。\par

你或许对叛徒的如此愚昧感到惊奇；因此要畏惧那伤害他的东西。他成为这样是由于贪吝，由于贪爱金钱。要切除这情欲，因为它滋生这些疾病；它使我们不虔敬，使我们不认识天主，尽管我们从他手中领受了无数恩惠。我恳求你，切除它，这不是普通的疾病，它知道如何产生无数致命的死亡。我们看到了他的悲剧。让我们畏惧，免得我们也陷入同样的罗网。因为这些事记载下来，正是为了我们不遭遇同样的事。因此所有福音圣史记载了这事，为约束我们。要远远逃避它。贪婪不仅在于贪爱大量金钱，也在于贪爱金钱本身。贪求超过我们所需，是严重的贪吝。是成堆的金子说服了叛徒吗？为了三十块银钱他背叛了主。你们不记得我以前说的吗？贪婪不在于贪多，而在于贪小。看，他为了一点金子，不，不是为了金子，而是为了银钱，犯下了何等大罪。\par

一个贪吝的人决不可能看见基督的面容！这是不可能的事之一。它是万恶的根源，如果一个人拥有一样邪恶就失去那荣耀，那带着这根源的人将站在何处？金钱的仆人不能成为基督的真仆人。基督亲自宣告这是不可能的。他说：\Quote{你们不能事奉天主而又事奉玛门}\BibleRef{玛 6:24}，又说：\Quote{没有人能事奉两个主人}\BibleRef{玛 6:24}，因为他们给我们下达相反的命令。基督说：\Quote{怜悯穷人}；玛门说：\Quote{连赤身露体者身上的东西也要剥夺。}基督说：\Quote{舍弃你所拥有的}；玛门说：\Quote{夺取你所未有的。}你看到对立，你看到争斗吗？人怎能轻易服从两者，而必须轻视一个？其实，这不需要证明吗？难道我们不是在实际中看到，基督被轻视，玛门被尊敬？你们没有觉察到吗？这些话本身就令人痛苦，更何况事情本身？但实际中并不那么痛苦，因为我们被疾病所控制。如果灵魂稍许洁净这疾病，只要它还留在此世，就能正确判断；但当它离开到别处，被热病抓住，沉溺于事物的快乐时，它的知觉就不清晰，它的法庭就不公正。基督说：\Quote{你们中不论是谁，如不舍弃他的一切所有，不能做我的门徒}\BibleRef{路 14:33}；玛门说：\Quote{夺取饥饿者口中的食物。}基督说：\Quote{遮盖赤身露体的人}\BibleRef{依 58:7}；另一个说：\Quote{剥夺赤身露体者。}基督说：\Quote{不要掩面不顾你的骨肉}\BibleRef{依 58:7}，和你家中的人；玛门说：\Quote{不要怜悯你的骨肉；即使你看见你的母亲或父亲穷困，也要轻看他们。}我为什么说父亲或母亲？他说：\Quote{连你的灵魂也要毁掉。}而他竟被听从！哀哉！那命令我们残忍、疯狂、野蛮之事的主，竟比那命令我们温柔、有益之事的主更被听从！为此预备了地狱；为此，有火；为此，有火河；为此，有永不死的虫子。\par

我知道许多人听我说这些话时感到痛苦，的确，我说这些话时也并非没有痛苦。但我何必说这些话呢？我宁愿我的讲论始终关乎天国，关乎安息，关乎安息的水滨，关乎青绿的草场，正如圣经所说：\BibleQuote{祂使我卧在青绿的草场上，领我走近幽静的水旁}\BibleRef{咏 22:2}，祂使我在那里安居。我宁愿谈论那地方，\BibleQuote{忧愁和叹息将要远遁}\BibleRef{依 51:11}。\par

我宁愿谈论与基督同在的喜乐，尽管这超越一切言语和一切理解。但我仍愿按我的能力谈论这些事。然而我该怎么办呢？不可能对一个患病发烧的人谈论王国；那时我们必须谈论健康。不可能对一个受审的人谈论尊荣，因为那时他渴望的是从审判、刑罚和惩罚中解脱。如果这事不成，另一事又如何能成？正因如此，我不断地谈论这些事，好使我们能更快地转向那些事。正因如此，天主以地狱为威胁，好使无人堕入地狱，使我们所有人都能获得天国；正因如此，我们也常常提及地狱，好将你们推向天国，当我们以恐惧软化你们的心时，引领你们作出与天国相称的行为。不要因我们言语的沉重而不悦，因为这些言语的沉重使我们的灵魂从罪恶中得轻省。铁是沉重的，锤子也是沉重的，但它们能制成合用的器皿，无论是金的还是银的，并能矫正弯曲之物；如果它们不沉重，就无法矫正扭曲的实质。同样，我们沉重的言语有能力使灵魂恢复其适当的音调。那么，我们不要逃避沉重的言语，也不要逃避它给予的打击；这打击不是要粉碎或撕裂灵魂，而是要矫正它。我们知道如何打击，知道如何因天主的恩宠施加打击，不至于压碎器皿，而是打磨它，使它笔直，合乎主用，在那一日火河之前，将它呈现得光亮完好、精工巧作，无需那焚烧的柴堆。因为如果我们在此世不把自己置于火中，那么我们在那里必然被焚烧，不可能有别的结果；\BibleQuote{因为主的日子要在火中启示}\BibleRef{格前 3:13}。你们被我们的言语焚烧片刻，总比永远在那火焰中焚烧更好。这确实会发生，是显而易见的，而且我已屡次向你们提出无可辩驳的理由。我们本当真正信服圣经，但既然有些人好争辩，我们也从理性引出了许多论据。现在提一下它们也无妨，它们是什么呢？天主是公义的。我们所有人都承认这一点，无论希腊人、犹太人、异端者还是基督徒。但许多罪人离世时未受惩罚，许多义人离世时却遭受了千万苦难。那么，如果天主是公义的，祂将在哪里赏报前者之善，惩罚后者之恶呢？如果没有地狱，如果没有复活？所以，你们要不断向他人并向自己重复这个理由，它不会让你们不信复活；而不信复活的人必会谨慎生活，以求得永生之福，愿我们众人因我们的主耶稣基督的恩宠和慈爱，堪当获得这福分。愿光荣归于祂，等等。\par

\SourceRecord{Translated by John A. Broadus.；Nicene and Post-Nicene Fathers, First Series；Vol. 13.；Edited by Philip Schaff.；Buffalo, NY: Christian Literature Publishing Co.,；1889.；<http://www.newadvent.org/fathers/230206.htm>.}

% structure: p0001 source=h1 target=section reason=page-title

\section{斐理伯书讲道词第七篇}

\BibleRef{斐 2:5-11}\par

\BibleQuote{你们该怀有基督耶稣所怀有的心情：他虽具有天主的形体，并没有以自己与天主同等，为应当把持不舍的，却使自己空虚，取了奴仆的形体，与人相似，形状也一见如人；他贬抑自己，听命至死，且死在十字架上。为此，天主极其举扬他，赐给了他一个名字，超越其他所有的名字，致使上天、地上和地下的一切，一听到耶稣的名字，无不屈膝叩拜；一切唇舌无不明认耶稣基督是主，以光荣天主圣父。}\par

我已经陈述了异端者的观点。现在适宜谈谈我们自己的看法。他们说：\Quote{不以之为强夺的}，是指抢劫。我们已经证明，这完全是空洞无力的，也是不相干的，因为没有人会以此为由劝人谦卑，也不会以此称赞天主，甚至不会以此称赞人。那么，亲爱的，是什么意思呢？请留心我现在说的话。因为许多人认为，当他们卑下时，就被剥夺了应有的权利，降了格；保禄为消除这种恐惧，并表明我们不该如此受影响，便说：天主，那独生子，具有\OriginalTerm{天主的形象}{form of God}，丝毫不亚于圣父，与祂同等，\InnerQuote{不以\OriginalTerm{与天主同等}{equality with God}为强夺的}。\par

现在请明白这是什么意思。凡是人抢夺、违反权力所取得的，他不敢放下，因为害怕它会失落而脱离他的占有，而是紧紧抓住不放。那拥有与生俱来的尊严的人，并不害怕降下这个尊严，因为他确信不会发生这样的事。例如，阿贝沙隆篡夺了政权，此后就不敢放下。我们将举另一个例子，但如果例子不能把全部情况呈现在你面前，请不要见怪，因为这就是例子的本质：它们留下大部分内容让想象力去推理。一个人反叛他的君主，篡夺了王国：他不敢放下和隐藏这事，因为一旦隐藏，立刻就会失去。让我们再举一个例子：如果一个人暴力夺取某物，他会紧紧抓住不放，因为如果放下，立刻就会失去。总的来说，靠\OriginalTerm{掠夺}{rapine}获得东西的人都害怕放下、隐藏，或不敢不持续保持他们已占据的状态。但那些不是靠掠夺获得财产的人则不然，像人拥有理性存在的尊严。但在这里例子让我失效，因为在我们中间没有天生的优越性，没有一样好东西是天生属于自己的；但这些都内在于天主的本性中。那么一个人说什么呢？天主子不怕从祂的权利上降下，因为祂不把神性视为被\OriginalTerm{强夺之物}{prize}。祂不害怕任何人会剥夺祂的本性或权利，因此祂放下它，确信会再次拿起。祂隐藏了它，知道这样做并不会使祂变得低下。因此，保禄不说：\Quote{祂没有抢夺}，而是说：\Quote{祂不以之为强夺的}；祂不是靠抢夺拥有那地位，而是与生俱来，并非赐予，那是持久且安全的。因此，祂不拒绝取下位者的形象。暴君在战争中害怕脱下紫袍，而君王却非常安全地这样做。为什么呢？因为他的权力不是靠夺取来的。祂并没有像篡位者一样拒绝放下，而是因为祂天然拥有它，且永远不会与祂分离，祂便隐藏了它。\par

那些断言祂受了强制、祂是\OriginalTerm{受制}{subjected}的人在哪里呢？圣经说：\Quote{祂空虚了自己，贬抑自己，成了服从的，致死。}祂是如何空虚自己的？取了\Quote{奴仆的形象，成了人的模样，且形状一见如人。}经上写道：\Quote{祂空虚了自己}，是参照这段话：\Quote{各人看别人比自己强。}因为如果祂是受制的，不是出于自己的同意、自己的自由意志，那就不是谦卑的行为。如果祂不知道必须如此，祂就不完美。如果不知道时机，祂就不认识时机。但如果祂既知道必须如此，也知道何时该如此，为什么祂要服从受制呢？他们说，是为了显示圣父的优越。但这显示的不是圣父的优越，而是祂自己的卑下。难道圣父的名号不足以显示圣父的优先地位吗？因为除了祂，圣子拥有一切相同的事物。因为这尊荣不能从圣父转移给圣子。\par

那么异端者说什么？他们说：看，祂没有成为人。我指的是马西昂派。为什么呢？祂只是\Quote{成了人的模样}。但一个人如何能\InnerQuote{成了人的模样}？穿上影子吗？但那是幻影，不再是人模样，因为人的模样就是另一个人。那么你要如何回答若望，他说：\Quote{圣言成了血肉}？\BibleRef{若 1:14}但这位蒙福的同一位在另一处也说：\Quote{成了罪恶肉身的模样}。\BibleRef{罗 8:3}\par

\BibleQuote{在样式上如同人。}你们看，他们既说\BibleQuote{样式}，又说\BibleQuote{如同人}。如同人，并在样式上为人，并不是真正的人。在样式上为人，就不是按本性为人。你们看我多么坦率地陈明敌人的话，因为当我们不掩藏他们似乎强劲的论点时，那才是辉煌的胜利，且是充分的胜利。这乃是诡诈而非胜利。那么他们说什么呢？让我重复他们的论证。在样式上为人，就不是按本性为人；并且如同人，且在人的样式里，就不是人。那么，取奴仆的形体，就等于不取奴仆的形体。这里就有矛盾；你为何不首先解决这个难题？因为正如你们认为这反驳我们，我们也说那反驳你们。他不是说\BibleQuote{如同奴仆的形体}，也不是\BibleQuote{相似于奴仆的形体}，也不是\BibleQuote{在奴仆的形体之样式}，而是说\BibleQuote{祂取了奴仆的形体}。那么这是什么？因为有了矛盾。没有矛盾。愿天主禁止！他们的论证是冰冷而可笑的。他们说，祂取了奴仆的形体，是在祂拿毛巾束腰，洗门徒的脚之时。这是奴仆的形体吗？不，这不是形体，而是奴仆的工作。有奴仆的工作是一回事，取奴仆的形体是另一回事。为什么他不说\Quote{祂做了奴仆的工作}，那样更清楚？可是在圣经中，从来没有用\Quote{形体}代替\Quote{工作}，因为区别很大：一个是本性的结果，另一个是行动的结果。在日常用语中，我们也从不把\Quote{形体}叫做\Quote{工作}。此外，按照他们的说法，祂甚至没有取奴仆的工作，也没有束腰。因为如果一切都是幻影，就没有实在。如果祂没有真的手，怎能洗他们的脚？如果祂没有真的腰，怎能拿毛巾束腰？祂取了什么衣服？因为圣经说：\BibleQuote{祂拿起自己的衣裳}。\BibleRef{若 13:12}如此看来，就连工作也不曾真实发生，而全是欺骗，祂甚至没有洗门徒的脚。因为如果那无形体的本性没有显现，它就不在身体里。那么，是谁洗了门徒的脚？\par

再次，反对\Person{萨莫萨塔的保禄}要说什么？因为他肯定了什么呢？完全一样。但一个属于人性、仅是人的人，去洗他的同仆，这并不是空虚自己。因为我们反对亚略派所说的话，也必须反对这些人，因为他们彼此并无区别，只差一段时间；两者都肯定天主子是受造物。那么我们要对他们说什么呢？如果祂作为人洗了人，祂没有空虚自己，祂没有谦卑自己。如果祂作为人不坚持与天主同等，祂不配受赞美。天主成为人，是伟大、难以言喻、无法表达的谦卑；但那个人是个人，做人的工作，又有什么谦卑呢？哪里有天主的工作称为\Quote{天主的形体}？因为如果祂只是一个人，并且因祂的工作被称为天主的形体，我们为什么不也对\Person{伯多禄}这样做呢？因为他行了比基督自己更大的事迹。你们为何不说\Person{保禄}有天主的形体呢？为什么保禄不拿自己作例子？因为他行了上千个奴仆的工作，甚至毫不拒绝说：\BibleQuote{我们不是宣讲我们自己，而是宣讲基督耶稣为主，并因耶稣的缘故作你们的奴仆。}\BibleRef{格后 4:5}\par

这些都是荒谬和琐碎的！圣经说祂\BibleQuote{空虚自己}。祂如何空虚自己？请告诉我。祂的空虚是什么？祂的谦卑是什么？是因为行了奇事吗？这\Person{保禄}和\Person{伯多禄}也都行过，所以这并非圣子所特有的。那么，\BibleQuote{相似于人}是什么意思？祂有许多属于我们的东西，也有许多祂没有的；例如，祂不是由婚姻所生。祂没有犯罪。这些是祂有而没有人有的东西。祂不仅仅看起来是祂所是的，祂也是天主；祂看起来是一个人，但祂不像大众的人。因为祂在肉体上与他们相似。所以祂的意思是祂不仅仅是一个人。因此祂说\BibleQuote{相似于人}。因为我们确实是灵魂和身体，但祂是天主，也是灵魂和身体，因此祂说\BibleQuote{相似}。因为免得你们听到祂空虚自己，就以为这里有某种改变、堕落和损失；他说，当祂保持祂所是时，祂取了祂所不是的，并且\BibleQuote{成了血肉}后，祂仍是天主，因为祂是圣言。\BibleRef{若 1:14}\par

在这方面，祂与人相似，为此\Person{保禄}说：\BibleQuote{并样式上。}不是祂的本性堕落，也不是发生任何混淆，而是祂在样式上成为人。因为当祂说了\BibleQuote{祂取了奴仆的形体}之后，祂也大胆说了这个，看到前者会使一切反对者无言以对；因为当祂说\BibleQuote{在罪恶肉身的相似上}，祂不是说祂没有肉身，而是说那肉身没有犯罪，却相似于罪恶的肉身。相似在什么？在本性上，不在罪上，因此祂的肉身如同一个有罪的灵魂。如同在前一个情形下用了相似一词，因为祂并非在一切上平等，同样这里也有相似，因为祂并非在一切上平等，例如祂不是由婚姻所生，祂没有罪，祂不仅是一个人。祂说得对\BibleQuote{如同一人}，因为祂不是众人中的一个，而是\BibleQuote{如同}众人中的一个。那圣言，祂是天主，并没有降格为人，祂的实体也没有改变，而是显现为人；不是以幻影欺骗我们，而是教导我们谦卑。因此当祂说\BibleQuote{如同一人}时，这就是祂的意思；因为祂在别处也称呼祂为人，当祂说\BibleQuote{只有一个天主，在天主与人之间也只有一个中保，就是那人基督耶稣}时。\BibleRef{弟前 2:5}\par

到此为止是对这些异端者的驳斥。现在我必须针对那些否认祂拥有灵魂的人发言。倘若\Quote{天主的形体}是\Quote{完全的天主}，那么\Quote{奴仆的形体}就是\Quote{完全的奴仆}。再者，驳斥\OriginalTerm{亚略人}{Arians}。此处论及祂的天主性，我们不再见到\Quote{祂成了}、\Quote{祂取了}，而是\BibleQuote{祂空虚了自己，取了奴仆的形体，成了人的模样}；此处论及祂的人性，我们见到\Quote{祂取了}、\Quote{祂成了}。祂成了后者，祂取了后者；祂原是前者。因此让我们既不混淆也不分割本性。只有一位天主，只有一位基督，天主之子；当我说\Quote{一}时，我意指结合，而非混淆；一种本性并未退化成另一种，而是与之结合。\par

\BibleQuote{祂谦抑自己，服从至死，甚至死在十字架上。}你看，有人说，祂是自愿服从的；祂并不等同于祂所服从的那一位。你们这些顽固不悟的人！这丝毫未贬低祂。因我们也服从我们的朋友，但这并无影响。祂以圣子的身份服从圣父；祂并未因此陷入奴役状态，反而藉此行为尤其维护了祂奇妙的圣子身份，因祂这样极大地尊崇了圣父。祂尊崇圣父，不是要你们侮辱祂，而是使你们更钦佩祂，并从这个行为学到祂是真正的圣子，因祂尊崇圣父超过所有其他。没有人曾如此尊崇天主。祂有多崇高，就相应经历了多大的谦卑。因祂大于所有，无人与祂相等，同样，在尊崇圣父方面，祂超越了所有，不是出于必要，也非不情愿，这本身也是祂卓越的一部分；是的，言语无法表达。这确实是一件伟大而不可言喻的事：祂成了奴仆；祂经历了死亡，这更为伟大；但还有更伟大、更奇妙的；何以如此？并非所有死亡都一样；祂的死亡似乎是所有死亡中最可耻的，充满羞辱，是被诅咒的；因为经上记载：\BibleQuote{凡挂在树上的，都是可咒骂的。}\BibleRef{申 21:23}为此，犹太人也急切地要以这种方式杀害祂，使祂受辱，即使无人因祂的死亡而背弃祂，但或许他们会因祂死亡的方式而背弃。为此，两个强盗与祂同钉，祂在中间，使祂分享他们的恶名，并使经上所言得以应验：\BibleQuote{祂被列于罪犯之中。}\BibleRef{依 53:12}然而真理越发闪耀，越发光明；因为当祂的敌人如此图谋对抗祂的光荣，但光荣仍照耀出来，这事就更显伟大。他们以为不是通过杀害祂，而是通过以这种方式杀害祂，使祂成为可憎的，证明祂比所有人都更可憎，但他们一无所获。那两个强盗也是极为不虔敬的人（因为后来其中一人悔改了），以致即使在十字架上，他们还辱骂祂；对自己罪过的意识、当前的刑罚、以及自己同样受苦，都不能抑制他们的疯狂。因此一个对另一个说，并制止他：\BibleQuote{你连天主都不怕吗？你既然受同样的判决。}\BibleRef{路 23:40}他们的邪恶如此之大。因此经上记载：\BibleQuote{为此，天主极其举扬祂，赐给祂那超乎万名之上的名。}当圣保禄提及肉体时，他无畏地讲述祂所有的谦卑。因为直到他提及祂取了奴仆的形体，并且当他谈论祂的天主性时，看，他是多么崇高（崇高，我是指按他的能力；因为他并非按自己的配得而言，因他不能）。\BibleQuote{祂虽具有天主的形体，却没有以自己与天主同等为持守的宝物。}但当他说到祂成了人时，从此他便无畏地谈论祂的卑微，深信提及祂的卑微不会损害祂的天主性，因为祂的肉体容许了这点。\par

第9-11节。\BibleQuote{为此，天主极其举扬祂，赐给祂那超乎万名之上的名，使一切在天上的、地上的和地下的，因耶稣的名，屈膝叩拜；一切唇舌都明认耶稣基督是主，以光荣天主圣父。}我们要对异端者说：如果这是论及一位未降生成人的，即论及天主圣言，那么天主如何举扬祂？好像祂赐给祂比先前更多的东西？那么祂在这方面原是不完美的，并为了我们而成为完美！\Quote{赐给祂那名。}看，照你们所说，祂连个名字都没有！但如果祂是按应得而领受的，又怎会在此发现祂是藉恩宠、作为恩赐而领受的呢？那\Quote{超乎万名之上的名}——让我们看看，这名是怎样的？\Quote{因耶稣的名，}祂说，\Quote{一切屈膝叩拜。}他们（异端者）解释\Quote{名}为光荣。那么这光荣是超越一切光荣的，且这光荣简言之就是万有敬拜祂！但你们自外于天主的伟大，你们以为认识天主如祂认识自己一样，由此显然可见你们离天主的正确思想有多远。这从以下就可明白。告诉我，这是光荣吗？那么在世人受造前，在天使或总领天使受造前，祂未在光荣中？如果这是那超越一切光荣的光荣（因为这名是\Quote{超乎万名的}），那么即使祂先前已在光荣中，祂也是在低于此的光荣中。那么祂创造万物，就是为了被举扬到光荣中，不是出于祂自己的美善，而是因为祂需要从我们这里得到光荣！你们不见他们的愚昧吗？你们不见他们的不虔敬吗？\BibleRef{斐 2:9}\par

如今，倘若他们所说的是关于那降生成人的那一位，那倒有理由，因为天主圣言容许这样的话论及祂的肉躯。这并不涉及祂的神性，而是完全关乎救恩计划。\Quote{天上、地上和地底下的事物}是什么意思？意指整个世界，以及天使、人和魔鬼；或者指义人、活人和罪人。\par

\Quote{众口}要\Quote{承认耶稣基督是主，以光荣天主圣父。}也就是说，所有人都应如此说；而这正是圣父的光荣。你看见了吗？无论圣子在哪里受光荣，圣父也同样受光荣。同样，当圣子受羞辱时，圣父也受羞辱。若在我们身上尚且如此，尽管父子间的差别很大，但在天主身上，既然没有差别，光荣与羞辱就更会传递给祂。若世界服从圣子，这就是圣父的光荣。因此，当我们说祂是完美的、一无所缺、不次于圣父时，这就是圣父的光荣，因为祂生了这样的一位。这也是祂大能、良善和智慧的一大证明：祂生了一位在智慧与良善上毫不逊色的。当我说祂与圣父一样智慧、毫不逊色时，这是对圣父大智慧的证明；当我说祂与圣父一样大能时，这是对圣父大能的证明。当我说祂与圣父一样良善时，这是圣父良善的最大证据，因为祂生了这样一位（圣子），在各方面都不小于或次于祂自己。当我说祂所生的圣子在本性上不次于祂、而是相等，并且不是出于另一个本性时，我再次惊叹于天主，祂的大能、良善和智慧，因为祂向我们显现了另一位，出自祂自己，如同祂自己，除了祂不是圣父。因此，凡我对圣子所说的伟大之事，都传递给圣父。如果这件微小而轻松的事（因为世界敬拜祂，对天主的荣耀只是小事）是为了天主的荣耀，那么其他一切事情岂不更是如此？\par

那么，让我们为祂的光荣而相信，为祂的光荣而生活，因为二者缺一不可；当我们正确光荣祂，却生活不正时，我们便特别羞辱祂，因为我们归祂为师傅和导师，却藐视祂，不畏惧那可怕的审判宝座。外教人生活不洁并不奇怪；这并不值得如此谴责。但基督徒，分享如此伟大奥迹、享受如此光荣的人，竟然如此不洁地生活，这是最糟的，也是无法承受的。请告诉我：祂服从到底，因此接受了高处的荣耀。祂成了仆人，因此成为万有之主，包括天使和所有其他。所以，我们也不要以为我们谦卑时就是降低了自己。因为这样我们反而更能被高举；这是有理由的；那时我们才特别变得可敬。因为骄傲的人实在卑下，而谦卑的人被高举，基督的话已充分表明。然而，让我们审视事情本身。什么是谦卑？岂不是受责备、被控告、被诽谤？什么是被高举？就是受尊敬、受称赞、得荣耀。好吧。让我们看看事情如何。撒殚是一个天使，他高举了自己。结果如何？他不是比所有其他更卑下吗？他不是以大地为居所吗？他不是被所有人定罪和控告吗？保禄是一个人，他谦卑了自己。结果如何？他不是被钦佩吗？不是被称赞吗？不是被赞美吗？不是基督的朋友吗？他岂不行了比基督更大的事？他岂不常常如同指挥俘虏奴隶一样命令魔鬼？他岂不把他当作刽子手一样驱使？他岂不把他嘲笑？他岂不把魔鬼的头踩在脚下？他岂不大胆地求天主也让别人做同样的事？我为何这样说？阿贝沙隆高举了自己，达味谦卑了自己；哪一位被高举了？哪一位成了光荣的？还有什么比那位有福先知论及史米时所说的话更能证明谦卑呢？\BibleQuote{由他咒骂吧，因为上主吩咐了他。}\BibleRef{撒下 16:11}你若愿意，我们可亲自审视这些事例。税吏谦卑了自己，尽管这几乎不能称为谦卑；但如何？他以正确的心态回应。法利塞人高举了自己。结果如何？我们也来审视这些事。假设有两个人，都富有、极受尊敬、因智慧和能力以及世俗其他优势而显赫；然后其中一人寻求所有人的尊敬，若不得便发怒，要求超过应得的，并高举自己；另一人轻视这一切，不因此对任何人怀有恶意，并躲避别人给予的尊敬。因为除了逃避光荣，别无他法获得光荣；只要我们追求光荣，它就逃离我们；但当我们躲避它时，它就追逐我们。你若要有光荣，就不要贪求光荣。你若要崇高，就不要自高。况且，所有人都尊敬那不贪求光荣的人，而鄙视那追求光荣的人。因为人类天性多少有些好争辩，倾向于相反的情感。因此，让我们轻视光荣，这样我们便能成为谦卑，或者更确切地说，成为崇高。不要高举自己，好让你被他人高举；自己高举自己的人不被他人高举，自己谦卑的人不被他人谦卑。傲慢是一大恶，做一个愚人比傲慢更好；因为在前一种情况下，愚昧只是理智的错乱，而在后一种情况下则更糟，是愚昧加上疯狂：愚人是自己的祸害，但傲慢的人也是别人的祸害。这种悲惨来自无知。人不可能傲慢而不愚昧；充满愚昧的人便是傲慢的。\par

请听智者的话，他说：\Quote{我见自恃为智者的人，愚昧人比他更有希望。} \BibleRef{箴 26:12} 你看见吗？我没有无缘无故地说，我所说的邪恶比愚蠢更严重，因为经上记载：\Quote{愚昧人比他更有希望。} 因此，\Person{保禄}也说：\Quote{不要自作聪明。} \BibleRef{罗 12:16} 请告诉我，我们说哪些身体是健康的？是那些膨胀、内里充满空气和水分的，还是那些保持低平、表面显出克制的？显然，我们会选择后者。同样，灵魂也是如此：骄傲膨胀的灵魂比水肿病更糟糕，而克制的灵魂则摆脱了一切邪恶。那么，谦卑带来多少益处！你要什么？宽容？不轻易发怒？爱人之心？清醒？专注？所有这些善行都来自谦卑，而它们的反面来自傲慢：傲慢的人必然也是无礼的、好争吵的、易怒的、刻薄的、阴郁的，更像野兽而非人。你强壮并因此骄傲吗？你反而应当因此谦卑。你为什么为虚无之事骄傲？连狮子都比你勇敢，野猪比你强壮，你甚至比不上它们面前的一只苍蝇。强盗、盗墓者、角斗士，甚至你自己的奴隶，也许还有更愚蠢的人，都比你强壮。这是值得赞美的事吗？你为这样的事骄傲吗？你该羞愧地隐藏起来！\newline{}\par

你容貌俊美吗？这是乌鸦的夸耀！你并不比孔雀更美，无论颜色还是羽毛；鸟儿在羽毛上胜过你，在羽毛和颜色上远胜于你。天鹅也很美，还有许多别的鸟，如果与它们相比，你会发现自己一无是处。常常，无用的男孩、未婚少女、妓女和伪男子也有这种夸耀；这难道是骄傲的理由吗？你富有吗？从哪里来的？你有什么？金银宝石！这也是强盗、杀人犯、矿工的夸耀。罪犯的劳动成了你的夸耀！你装饰打扮自己吗？我们也可以看到马匹被打扮，在波斯人中骆驼也是如此，至于人，所有与舞台有关的人也是如此。难道你不为这些事夸耀感到羞耻吗？若没有理性的动物、奴隶、杀人犯、伪男子、强盗、盗墓者与你共享这些？你建造华丽的宫殿吗？那又如何？许多穴鸟住在更华丽的房子里，有更尊贵的居所。你没有看见许多对金钱疯狂的人在田野和荒僻之处建造房屋，成为穴鸟的居所吗？你因自己的声音骄傲吗？你无论如何也不能比天鹅或夜莺唱得更清晰、更甜美。是因你多种的艺术知识吗？但蜜蜂在这方面是何等有智慧！什么刺绣工、画家、几何学家能仿效它的作品？是因你衣服的精致吗？但蜘蛛在这方面胜过你。是因你脚步的敏捷吗？再次，第一奖赏属于没有理性的动物，如野兔、瞪羚，以及其他不缺脚步敏捷的野兽。你旅行了很多地方吗？并不比鸟儿更多；它们的迁徙更容易，它们不需要路上的食物，也不需要驮兽，因为它们的翅膀足够用；这是它们的船，这是它们的驮兽，这是它们的车，这也是它们风，总之，人能命名的一切。你视力敏锐吗？不如瞪羚，不如鹰。你听觉灵敏吗？驴更灵敏。嗅觉？猎犬不让你超过它。你善于储备吗？但你不如蚂蚁。你收集金子吗？但不如印度的蚂蚁。你因健康骄傲吗？没有理性的动物在身体状况和独立性上都比我们好得多；它们不怕贫穷。\Quote{你们看天空的飞鸟，不播种，也不收割，也不收集在仓里。} \BibleRef{玛 6:26} \Quote{的确，} 祂的意思是，\Quote{天主并没有使无理性的动物比我们优越。} 你注意到这里缺乏思考吗？你观察到完全缺乏探究吗？你观察到探究这些点带来的巨大益处吗？那心中高过众人的人，竟然被发现比无理性的动物还低下。\newline{}\par

但我们必会怜悯他，不效法他的榜样；也不因我们必死本性的界限容不下他的自高自大，便将他的地位降低到无理性的禽兽之列，反要从那里提升他，并非为他自己的缘故（他不配更好的命运），而是为彰显天主的慈爱和祂赐予我们的尊荣。因为有些事，是的，有些事是无理性的动物不能与我们共享的。那是哪些事呢？虔敬和基于德行的人生。在这里你绝不能说淫乱者、伪男子、杀人犯，因为我们已与他们分离。那么这里有什么呢？我们认识天主，承认祂的眷顾，并对永生怀有真正的智慧。于此，让无理性动物退让。它们在这些点上不能与我们争竞。我们生活在自制中。于此，无理性动物与我们毫无共同之处。因为，虽然我们在所有方面都落后于它们，却管辖它们；因为我们管辖的优越性就在这里：我们虽然落后于它们，却仍统辖它们；好叫你明白这些事的原因不在于你自己，而在于创造你、赐给你理性的天主。我们为它们设置网罗和陷阱，驱赶它们，它们任我们处置。\newline{}\par

自制、顺从的性情、温和、轻视钱财，都是我们人类的天赋；但你既是被傲慢蒙蔽之人，这些一样也没有，那么你持有超出人类水平或低于没有理性之受造物的观念，倒也合适。因为愚昧和鲁莽的本性就是如此：要么过分抬高，要么过分贬低，从不过度。在这方面我们与天使平等：我们被应许了一个王国，就是基督所加入的行列。一个人可能被鞭打，却不会屈服。一个人嘲笑死亡，不懂恐惧和战栗，他不贪图超过他所有的。因此，所有不像这样的人就低于没有理性的动物。因为当你在身体之事上想要占优势，却在灵魂之事上没有任何优势时，你除了比没有理性的动物更低劣之外，还能是什么呢？请举出一个邪恶且无思想的人，一个生活在放纵和自私中的人。马在好战精神上超过他，野猪在力量上，野兔在敏捷上，孔雀在优雅上，天鹅在声音的优美上，大象在体型上，鹰在视力的敏锐上，所有的飞鸟在财富上。那么你从何处获得统治没有理性受造物的资格呢？从理性吗？但你却没有它？因为凡停止妥善使用它的人，就再次变得比它们更低劣。因为当他虽然拥有理性却比它们更无理性时，倒不如他一开始就没有能够运用理性。因为接受了统治权后辜负信任，与错过接受它的时机，并不是同一回事。那位比他的卫兵水平还低的君主，倒不如从未穿过紫袍。在这个情况下也是如此。既然知道没有德行我们就低于没有理性的动物，就让我们在其中操练自己，使我们能成为人，是的，甚至成为天使，并享受所应许的福分，藉着我们的主耶稣基督的恩宠和慈爱，愿光荣归于祂，等等。\par

\SourceRecord{Translated by John A. Broadus.；Nicene and Post-Nicene Fathers, First Series；Vol. 13.；Edited by Philip Schaff.；Buffalo, NY: Christian Literature Publishing Co.,；1889.；<http://www.newadvent.org/fathers/230207.htm>.}

% structure: p0001 source=h1 target=section reason=page-title

\section{斐理伯书讲道集 第八讲}

\BibleRef{斐 2:12-16}\par

\Quote{所以，我亲爱的，你们既然常常听从，不但在我面前，如今不在我面前时更加顺服，就应当恐惧战栗做成你们自己的救恩；因为是天主在你们内运作使你们愿意和实行，为成就祂的美意。你们做一切事都不可发怨言和争论：好使你们无可指摘和纯洁，在这弯曲悖谬的世代中作天主无瑕疵的儿女；你们在其中好像世上的光，持守生命的话，使我在基督的日子有可夸耀的。}\par

我们所给的劝告应伴随称赞；因为这样当我们将受劝告者引向他们自己所表现的热忱时，劝告就变得受欢迎了；正如保禄在此所做的一样，且看他多么有智慧地说：\Quote{所以，我亲爱的，}他说；他没有简单地说\Quote{要听从}，直到他先用这些话称赞他们：\Quote{你们既然常常听从}；也就是说，\Quote{我请你们效法的不是别人，而是你们自己。} \Quote{不但在我面前，如今不在我面前时更加顺服。} 为什么呢？\Quote{更加不在我面前？} \Quote{你们那时或许是为了尊重我，出于羞耻而做一切事，现在却不再是这样；如果你们证明现在更加努力，那么也证明那时也不是为了我的缘故，而是为了天主的缘故。} 告诉我，你要什么？\Quote{不是你们听从我，而是你们\InnerQuote{以恐惧战栗做成自己的救恩}}；因为一个没有恐惧的人不可能树立任何崇高或令人信服的榜样；他所说的不只是\Quote{以恐惧}，而是\Quote{以战栗}，这是过度的恐惧。保禄就有这样的恐惧，因此他说：我害怕\Quote{免得我传给别人，自己反被弃绝。} \BibleRef{格前 9:27} 因为若没有恐惧，现世的事尚且不能成就，何况属灵的事呢？因为我想知道，有谁不带着恐惧学会字母呢？有谁不带着恐惧精通任何技艺呢？但是，当魔鬼不阻拦，只有懒惰是唯一障碍时，尚且需要如此多的恐惧，只是为使我们能克服我们本性的懒惰；那么当有如此激烈的战争，如此巨大的阻碍时，我们怎么可能不靠恐惧而得救呢？\par

这恐惧如何产生呢？只要我们思想天主无所不在，听见一切，看见一切，不仅是所行所说的，连心中和灵魂深处的也都看见，因为祂是\Quote{快利刺入心思和意图的} \BibleRef{希 4:12}；我们若这样自我约束，就不会行、说或想任何邪恶的事。告诉我，如果你要不断站在一位统治者身旁，你难道不带着恐惧站在那里吗？既然站在天主面前，你怎么能大笑、向后仰，而不感到恐惧和战栗呢？切不要藐视祂的忍耐，因为祂之所以容忍，是要引你悔改。每当你吃的时候，要思想天主临在，因为祂临在；每当你准备睡觉或放纵情欲、抢劫他人、沉溺奢侈、或做任何事时，你就不会大笑，不会发怒。你若常存此念，就会常处于\Quote{恐惧和战栗}中，因为你正站在君王旁边。建筑师尽管有经验，尽管精通技艺，仍要以\Quote{恐惧和战栗}站着，免得从建筑物上坠落。你已相信，你行了无数善事，你登上了高处：要确保自己，站时要恐惧，并谨慎留意，免得从那里坠落。因为有许多属灵的邪恶种类企图把你摔下来。\BibleRef{弗 6:12} \Quote{当以恐惧事奉主，}他说，\Quote{并且用战栗向祂欢呼。} \BibleRef{咏 2:11} 欢呼与\Quote{战栗}怎能兼容呢？但这确是唯一的欢呼；因为当我们做些善事，且是那些配得上以\Quote{战栗}行事的人所当作的事时，我们才欢呼。\Quote{以恐惧战栗做成自己的救恩}：他说的不是\Quote{做}，而是\Quote{做成}，即带着极大的热忱和勤奋；但既然他说了\Quote{以恐惧和战栗}，请看他是如何解除他们的焦虑的：因为他说什么？\Quote{是天主在你们内工作。} 不要因我说了\Quote{以恐惧和战栗}而害怕。我说这话不是要你们灰心绝望，以为德行难以达到，而是要你们追求它，不至于在虚妄的事上耗尽自己；若如此，天主就会成就一切。你只管放胆，\Quote{因为是天主在你们内工作。} 既然如此，我们的责任就是保持一颗永远坚决、握紧而不放松的心。\Quote{因为是天主在你们内运作使你们愿意和实行。} \Quote{如果祂自己在我们内运作意愿，你怎么对我们劝勉呢？因为如果祂自己运作意愿，你对我说的话就没有意义：\InnerQuote{你们听从了}；因为我们并没有\InnerQuote{听从}；你说\InnerQuote{以恐惧和战栗}也没有意义；因为全出于天主。} 我对你们说\Quote{因为祂在你们内运作使你们愿意和实行，}不是为此，而是为解除你们的焦虑。如果你愿意，祂就会\Quote{在你内运作意愿。} 不要害怕，你不会被击败；诚心的愿望和成就都是来自祂的恩赐：因为哪里有了意愿，祂就会增加我们的意愿。例如，我愿意行些善事：祂成就了善事本身，并藉着它成就了意愿。或者他是在极度虔敬中如此说，正如他宣称我们的善行是恩宠的恩赐一样。\par

同样，当祂赐予这些恩赐时，并未将我们排除在自由意志之外，反而赋予我们自由意志，因此当祂说：\Quote{在我们内工作意愿}，祂并未剥夺我们的自由意志，而是表明，通过实际行善，我们大大增进了行善的意愿。因为行善生出行善，不行善生出不行善。你施舍了吗？你便更受激励去施舍。你拒绝施舍了吗？你便更加不愿施舍。你节制了一天吗？你便同样有激励去节制次日。你纵欲过度了吗？你便增加了自我放纵的倾向。\Quote{恶人陷入邪恶的深渊时，便轻视一切。}\BibleRef{箴 18:3}同样，当人陷入不义的深渊时，就会变成轻蔑者；而当人进入良善的深处时，则会加紧迫切。因为前者因绝望而放肆，后者因感受到众多恩惠而更加努力，唯恐丧失全部。\Quote{为了祂的美意}，祂说，也就是说，\Quote{为了爱}，为了取悦祂；目的是使祂所悦纳的事得以成就，使万事依照祂的旨意成就。在此，祂指明并提供了信靠的基础：祂必定在我们内工作，因为祂愿意我们按祂所期望的方式生活，而如果祂愿意，祂亲自为此在我们内工作，并必定完成之；因为祂愿意我们生活正直。你看出祂如何未剥夺我们的自由意志吗？\par

\Quote{你们做一切事，总不可发出怨言和争论。}魔鬼发现他无力使我们离开正行，便想用其他方法破坏我们的赏报。因为他趁机潜入骄傲或虚荣，若这些都不行，就潜入怨言，或若不如是，就潜入疑惧。现在请看\Person{保禄}如何扫除这一切。他讲论谦卑时所说的一切，都是为了推翻骄傲；他谈及虚荣，即：\Quote{不但在我的面前如此}；他在这里提到\Quote{怨言和争论}。但我想知道，为何他在对待\Person{格林多人}时致力根除这恶习，便提醒他们\Person{以色列子民}的事例，而在这里却只字未提，仅简单告诫他们？因为前者的祸患已经造成，所以需要更严厉的打击和更尖锐的斥责；但在这里，他是在给予劝诫以防患未然。因此，对于尚未犯罪的人，不需要采取严厉措施；正如在引导他们谦卑时，他没有引用\WorkTitle{福音}中骄傲者受罚的事例，而是如同从天主口中发出的警告；\BibleRef{路 16:23}并且他以自由人、纯正子女的身份对待他们，而非奴仆；因为在修德行善方面，正派慷慨之人受善人影响，而品性不良之人受恶人影响；前者因顾及荣誉，后者因畏惧惩罚。为此，他写信给\Person{希伯来人}时，举出\Person{厄撒乌}为例：\Quote{他为了一碗食物，出卖了自己长子的名分}\BibleRef{希 12:16}；又说：\Quote{若退避，我的心就不喜悦他。}\BibleRef{希 10:38}在\Person{格林多人}中，有许多人曾犯奸淫。因此他说：\Quote{免得我再来时，我的天主使我在你们面前自卑，并为许多从前犯罪、而未悔改他们所学的不洁、奸淫和放荡的人哀哭。}\BibleRef{格后 12:21}\Quote{使你们成为无可指摘}，他说，\Quote{和无害}；\Quote{作天主无瑕疵的儿女，在歪曲悖逆的世代中，你们好像世上的光，持守生命的话，使我在基督的日子可以夸耀。}你注意到他是在教导他们不要发怨言吗？可见怨言是留给无原则、无恩宠的奴仆的。请告诉我，有哪个儿子，在自己经管父亲的事务、为自己谋利时，竟会发怨言？他说，你要思想，你是在为自己劳作，你是在为自己积蓄；只有那些别人享受他们的劳动成果、别人收获果实而他们却背负重担的人，才会发怨言；但为自己积聚的人，为何要发怨言？因为他的财富没有增加？但事实并非如此。为何自愿而非被迫行动的人要发怨言？与其带着怨言做事，不如什么都不做，因为连事情本身也被破坏了。你难道没有注意到，在我们自己的家庭中，我们常常这样说：\Quote{这些事宁可从来不做，也胜过带着怨言去做}？而且我们往往宁愿不被某人还欠我们的服务，也不愿忍受他发怨言的麻烦。因为怨言是难以忍受的，极其难以忍受；它近乎亵渎。不然，那些人为何要受如此严厉的惩罚？这是忘恩负义的证明；发怨言者对天主忘恩负义，而对天主忘恩负义的人因而成为亵渎者。那时，如果曾有过连绵不断的患难、无休止的危险，那就确是如此；没有停顿，没有缓解；从四面八方压来无数恐怖；但现在我们享有深沉的平安，完全的宁静。\par

那么，为何要抱怨呢？是因为你贫穷吗？但请想想约伯。或者是因为你患病？那么，假如你像那位圣人一样，拥有同样多的美德与崇高的成就，却遭受这样的苦难，你又会如何？再想想他吧：他长久以来不停地生虫，坐在粪堆上刮自己的疮；因为经上记载说：\Quote{（过了许多日子，）他的妻子对他说：\Quote{你还要坚持多久，说\InnerQuote{我还稍等片刻，怀着期望}？你说一句反对主的话，死了吧！}} \BibleRef{约 2:9}但你的孩子死了？那么，假如你像他一样，所有孩子都惨遭不幸而死，你又会怎样？因为你知道，你很清楚，陪伴在病人身边，合上他的嘴，闭上他的眼，抚摸他的胡须，听他最后的言语，这些并不是小小的安慰；但那位义人却没有得到丝毫安慰，他们一下子全都失去了。\newline{}我还说什么呢？假如你自己奉命去杀死并献上你的独生子，亲眼看着他的身体被焚烧，就像那位有福的圣祖一样，那么当你筑起祭坛、放上木柴、捆绑孩子时，你又会作何感受？\newline{}但有人辱骂你？那么，假如你的朋友们前来安慰你，却说出像约伯的朋友们那样的话，你又会作何感受？因为现在，我们有无数的罪过，我们理应受责难；但在那种情况下，那位真实、正义、虔敬、远离一切恶行的人，却从朋友那里听到了与他所受指控相反的话。那么，请告诉我，假如你听到你的妻子用责备的口吻喊道：\Quote{我成了流浪者和仆人，从一处到另一处，从一家到另一家，等待太阳下山，好让我从包围我的苦难中休息。} \BibleRef{约 2:9}愚蠢的女人啊，你为什么这样说？难道你的丈夫该为这些事负责吗？不，是魔鬼。\Quote{你说一句反对天主的话，}她说，\Quote{然后死去吧}——如果那受打击的人当时就咒骂并死去，你又能得到什么好处呢？\newline{}没有一种疾病能比他的更糟糕，即使你列举成千上万种。他的病如此严重，以至于他不能再待在屋里或有遮盖的地方；所有人都放弃了他。如果不是他已无可救药，他绝不会坐在城外，比那些患麻风病的人更可怜；因为那些人还可以进屋，他们聚集在一起；但他露天过夜，赤身裸体在粪堆上，甚至连衣服也穿不住。为什么这样？因为那只会增加他的痛苦。因为他说：\Quote{我刮我的疮，融化了地上的土块。} \BibleRef{约 7:5}他的肉体不断生出疮和虫子。你看到我们每个人听到这些时是如何作呕的吗？如果听都受不了，看到岂不是更难以忍受？如果看到都难以忍受，亲身经历岂不是更加难以忍受？然而，那位义人确实经历了这些，不是两三天，而是很长时间，并且他没有犯罪，连嘴唇也没有犯罪。\newline{}你能向我描述一种比这更痛苦的疾病吗？难道不比失明更糟糕吗？他说：\Quote{我看我的食物如同腐臭的团块。} \BibleRef{约 6:7}不仅如此，而且对别人来说能带来休息的夜晚和睡眠，对他却没有丝毫缓解，反而比任何折磨更痛苦。请听他的话：\Quote{你为何用噩梦惊吓我，用异象恐吓我？若是早晨，我就说：\Quote{何时到晚上？}} \BibleRef{约 7:14}而他并没有抱怨。还不止这些，还有世上的名声；因为人们根据他所受的一切痛苦，立刻断定他犯下了无数罪行。正因为如此，他的朋友们这样劝他：\Quote{所以，你当知道，天主向你追讨的，比你应得的还少。} \BibleRef{约 11:6}因此他自己也说：\Quote{但现在比我年少的人嘲笑我，他们的父亲我曾不屑与我的牧羊犬同列。} \BibleRef{约 30:1}这难道不比多次死亡更糟糕吗？然而，当他四面被这样的洪流冲击，周围风暴可怕，乌云、雨水、闪电、旋风、水柱一起肆虐时，他自己却岿然不动，仿佛坐在那可怕而汹涌的波涛中央，如一片平静之中，没有一声抱怨从他口中发出；而且这是在没有恩宠的时代，在关于复活、关于地狱、关于惩罚和报应的事丝毫未曾启示之前。可是我们，听到了先知、宗徒和福音传道者对我们说话，面前有无数榜样，又接受了复活的喜讯，却仍然心怀不满，尽管没有人能说自己也遭遇了这样的命运。因为如果有人损失了钱财，却并没有失去那么多儿女；或者即使失去了，也许是因为他犯了罪；但约伯却突然失去了他们，在他献祭、正在服侍天主的时候。如果有人一次损失了同样多的财产（虽然这不可能），但他并没有全身生疮的痛苦，没有刮去覆盖他的体液；或者即使他也遭遇了同样的命运，他并没有遭到人们的指责和辱骂，而后者比我们所受的灾祸更伤人。因为当我们遭遇不幸时，如果有人来安慰鼓励我们，给予我们美好的希望，我们尚且沮丧，那么想想有人辱骂他是什么滋味。如果\Quote{我指望有人怜悯，却没有一个人；指望有人安慰，却找不着一个} \BibleRef{咏 68:20}描述了无法忍受的痛苦，那么有辱骂者代替安慰者，是多么大的加重啊！他说：\Quote{你们都是令人厌烦的安慰者。} \BibleRef{约 16:2}如果我们不断地在心中思考这些事，好好权衡，那么当我们将目光转向那位战士、那钢铁般的灵魂、那铜墙铁壁般的精神时，现世的一切苦难就绝不能扰乱我们的平安。因为他仿佛带着一副铜或石的身体，以高贵而坚定的精神面对一切遭遇。\par

将这些事放在心上，让我们凡事\Quote{毫无怨言，也无争议}。你面前有善工，却抱怨吗？为什么？难道你被迫如此？因为我知道有许多周围的人强迫你抱怨，他说。这暗示了他说\Quote{在这弯曲悖谬的世代中}，但值得钦佩的是，我们在遭受令人恼怒的挑衅时，仍不容许自己有这样的感受。因为星星也在夜间发光，在黑暗中闪耀，且不损害自己的美丽，甚至更加明亮；但当光明返回时，它们就不再那样闪耀。同样，当你弯曲中持守正直时，你也显得更加光辉。这值得我们钦佩，就是\Quote{无可指摘}；因为他们可能以此为借口，他预先写下了这一点。所谓\Quote{持守生命之言}是什么意思？即\Quote{注定要活，属于那些正在获得救恩的人}。注意他如何立即附上存留的赏报。他说，光体保持光的原则；你们也要保持生命的原则。\Quote{生命之言}是什么意思？拥有生命的种子，即拥有生命的凭据，持守生命本身，即\Quote{你们自己拥有生命的种子}，他称之为\Quote{生命之言}。因此，其余的人都是死的，因为这些话他暗示了这一点；否则那些人也就会持守\Quote{生命之言}。\Quote{好叫我在基督的日子有可夸耀的}，他说；这是什么？他说，我也在你们的善行中有分。你们的德行如此伟大，不仅拯救了你们自己，也使我显赫。你这有福的保禄，真是奇特的\Quote{夸耀}！你为我们受鞭打、被驱赶、被辱骂：因此他补充说，\Quote{在基督的日子，我没有空跑}，他说，\Quote{也没有徒劳劳碌}，但我始终有可夸耀的理由，他说，我没有空跑。\par

\Quote{并且，即使我被奠祭}。他没有说\Quote{即使我死}，也没有在写给弟茂德时这样说，因为在那裡他也用了同样的表达：\Quote{因为我已被奠祭}\BibleRef{弟后 4:6}。他既安慰他们关于自己的死亡，又教导他们要欣然接受为基督而死的死亡。他说，我成了好像奠祭和牺牲。有福的灵魂！他把带领他们归向天主称为祭献。献上一个灵魂比献上公牛好得多。他说，\Quote{那么，如果在这祭献之上，我再献上自己，我的死亡作为奠祭，我就喜乐}。因为当他这样说时，他暗示了这一点：\Quote{并且，即使我被奠祭在你们的祭献和事奉上，我也喜乐，并与你们众人一同喜乐；同样，你们也要与我一同喜乐}。你为什么与他们一同喜乐？你看他表明他们应当喜乐吗？一方面，我因成为奠祭而喜乐；另一方面，我因献上了祭献而与你们一同喜乐；\Quote{同样，你们也要与我一同喜乐}，因为我被献上；\Quote{与我一同喜乐}，指的是\Quote{那在我自身喜乐的}。所以义人的死不是流泪的理由，而是喜乐的理由。如果他们喜乐，我们也应当与他们一同喜乐。他们喜乐时我们哭泣是不合适的。\Quote{但是}，有人争辩说，\Quote{我们渴望往常的交往}。这只是一个借口和托辞；要证明这一点，请注意他所吩咐的：\Quote{你们要与我一同喜乐，并欢喜}。你渴望往常的交往吗？如果你自己注定要留在此地，你这样说还有道理，但若在短暂的时光后你将追上那已离去的，你所寻求的交往又是什么呢？因为只有当一个人永远与他分离时，才会思念另一个人的陪伴；但如果他走的路你也要走，你所渴望的交往又是什么呢？我们为什么不哀哭所有外出旅行的人呢？我们不是只哭一会儿，在第一或第二天后就停止了吗？如果你渴望与他往常的交往，只哭那么久吧。他说：\Quote{我所受的不是灾难，我甚至因到基督那里去而喜乐，而你们却不喜乐。}\Quote{与我一同喜乐}。让我们也这样，当我们看见一个义人死去时就喜乐，甚至更当看见一个极度邪恶的人死去时；因为前者要去领受他劳苦的赏报，后者则减轻了一些罪债。但有人说，如果他活着，也许他会改变。但如果真有改变的可能，天主绝不可能取走他。因为那安排一切事件以使我们得救的，若有人答应要中悦祂，祂为何不给他机会呢？如果祂任凭那些永不改变的人，更会任凭那些改变的人。那么，让我们从各方面除去悲痛的锐利，让哀号之声止息。让我们在一切事上感谢天主；让我们凡事不出怨言；让我们喜乐，并在一切事上中悦祂，好使我们也能获得将来的美善，藉着我们的主耶稣基督的恩宠和慈爱，愿光荣归于祂，等等。\par

\SourceRecord{Translated by John A. Broadus.；Nicene and Post-Nicene Fathers, First Series；Vol. 13.；Edited by Philip Schaff.；Buffalo, NY: Christian Literature Publishing Co.,；1889.；<http://www.newadvent.org/fathers/230208.htm>.}

% structure: p0001 source=h1 target=section reason=page-title

\section{《斐理伯书》讲道辞第九篇}

\BibleRef{斐 2:19-21}\par

\BibleQuote{但在主耶稣内我希望能尽快打发弟茂德到你们那里去，使我因知道你们的情况而得到安慰；因为我没有一个同心人，会真诚地关心你们的情况；因为他们各求自己的事，不求耶稣基督的事。}\par

他曾说：\BibleQuote{这些事反而使福音更见进展，以致我的锁链，在禁卫军全军和其余的人中，都显明是为基督的缘故。}\BibleRef{斐 1:12}又说：\BibleQuote{即使我当以你们的信德为祭品、服务和牺牲，我也喜欢。}\BibleRef{斐 2:17}他用这些话坚固他们。也许他们怀疑他先前的话只是为了安慰他们。那么怎么样呢？他说：\BibleQuote{我打发弟茂德到你们那里去。}因为他们渴望听到关于他的一切。为什么他不说：\BibleQuote{叫你们知道我的情况，}而说：\BibleQuote{使我知道你们的情况}？因为厄帕夫洛狄托会在弟茂德到达之前报告他的情况。因此他在后面说：\BibleQuote{但我认为必须打发厄帕夫洛狄托，我的兄弟}\BibleRef{斐 2:25}；但我希望了解你们的事情。很可能他因身体软弱，在保禄那里住了很久。所以他说，我希望\BibleQuote{知道你们的情况。}看，他把一切事都归于基督，甚至打发弟茂德也是如此，他说：\BibleQuote{我在主耶稣内希望，}就是说，我确信天主会为我促成这事，使我知道你们的情况而得到安慰。正如我因你们听见关于我的事——就是你们所祈祷的，福音有了进展，它的敌人蒙受羞辱，他们认为能伤害我的手段反而使我喜乐——而安慰了你们；同样，我也希望了解你们的事，使我知道你们的情况而得到安慰。在此他表明，他们应当为他的锁链而喜乐，并与他一致，因为这些锁链为他产生了极大的安慰；因为\BibleQuote{我也能得到安慰}这句话暗示，就像你们一样。\par

哦，他对\Place{马其顿}有多么渴望！他向得撒洛尼人同样证明了这一点，就如他所说：\BibleQuote{但我们弟兄们，对你们暂被夺去，等等。}\BibleRef{得前 2:17}在这里他说：我希望打发弟茂德，使我能知道你们的情况，这是极度关怀的证明：因为他自己不能与他们同在时，就打发他的门徒，因为他不能忍受，哪怕片刻，对他们的状况一无所知。因为他并非透过圣神的启示知道一切，对此我们能够看出一些理由；因为如果门徒们相信事情是这样，他们就会失去羞耻感，但如今由于期望隐瞒，他们更容易被纠正。他通过说\BibleQuote{我也能得到安慰}极高程度地引起了他们的注意，并使他们更加热切，这样，当弟茂德到来时，可能不会发现另一种情况，并向他报告。他本人似乎也以同样方式行事，当他延迟来到\Place{格林多}人那里，好使他们悔改；因此他写道：\BibleQuote{为宽免你们，我暂缓到\Place{格林多}去。}\BibleRef{格后 1:23}因为他的爱不仅显现在报告自己的情况，也显现在渴望了解他们的情况；因为这是一个关怀他人、为他们着想、始终为他们奋斗的灵魂所具有的特质。\par

同时，他也通过打发弟茂德来尊荣他们。你说什么？你打发弟茂德？为什么？因为\BibleQuote{我没有一个同心的人，}就是说，没有一个关心如同我关心的人，没有一个人会\BibleQuote{真诚地关心你们。}\BibleRef{斐 2:20}那么他那些与他同在的人中难道没有一个人吗？没有一个同心的人，就是没有一个人像他那样为你们挂念、为你们着想。他说，没有人会轻易选择为此目的做这么长的旅程。弟茂德是与我一起爱你们的那一位。因为我本可以打发别人，但没有人像他一样。这就是同心，爱门徒如同主人爱他们。\BibleQuote{谁，}他说，\BibleQuote{会真诚地关心你们，}即如同父亲。\BibleQuote{因为各求自己的事，不求耶稣基督的事，}\BibleRef{斐 2:21}他们求自己的舒适，自己的安全。他也这样写给弟茂德。但他为什么哀叹这样的事？为了教导我们听众不要同样跌倒，教导听众不要寻求免除劳苦；因为寻求免除劳苦的人，不是寻求基督的事，而是自己的事。我们应当准备面对一切劳苦，一切患难。\par

第22节。\BibleRef{斐 2:22} \BibleQuote{你们知道他的明证，他怎样为了福音与我一同服侍，如同儿子服侍父亲一样。}\par

而且我并非信口而言，他说：\BibleQuote{你们自己，}他说，\BibleQuote{知道，他怎样为了福音与我一同服侍，如同儿子服侍父亲一样。}\BibleRef{斐 2:22}然后他将弟茂德推荐给他们，有理有据，使他能受到他们的尊敬。他写信给\Place{格林多}人时也这样做，他说：\BibleQuote{所以，谁也不可轻视他，因为他做主的工，正如我也做。}\BibleRef{格前 16:10}他说这话不是为照顾他，而是为那些接待他的人，使他们能得到大赏报。\par

第23节。\BibleRef{斐 2:23}他说：\BibleQuote{因此，我希望立刻打发他，一旦我看清我的事将要如何，}就是说，当我看到我的处境，我的事将有怎样的结局。\par

第24节。\BibleRef{斐 2:24} \BibleQuote{但在主内我信赖，我自己不久也会到你们那里去。}\par

因此我打发他，并非好像我自己不会来，而是为了使我得知你们的情况而得到安慰，甚至在这期间我也不至一无所知。他说：\BibleQuote{但我信赖主。}看，他如何将所有事交托于天主，不凭己意说话。也就是说，若主愿意。\par

第25节。\BibleRef{斐 2:25} \BibleQuote{但我认为必须打发厄帕夫洛狄托，我的兄弟、同工和战友，到你们那里去。}\par

他也以同样称赞弟茂德的话打发这人与他们同去，因为他在这两点上称赞了他：第一，他爱他们，当他说\BibleQuote{他必真实地关心你们}时；第二，他在福音上蒙了悦纳。为此，他以同样的理由、同样的言辞也称赞这人：如何称赞呢？称他为弟兄、同工，并且不止于此，也称他为\BibleQuote{战友}，显示了他如何分担他的危险，并为他作证，正如他为自己作证一样。因为\BibleQuote{战友}比\BibleQuote{同工}更进一层；也许他是在平安的事上给予帮助，但在战争和危险中却不然；但说\BibleQuote{战友}，也显示了这一点。\par

第二十五节。\BibleQuote{派遣你们的使者，并供给我所需的人} \BibleRef{斐 2:25}这就是说，我将你们自己的给予你们，或者也许是你们的导师。他再次多加了许多关于他爱的话，说：\par

第二十六、二十七节。\BibleQuote{他切望你们众人，并且甚是忧苦，因为你们听见他病了。他实在是病到了快死的地步；然而天主怜悯了他，不但怜悯他，也怜悯了我，免得我忧上加忧} \BibleRef{斐 2:26-27}\par

在这里他指向更深的一点，表明厄帕夫洛狄托也十分清楚他们是如何爱他的。这在相爱的事上并非小事。你们知道他怎样病了，他说；他忧苦，因为在他康复后未能见到你们，使你们因他患病而有的忧愁得以解除。在这里他也给出了另一个迟延打发他们去的理由，并非由于任何懈怠，而是他留下了弟茂德，因为没有别人（正如他所写的，他没有\BibleQuote{同心合意的人}），而厄帕夫洛狄托则因他的疾病。他接着显示这是一场久病，耗费了许多时间，添上说\BibleQuote{他实在是病到了快死的地步}。你们看见保禄多么急切地为门徒切断一切轻视或蔑视的借口，以及一切认为他不来是因轻视他们的猜疑。因为没有什么比使门徒相信他的长上关心他、并为他忧伤更足以吸引门徒到长上那里，这是至爱的一部分。因为\BibleQuote{你们听见了}，他说，\BibleQuote{他病了；他实在是病到了快死的地步}。我并非在制造借口，请听下文。\BibleQuote{然而天主怜悯了他}。你说什么，异端者？这里写着，天主的怜悯留住了并带回了那快要离去的人。如果世界是邪恶的，那么让人留在邪恶中就不是怜悯。向异端者回答容易，但向基督徒说什么呢？因为他或许会质疑说：\BibleQuote{离开世界与基督同在，那是好得无比的}，怎么说他蒙了怜悯呢？我要问，为什么同一宗徒说\BibleQuote{然而我在肉身活着，为你们更是要紧的}？正如为他要紧，为这人也要紧，他将在以后带着更丰富的财富、更大的胆量去见天主。即使现在不如此，以后也会发生，但为那些一旦离开人世的人，赢取灵魂的事就终止了。在许多地方，保禄也按照听者的一般习惯说话，并非处处按照他天上的智慧；因为他必须对仍怕死的世俗人说话。于是他显示他怎样看重厄帕夫洛狄托，并由此为他赢得尊敬，说他的保全对自己如此有益，以致对厄帕夫洛狄托所施的怜悯也临到了他。此外，没有这一点，现世生命也是好的；若非如此，为何保禄将非命之死视为惩罚？正如他说：\BibleQuote{在你们中间有好些软弱的与患病的，死的也不少} \BibleRef{格前 11:30}。因为来世并非仅仅比邪恶状态好（那样的话它就不是好的），而是比好的状态更好。\par

\BibleQuote{免得我}，他说，\BibleQuote{忧上加忧}。因他死亡而有的忧愁，加上因他疾病而有的忧愁。由此他显示他何等珍视厄帕夫洛狄托。\par

第二十八节。\BibleQuote{所以我越发急速打发他去} \BibleRef{斐 2:28}。\BibleQuote{越发急速}是什么意思？即不拖延，不迟延，十分迅速，吩咐他放下一切，去到你们那里，使他得以免除忧愁；因为我们所爱的人的健康，我们听见后，不如看见他们时更喜乐，尤其当这事出乎意料时，正如厄帕夫洛狄托的情形。\par

\BibleQuote{所以我越发急速打发他去，叫你们再见他时，可以喜乐，我也可以少些忧愁}。\BibleQuote{少些忧愁}怎讲？因为如果你们喜乐，我也喜乐，他也喜乐于这样的喜乐，我便\BibleQuote{少些忧愁}。他没有说无忧愁，而是说\BibleQuote{少些忧愁}，以显示他的灵魂从不免于忧愁；因为他说\BibleQuote{有谁软弱，我不软弱呢？有谁跌倒，我不焦急呢？} \BibleRef{格后 11:29}这样一个人，何时能免于忧愁呢？这就是说，我现在摆脱了这一沮丧。\par

第二十九节。\BibleQuote{所以你们要在主里欢欢喜喜地接待他} \BibleRef{斐 2:29}\par

\BibleQuote{在主里}或指属灵地、极其热诚地，或更确切地，\BibleQuote{在主里}的意思是蒙主准许。以配得上圣人的方式接待他，如同圣人应当被欢喜接待一样。\par

这一切他都为他们做，不是为他的使者，因为行善者比受善者得益更多。\BibleQuote{且尊重这样的人}，就是以配得上圣人的方式接待他。\par

第三十节。\BibleQuote{因他为基督的工作，几乎至死，不顾性命，要补足你们供给我的不及之处} \BibleRef{斐 2:30}\par

这人是腓立比城公开差派的，他们来作保禄的仆役，并或许为他带来一些捐款；因为在这书信的末尾，他显示他也带钱来，当他说：\BibleQuote{我藉厄帕夫洛狄托收到了你们所送的} \BibleRef{斐 4:18}\par

这样看来，很可能当他到达罗马城时，发现保禄正处在极大的迫近危险中，以致那些惯常去探望他的人无法安全地这样做，反而因他们的探望自身也陷入危险；这通常发生在非常巨大的危险和君王极度的盛怒中（因为当有人触怒了君王，被投入监狱，并受到严密看守时，连他的仆人也无法接近；这很可能当时就发生在保禄身上）。而厄帕夫洛狄托，因其高贵的品性，藐视一切危险，以便能进去见他，服侍他，做一切所需的事。因此，他提出两个事实，借此为他赢得他们的尊重：一是他几乎冒死之险，他说，是为了我；二是在如此受苦时，他代表了他们的城邑，以致那因他冒险所得的报偿将算在派遣他的人身上，仿佛城邑派遣他作为使者一样，因而对他友善的接待并认可他所做的，更可被称为参与了他所敢行之事。他没有说\Quote{为了我}，而是为使他的话更可信，说\Quote{因为为了天主的工作}，因为他行事不是为了我，而是为了天主，\Quote{他几乎至于死}。那么如何？虽然因天主的眷顾他没有死，但他自己却不看重自己的性命，把自己交付给可能临到他的一切苦难，以便不松懈对我的服侍。如果他为了服侍保禄而舍身致死，那么他会为了福音的缘故更加忍受这事。或者更确切地说，这也算为了福音的缘故，即使是为保禄而死。因为殉道的冠冕不仅可以通过拒绝祭拜而得戴在我们头上，这样的原因也同样使死亡成为殉道，而且如果我可以说出令人惊讶的话，后者更甚于前者。因为那敢为较小原因面对死亡的人，必更愿为较大原因这样做。因此，当我们看见圣人们处在危险中时，我们不应顾惜自己的性命，因为不冒险就永远无法成就任何高贵的事，而是那预先为自己此世安全着想的人，必要从将来的安全中失落。\par

他说：\Quote{为补足你们服侍我的欠缺。}这是什么意思？城邑没有亲临，但藉着派遣他，城邑通过他完成了对我的一切服侍。因此他补足了你们服侍的欠缺，所以也因这个理由他理当享有许多尊荣，因为你们全体所应行的事，他代表你们实行了。这里他表明，处于安全中的人向处于危险中的人也有一种先行的服侍，因为他如此说\Quote{欠缺}和\Quote{服侍的欠缺}。你们看见宗徒的精神了吗？这些话不是出自傲慢，而是出于他对他们的极大关怀；因为他称这事为\Quote{服侍}和\Quote{欠缺}，使他们不致自高自大，而是保持谦逊，也不致以为他们行了什么大事，反而要谦卑自持。\par

因为我们对圣人欠有债务，并非在施恩于他们。正如那些处于和平、未参战的人有义务向那些从军作战的人提供补给（因为他们是代表这些人而立），此处也是如此。因为如果保禄没有教导，谁会把他投入监狱？所以我们应当服侍圣人们。难道不是荒谬吗？对世俗君王在作战时贡献他所需的一切，如衣服和食物，不仅按他的需要，而且丰裕地；而对天上的君王，当祂在作战、对抗远更凶恶的仇敌时（因为经上记载：\Quote{我们的战斗不是对抗血肉}\BibleRef{弗 6:12}），我们却不愿供应急迫的需要？这是何等的愚昧！何等的忘恩！何等的贪财！但看来，对人的恐惧在我们身上比地狱和将来的刑罚更有力量。为此，诚然，一切都颠倒了；因为世俗的事务每天都以极大的热忱被完成，一个都不能落后，而对属灵的事却全然不顾；而那些对我们来说是被迫的、强制的、如同奴隶般不情愿地被要求的事物，我们却非常乐意地交出；而那些出自甘心、如同自由人般被请求的事物，却又缺乏。我并非对所有人说话，而是对那些在供应上落后的人。难道天主不能使这些贡献成为强制的吗？但祂不愿如此，因为祂对你们的关心超过对你们所供养的人。因此祂不愿你们不得己而贡献，因为那样就没有报酬了。然而，这里站着的许多人比犹太人更低劣。想想犹太人奉献了多少大事：什一之物、初果、再次什一、又另加什一，此外还有第十三份、协刻耳，没有人说他们吞吃了多少；因为他们收得越多，赏报就越大。他们不说：他们收得多，他们是贪食者；这些话我现在从某些人那里听见。他们那边，在建造房屋、购置田产时，仍以为一无所有；但如果某位司铎穿着比平时更光鲜的衣服，享受超出维持生计所需之物，或有一个仆从，以免自己被迫做出不合宜的举动，他们就把这事定为富有。而事实上，即便是这样的程度我们也算富有，他们却勉强承认；因为我们虽然只有一点点，却是富有的；而他们虽然得到了一切围绕的东西，却是贫穷的。\par

我们的愚昧要到什么程度呢？难道我们不行善事还不够受罚，还必须加上恶言相向的惩罚吗？因为如果他所有的原是你所赐的，你因以赐予的事责备他，便丧失了你的赏报。简言之，你若赐予了，为何还要责备他？你已经以说他所有的都是你所赐的，见证了他是贫穷。那么，你为何还要责备呢？你若打算这样做，当初就不该施予。但若是别人施予时你这样说呢？那就更严重了，因为你本人没有施予，却责备别人的善行。你想这样被谈论的人会得到多么大的赏报？他们是为了天主的缘故而这样受苦。怎么、为什么？如果他们愿意，即使不从祖先那里得到，他们也可以从事商人的生活。因为我常听见许多人随口这样说，当我们说某人是穷人时，他们说：他若愿意，本可富有，然后讥讽地加上一句：他的父亲、祖父，我不知道还有谁，都是如此；但现在看他穿的那件袍子！可是，告诉我，他应当赤身露体地行走吗？你竟对这些事吹毛求疵，但小心不要因此而反对你自己。请听基督的劝诫，祂说：\BibleQuote{你们不要判断，以免受判断。}\BibleRef{玛 7:1}诚然，他若愿意，本可过商贾或买卖人的生活，必不至缺乏。但他不愿意。那么，有人说，他在这里有何益处呢？告诉我，他有何益处？他穿绫罗绸缎了吗？他傲慢地前呼后拥穿过广场吗？他骑马出行吗？他盖房子吗？既然已有住处。若他如此行事，我也指责他，毫不宽容，并声明他不配作司铎。因为连自己都无法劝诫之人，怎能劝诫他人不要在这些多余的事情上浪费时间呢？但若他已有足够的供养，难道就做错了吗？你难道要他过流浪的生活，去讨饭吗？你，他的门徒，难道不也感到羞耻吗？如果你肉身的父亲如此行，你会觉得此事可耻。如果你的神师被迫如此行，你难道不会掩面，甚至觉得要钻入地缝吗？经上记载：\BibleQuote{父亲的不名誉是子女的耻辱。}\BibleRef{德 3:11}但怎办呢？难道要他饿死吗？这不像是虔诚人的行为，因为天主不要这样。但他们立刻如何论断呢？他们说，经上记载：\BibleQuote{你们不要在腰带里备下金、银、铜钱，也不要两件内衣，也不要鞋，也不要棍杖。}\BibleRef{玛 10:9}而这些人却有三四件衣服，还有铺设好的床。我现在不得不痛苦地叹息，若非失礼，我甚至要哭泣！为什么？因为我们如此好奇地察看别人眼中的木屑，却感觉不到自己眼中的大梁。告诉我，你为何不将这话对自己说？回答是：因为这条命令只赋予我们的导师。那么，当保禄说：\BibleQuote{只要有吃有穿，就当知足。}\BibleRef{弟前 6:8}他只是对导师说的吗？绝非如此，而是对众人；这一点若我们往前追溯，就很清楚。因为他说什么？\BibleQuote{虔敬本是获益的根源，只要有知足的心。}\BibleRef{弟前 6:6}\BibleQuote{因为我们没有带什么到世上来，也确实不能带走什么。}\BibleRef{弟前 6:7}他随即加上说：\BibleQuote{只要有吃有穿，就当知足。但那些想致富的人，就陷于诱惑，陷入罗网和许多无知有害的欲望中。}\BibleRef{弟前 6:8}你看，这是对众人说的。再者，他另一处说：\BibleQuote{不要为肉体安排，去放纵私欲。}\BibleRef{罗 13:14}难道这不是绝对对众人说的吗？以及他说：\BibleQuote{食物是为肚腹，肚腹是为食物；但天主必要使这两样都归于无有。}\BibleRef{格前 6:13}或者他说：\BibleQuote{那好宴乐的寡妇，虽生犹死。}\BibleRef{弟前 5:6}这是论及寡妇。难道寡妇是导师吗？他自己岂没有说过：\BibleQuote{我不准许女人施教，或管辖男人。}\BibleRef{弟前 2:12}但若一位寡妇到了老年（年老需要大量照顾），何况女性本性（因为女性体质软弱，需要更多休息），若在年老和女性本性两者兼具的情况下，他不容许她过着奢侈的生活，甚至说她是死的（因为他并非简单禁止奢侈生活，而是说：\BibleQuote{那好宴乐的寡妇，虽生犹死}，这样就断绝了她（因为死人已被断绝）），那么，凡做那些连妇女和老者都要受罚之事的人，将得到怎样的宽宥呢？\par

然而，没有人思想这些事，没有人追究这些事。我不得不如此说，并非有意为司铎开脱这些指责，乃是为了顾惜你们。他们实在不会因你们而受损害，即使你们有理有据地指控他们贪得无厌；因为不论你们说或不说，他们终须向审判者交账，所以你们的言语丝毫不伤害他们；但若你们的言语又是虚假的，他们反而因这些诬告而得利，你们却因此伤害自己。但于你们却不是这样；因为无论你们对他们的指控是真是假，你们说他们的坏话都是伤害自己。怎么是呢？若是真的，你们因判断导师、扰乱秩序而伤害自己。因为我们若不可判断弟兄，何况导师呢？但若是假的，那惩罚和报应就无法忍受；因为\BibleQuote{人所说的每句废话，在审判之日都要还清。}\BibleRef{玛 12:36}所以我为你们的缘故这样行事、劳碌。\par

但是，如同我所说的，没有人探究这些事，没有人忙碌于这些事，没有人独自思量这些事。你们愿意我再加些吗？\Quote{凡不舍弃一切所有的，基督说，便不配作我的门徒。} \BibleRef{路 14:33} 当祂说：\Quote{富人进入天国是难的}？ \BibleRef{玛 19:23} 又当祂说：\Quote{你们富人有祸了，因为你们已经得到了你们的安慰}？ \BibleRef{路 6:24} 没有人探究这一点，没有人铭记于心，没有人自我推理，但所有人都像严厉的审问官一样坐在他人的案子上。然而，这使他们成了指控的同谋。但请听，为了你们自己，我愿使司铎们从你们所说的指控中得释放，因为你们说服自己他们违背了天主的法律，这使你们不无朝着邪恶倾斜。来吧，让我们查考这件事。基督说：\Quote{不要带金子，也不要带银子，也不要带两件内衣，也不要带鞋子，也不要带腰带，也不要带手杖。} \BibleRef{玛 10:9} 那么怎样？告诉我，\Person{伯多禄}有没有违背这条命令？他确实违背了，因为他有腰带、衣服和鞋子，因为请听天使的话：\Quote{束上腰带，穿上凉鞋。} \BibleRef{宗 12:8} 然而他并不那么需要凉鞋，因为在那时节，人甚至可以赤脚行走；它们最大的用处是在冬天，但他却有了。我们对\Person{保禄}怎么说呢？当他这样给\Person{弟茂德}写信说：\Quote{你要赶紧在冬天以前来}？ \BibleRef{弟后 4:21} 他还给他的命令说：\Quote{我在\Place{特洛阿}留在\Person{卡尔颇}那里的外衣，你带来的时候，也把那些书，尤其是羊皮卷带来。} \BibleRef{弟后 4:13} 看，他说到了外衣，没有人能说他没有任何另一件可以穿的；因为如果他根本不穿一件，命令人带来这一件就是多余的，而如果他不能没有一件可穿的，很明显他有第二件。\par

我们怎么说他\Quote{整整两年住在他自己所租的房子里}？ \BibleRef{宗 28:30} 那么，这个蒙选的器皿竟然违背基督吗？这个人曾说过：\Quote{我活着，但不再是我，而是基督在我内活着}， \BibleRef{迦 2:20} 关于他基督曾作证说：\Quote{他是我所拣选的器皿}？ \BibleRef{宗 9:15} 我应该把这个难题留给你们，不提供任何解答。我应该向你们索取这份惩罚，因为你们疏忽了圣经，因为这就是所有这类困难的根源。因为我们不知道圣经，我们没有受过天主法律的训练，因此我们成了别人过错的挑剔的调查者，却不顾自己的过犯。那么，我应该向你们索取这份惩罚。但我该怎么办呢？父亲们会随意给儿子们许多超出适当的事：当他们父性的怜悯被点燃，看见自己的孩子垂头丧气、因忧伤而憔悴时，他们自己也感受到比孩子更尖锐的痛苦，并且不安息，直到除去他沮丧的原因。\par

但愿至少在此处，你们至少在未领受时感到沮丧，好使你们能善领受。那么这是什么？他们没有反对，绝无此事；而是勤奋地遵守基督的命令，因为那些命令只是暂时的，并不持久；我说这话不是出于猜测，而是出于神圣的圣经。怎么呢？\Person{路加}记载，基督对祂的门徒说：\Quote{当我派遣你们出去，没有钱囊、口袋、腰带和鞋子，你们缺少什么没有？他们说：什么都没有。但今后，你们要预备这些。} \BibleRef{路 22:35} 但请告诉我，他能做什么？他只能有一件外衣吗？那么怎样呢？如果需要洗这件外衣，他因为没有外衣而应该待在家里吗？当需要召唤时，他应该没有外衣而不体面地外出吗？想一想，\Person{保禄}以如此大的成功周游世界，却因为缺乏衣服而待在家里，从而阻碍他高贵的工作，那会是怎样的事。而且如果严寒来临，或者雨水浸湿了它，或者也许冻住了，他怎么能弄干他的衣服？他必须再次没有衣服吗？而且如果寒冷夺去了他身体的力气？他必须因疾病消瘦而无法说话吗？因为请听他给\Person{弟茂德}说的话，以证明他们不是有金刚不坏之身：\Quote{为了你的胃，和你的屡次疾病，稍微用点酒。} \BibleRef{弟前 5:23} 再者，当他谈到另一个人时，他说：\Quote{我认为有必要派你们的使者到你们那里，并供应我的需用。} \BibleRef{斐 2:25} \Quote{因为他确实病了，几乎要死；但天主怜悯了他，不仅怜悯了他，也怜悯了我。} \BibleRef{斐 2:27} 由此可见，他们容易患上各种疾病。那么怎样？他们必须死吗？绝不。那么当时基督为什么给他们那条命令？为了显示祂自己的权能，并证明在以后的时间祂能够这样做，虽然祂没有这样做。但为什么祂没有这样做？他们比以色列人更值得赞美，以色列人的鞋没有穿破，衣服也没有穿破，而且他们是在那酷热阳光的沙漠中旅行，阳光炽热足以熔化石头。\BibleRef{申 29:5} 那么祂为什么这样做？为了你们。因为既然你们不会保持健康，反而会满身伤痕，祂就给了你们那可以作为药的东西。由此显而易见：难道祂自己不能供养他们吗？那位赐给了你们——你们曾是祂的仇敌——的祂，难道不会更多赐给\Person{保禄}吗？那位赐给了以色列人——那些发怨言者、行淫者、拜偶像者——的祂，难道不会更多赐给\Person{伯多禄}吗？他为了祂舍弃了一切。那位允许恶人拥有任何事物的祂，难道不会更多白白赐给\Person{若望}吗？他为了祂甚至离弃了父亲？然而祂不愿意：祂借你们的手喂养他们，使你们得到圣化。请看祂慈爱的丰盛。祂选择让祂的门徒匮乏，好使你们能稍得安慰。\par

因为如果祂使他们免于一切匮乏，他们会更加令人钦佩，更加光荣。但那样，那对你们是救恩的事就会被切断。天主不愿意他们那样令人钦佩，为的是你们得救，而是宁愿他们受贬抑。祂允许他们被轻视，为的是你们能够得救。接受馈赠的老师并不同样受尊敬，而不接受馈赠的才更受尊荣。但在后一种情况中，门徒没有受益，反而被剥夺了果实。你看到如此爱人的天主的智慧了吗？因为正如祂不寻求自己的光荣，也不顾及自己，反而当祂在光荣中时，选择为了你们而受辱，同样，在你们老师的事上也是如此。当他们本可以受到高度尊敬时，祂却宁愿他们为你们的缘故受轻视，为的是你们能够受益，为的是你们能够富足。因为他在今世的事物中匮乏，为的是你们在属灵的事物中富足。那么，如果祂本可以使他们超乎一切匮乏，祂显出为你们的缘故，祂容许他们匮乏。既然知道这些事，让我们转向行善，而非指责。让我们不要过分好奇他人的过失，而要数算自己的；让我们牢记自己的过失，同时数算他人的优点；这样我们就能中悦天主。因为那注视他人过失和自己优点的人，在两方面受到损害：由后者他陷入骄傲，由前者他陷入懈怠。因为当他觉察到某人犯了罪，他很容易自己也会犯罪；当他觉察到自己在某方面优越，他很容易变得骄傲。那忘却自己的优点，只注视自己的过失，而对别人的优点而非罪过好奇探究的人，在许多方面受益。如何呢？当他看到某人做了卓越的事，他受到激励去效法；当他看到自己犯了罪，他变得谦卑和谦虚。如果我们这样做，如果我们这样规范自己，我们将能获得所应许的美善，借着我们的主耶稣基督的恩宠和慈爱，与祂一同，等等。\par

\SourceRecord{Translated by John A. Broadus.；Nicene and Post-Nicene Fathers, First Series；Vol. 13.；Edited by Philip Schaff.；Buffalo, NY: Christian Literature Publishing Co.,；1889.；<http://www.newadvent.org/fathers/230209.htm>.}

% structure: p0001 source=h1 target=section reason=page-title

\section{斐理伯书讲道词 第十篇}

\BibleRef{斐 3:1-3}\par

\Quote{最后，我的弟兄们，你们要因主喜乐。我把同样的事写给你们，对我并不厌烦，对你们却是稳妥的。你们要提防狗，提防作恶的人，提防切损。因为我们才是受割损的人，我们以天主的圣神敬拜，在基督耶稣里夸耀，而不信赖肉体。}\par

\Emph{沮丧}和忧虑，每当它们过度拉紧灵魂时，就夺去了灵魂本有的力量。因此，\Person{保禄}安慰斐理伯人，他们当时非常沮丧；他们沮丧是因为不知道\Person{保禄}的情况如何；他们沮丧是因为以为\Person{保禄}已经完了，因为传道的工作，因为\Person{厄帕夫洛狄托}的缘故。正是在这些事上给予他们保证时，他引入了这些话：\Quote{最后，我的弟兄们，你们要喜乐。} \Quote{你们不再有，}他说，\Quote{沮丧的理由了。你们有了\Person{厄帕夫洛狄托}，你们曾为他而忧伤；你们有了\Person{弟茂德}；我自己也要到你们那里去；福音正在进展。你们此后还缺什么呢？喜乐吧！}\par

现在，他确实称迦拉达人为\Quote{儿女}\BibleRef{迦 4:19}，却称这些人为\Quote{弟兄们}。因为他无论是为了纠正什么，或者为了表示喜爱，就称他们为\Quote{儿女}；但当他以更大的尊荣称呼他们时，就使用\Quote{弟兄们}的头衔。他说：\Quote{最后，我的弟兄们，}他说，\Quote{你们要因主喜乐。}他正确地说\Quote{在主内}，而不是\Quote{照着世俗}。因为这不是喜乐。他说，这些属于基督的苦难带来喜乐。\Quote{我把同样的事写给你们，对我并不厌烦，对你们却是稳妥的。你们要提防狗。}你注意到他是如何避免在开头就引入劝诫的吗？而是在给了他们许多称赞，表达了对他们的钦佩之后，他才这样做，并且又重复了他的称赞。因为这种说法似乎对他们有些严厉。因此他在各方面予以缓和。但他称谁为\Quote{狗}呢？在这个地方有一些人，就是他在所有书信中暗示的那些卑鄙可耻的犹太人，贪图不义之财，喜爱权力，他们想要引诱许多信徒，同时又宣讲基督教和犹太教，败坏了福音。因为他们不容易被识别，所以他说\Quote{要提防狗}：犹太人不再是儿女；从前外邦人被称为狗，但现在犹太人成了狗。为什么呢？因为正如外邦人原是疏远天主和基督的，这些人现在也变成了这样。他显明了他们的无耻和暴力，以及他们与儿女关系的无限距离，因为外邦人曾被称为\Quote{狗}，请听客纳罕妇人所说的话：\BibleQuote{是的，主，可是小狗也吃主人桌子上掉下来的碎屑。}\BibleRef{玛 15:27} 但为了使他们不占这个便宜，因为连狗也在桌子旁边，他补充了那句话，借此也指出他们是外人，说：\Quote{要提防作恶的人}；他妙不可言地表达了自己，\Quote{要提防作恶的人}；他的意思是，他们工作，但目的是恶的，而且工作比懒惰更糟糕，拔除已经有序建立起来的东西。他又说：\Quote{要提防切损；因为我们才是受割损的人}——如何？——\Quote{就是以天主的圣神敬拜，不信赖肉体的那些人。}\par

他说：\Quote{要提防，}他说，\Quote{切损。}割损礼在犹太人眼中是庄严的，因为连法律本身也为之让步，安息日也不如割损受尊重。因为为了行割损，可以违犯安息日；但为了守安息日，割损从未被违犯；并且，我请你们注意天主的安排。这发现甚至比安息日更隆重，因为在某些时候不被省略。那么，当它被废除时，安息日更被废除了。因此\Person{保禄}将名称割断，说：\Quote{要提防切损}；并且他没有说\Quote{割损是邪恶的，是多余的}，免得使那些人惊慌，而是更明智地处理，使他们脱离那事，却用词语讨好他们，不，更是用那事本身以更严肃的方式。但在迦拉达人的情况则不然，因为那里病势严重，他立即采取正面而大胆的切割疗法；但在这种情况下，因为他们没有做出那样的事，他赐予他们头衔的满足，他驱逐了其他人，并说：\Quote{要提防切损；因为我们才是受割损的人}——如何？——\Quote{就是以天主的圣神敬拜，不信赖肉体的那些人。}他没有说\Quote{我们检验这两种割损，哪一个更好}；他甚至不允许它在名称中有一份；但他说什么呢？那割损是\Quote{切损}。为什么？因为他们除了切割肉体之外什么也不做。因为所做的若不依据法律，就不过是切割和割裂肉体；那么他之所以这样称呼，或者因为这个原因，或者因为他们试图将教会一分为二；我们称那些随意、无目的、无技巧这样做的人为\Quote{切割}。现在，他说，如果你们必须寻求割损，你们会在我们中间找到，就是\Quote{以天主的圣神敬拜}的人，即属灵敬拜的人。\par

请你回答我：灵魂与身体，哪一样更高贵？显然是前者。因此，那割损也是更高贵的，或者不如说，它不再更高贵，因为这才是唯一的割损；当预象尚存时，祂恰当地将它一并提出，写道：\BibleQuote{你们要割损你们的心包皮。}\BibleRef{耶 4:4}同样，在\BibleBook{罗马}{书}中，他废除了它，说：\BibleQuote{其实，外表上做犹太人的，并不是真犹太人；外表上肉身的割损，也不是真割损；惟独内心做犹太人的，才是真犹太人，而割损是内心的，在于神，不在于文字。}\BibleRef{罗 2:28}最后，他甚至剥夺了它的名称，主张\Quote{它也不是割损}；因为当现实尚未来临时，预象被称为这名，但当现实来到时，它就不再保留这称号。如同在绘画中，有人勾勒了一位国王的轮廓；只要色彩尚未涂上，我们就说：\Quote{看哪，这就是国王}；但当色彩加上后，预象就被现实所吞没，不再显现。他并没有说：\Quote{因为割损在我们内}，而是说：\Quote{我们就是割损}，这是公正的；因为这才是人，有德行的割损，这才是真正的人。他没有说：\Quote{因为在他们那里有\InnerQuote{割损}}，因为他们从此处于毁灭和邪恶的状态。但他却说：不再是在身体上施行割损，而是在心中。\BibleQuote{并且不要倚靠肉体；虽然我若愿意，也可以倚靠肉体。}\BibleRef{斐 3:4}这里他所说的\Quote{倚靠}和\Quote{肉体}是指什么？夸耀、大胆、高傲。他加上这一点做得好；因为如果他来自外邦人，并且谴责割损，不仅割损，而且所有不合时宜接受割损的人，那么似乎他是在贬低它，因为他缺乏犹太教崇高血统的资格，仿佛是外人，对其庄严礼仪一无所知，也没有任何分内。但事实上，他虽是一位分享者，却仍指责他们，因此他不是因为没有分内而指责他们，而是因为否认它们；不是出于无知，而恰恰是因为对它们非常熟悉。因此，请观察他在\BibleBook{迦拉达}{书}中所说的；当他被迫不得不谈论自己的伟大之处时，即便在这种情况下，他所表现的除了谦卑别无他物。他说：\BibleQuote{你们听说过我从前在犹太教中的生活}\BibleRef{迦 1:13}；同样在此处：\BibleQuote{若别人以为自己可以倚靠肉体，我更可以。}\BibleRef{斐 3:4}他立即加上：\Quote{是希伯来人的希伯来人}。但他说：\Quote{若别人}，显示这是出于必要，显示他是为了他们才说的。他说：\Quote{如果你们能倚靠}，我也这样说，因为我是沉默的。注意他责备时毫无刻薄；通过不指名道姓，他甚至给了他们悔改的机会。\Quote{若有人以为自己可以倚靠}；他巧妙地用了\Quote{以为}，或者因为他们实际上并没有这样的倚靠，或者因为那倚靠并非真正的倚靠，因为一切都是出于必要，而非出于选择。\Quote{第八天受了割损}；他列出了第一件事，即他们主要夸耀的割损之礼。\Quote{属于以色列民族}。他指出了这两个事实：他既不是归化者，也不是归化者所生；因为从他第八天受割损可知他不是归化者，从他属于以色列民族可知他不是归化者的后代。但为了不让你想象他属于以色列民族是来自十个支派，他说：\Quote{是本雅明支派}。这样他属于更被认可的支派，因为司祭的职位就在这个支派的地业中。\Quote{是希伯来人的希伯来人}。因为他不是归化者，而是从古以来就是显赫的犹太人；因为他可能属于以色列，却不一定是\Quote{希伯来人的希伯来人}，因为许多人已经败坏了自己，语言不通，被其他民族包围；因此他要表明的或是这一点，或是他出身的高超。\Quote{就法律说，是法利塞人}。他现在转向取决于他自己意志的事；因为所有那些事都不在于他自己的意志：他受割损并非出于自己，他属于以色列民族，他属于本雅明支派。因此，即使在这些方面他也有着更大的份额，尽管实际上有许多人同他一样。那么\Quote{更}体现在哪里呢？特别在于他不是归化者；因为属于最显赫的支派和派别，并且从以前的祖先起就是，这是许多人不能享有的。但他来到了那些选择的事，其中我们有\Quote{更}。\Quote{就热诚说，迫害教会}。但这还不够；因为一个人甚至可以是法利塞人，却并不十分热衷。但他又加上了这一点；看吧，那\Quote{更}。\Quote{就正义说}。然而，一个人也可能冒险从事，或出于野心而非对于法律的热诚而行动，如同大司祭们所做。但也不是这种情况，而是：\Quote{就法律上的义说，是无瑕可指的。}因此，他问道：如果我在血统纯正、热心、习惯和生活方式上都超越众人，为什么我要放弃所有这些尊位？岂非因为我发现基督的事更好，而且好得多？故此他加上：\BibleQuote{然而凡以前对我有利益的事，我如今为了基督，都看作是损失。}\BibleRef{斐 3:7}\par

如此严格自律、自幼便开始的生涯，如此无可指摘的出身，如此的危险、阴谋、劳苦、热忱——保禄竟一概舍弃，\Quote{视之为损失}，而那些之前本是\Quote{利益}的，他为了\Quote{赢得基督}而丢弃。但我们连金钱也不肯轻视，好使我们\Quote{赢得基督}，却宁愿放弃来世的生命，也不愿失去现世的美好事物。然而这无非就是损失；因为请告诉我，让我们仔细检视财富的种种状况，看看它是否不是伴随烦恼的损失，且毫无益处。请说，那些成堆的昂贵衣物有什么好处？我们穿戴它们时得到了什么善？没有，反而有所损失。如何损失？因为连穷人身着廉价破旧的衣服，在炎热时受的烘烤也丝毫不比您差；甚至他忍受得更好，因为破旧单薄的衣服使身体更舒适，而新做的衣服，即使比蛛网更细，也不如此。此外，您因过度自傲，甚至穿两件、常常三件内衣，外加外套、腰带和裤子；但没人责备穷人若只穿一件内衣。因此穷人是最轻松忍受的人。正因如此，我们看见富人出汗，而穷人毫无此等困扰。既然穷人那件售价低廉的衣服起到同样甚至更好的作用，而那些衣服要人支付许多黄金，却只做同样的事，这种极大的多余岂不是损失？因为它在使用和服务上没有增添任何东西，而您的钱袋却因此少了那么多金币；用途和服务却相同。您这个富人花了一百金币甚至更多购买之物，穷人却用一点小钱就买到了。您看到损失了吗？不，因为您的骄傲不让您看见。您要我们在金饰方面也算这笔账吗？人们把金饰同样用在马匹和妻子身上。因为除了其他邪恶，金钱使人变得愚昧；他们视妻子与马匹值得同样的尊荣，二者的装饰相同；他们竟用与驮着他们的牲畜或他们乘坐的帐篷皮料相同的手段来打扮自己。用黄金装饰骡子或马匹有何用处？那位女士全身佩戴如此沉重的金子和珠宝，她又得到了什么？\Quote{但金饰不会磨损，}他答道。诚然，在澡堂和许多地方，宝石和金饰都会大大贬值。但就算如此，且承认它们不损坏，请告诉我，益处是什么？又当它们脱落丢失时，岂非遭受损失？而当它们招致嫉妒和阴谋时，岂非损失？因为当它们对佩戴者毫无益处，反而点燃嫉妒者的眼睛，刺激强盗时，岂非成了损失？再说，当一个人本可利用它们做有益的事，却因妻子的奢侈而无法做到，不得不节衣缩食，好让她浑身金饰，这难道不是损失？因为财物之所以得名\OriginalTerm{财物}{goods}，正是为了使用，而非要我们像金匠的样品那样摆弄它们，而是要用它们行善；那么，当爱财不允许这样做时，整个事情岂非损失？因为不敢使用的人，就像对待别人的财产一样抑制着不用，完全没有任何用途。\par

再者，当我们建造宏伟宽敞的宅邸，装饰以圆柱、大理石、柱廊、拱廊，用尽各种方式，到处设置雕像和塑像时，又如何？许多人甚至从这些——即雕像——中召叫邪魔，但我们暂且不查究这些。镀金的天花板有何意义？难道它给居处中等的人提供的需要不一样吗？\Quote{但其中有极大的愉悦，}他说。是的，第一天或第二天，之后完全没有了，只是空置而已。因为太阳若因习以为常而不令我们惊奇，那么艺术品更是如此；我们看它们如同看泥土之物。请告诉我，一排柱子、最美的雕像、或墙上的镀金，有何作用使您的居所优于他人？毫无作用；相反，这些来自奢侈、狂妄、过度的骄傲和愚蠢；因为那里的一切都应是必需且有用的，而非多余之物。您看见这件事是损失了吗？您看见它是多余且无益的吗？因为如果它不提供更多用途或愉悦，而经过一段时间，它\Quote{确实}带来厌倦，那它无非是损失；而虚荣是障碍，不让我们看见这一点。\par

保禄难道舍弃了他视为\Quote{利益}的东西，而我们竟不肯为基督舍弃我们的损失吗？我们还要被钉在地上多久？我们何时才能仰望天堂？你们没有注意到老年人，他们对过去的事记忆多么模糊吗？你们没有注意到那些行完人生路程的人，无论老年还是青年？你们没有看到那些在生命中途失去一切的人吗？我们为什么如此执着于不稳定的事物？为什么我们与变幻无常的东西捆在一起？我们还要多久才能抓住那永恒的事物？老年人难道不愿付出一切，只要能摆脱衰老吗？那么，渴望回到从前的青春，并乐意为此付出一切，以便变得年轻，这是多么不合情理！然而，当我们能获得一个永不衰老的青春，一个与巨大财富结合、更有精神的青春时，我们却不愿放弃一点琐事，紧抓着那些对现世生命毫无益处的东西。它们永远不能使你免于死亡，没有能力驱除疾病，阻止衰老，或任何那些必然发生、合乎自然规律的事件。你仍然紧抓着它们吗？告诉我，你得到了什么？醉酒、贪食、违反本性的快乐和各种各样的享乐，这些比最严厉的主人还要折磨人。\par

这就是我们从财富中获得的好处，此外没有别的，因为我们没有这样的心志；如果我们有心志，我们本可以通过财富赢得天堂本身作为产业。\Quote{那么，财富是好的，}他说。产生这效果的并不是财富，而是拥有者的意志；因为是意志成就这事，所以即使穷人也有能力赢得天堂。因为，如我常说的，天主不看礼物的数量，而是看给予者的心意；甚至一个穷人，只要给予少许，也能赢得一切，因为天主要求与我们能力相称的奉献。财富不能确保我们进天堂，贫穷也不能确保我们下地狱；但善意或恶意则决定二者之一。那么，让我们纠正这一点，让我们重新拥有这一点，让我们管束这一点，那么一切对我们来说都是容易的。\par

因为正如工匠用斧头加工木材，无论斧头是铁的还是金的，他都能同样工作，甚至用铁斧工作得更好；同样，在贫穷的状态下更容易保持正直的德行之路。关于财富，我们读到：\BibleQuote{骆驼穿过针眼，比富人进入天国更容易。}\BibleRef{玛 19:24} 但祂对于贫穷并没有作这样的声明；不，恰恰相反。\BibleQuote{变卖你的财产，施舍给穷人，然后来跟随我。}\BibleRef{玛 19:21} 好像跟随的行为是从变卖中产生出来的。\par

那么，我们绝不可逃避贫穷，视之为邪恶，因为它是天堂的导引者。同样，我们绝不可追求财富，视之为好事；因为它们是那些行走不慎之人的覆灭，但我们在一切事上要定睛仰望天主，按需要运用祂所恩赐的种种恩赐，包括身体的力气、大量的金钱和每一个其他恩赐；因为我们是为祂而存在的，若使这些事物有益于他人，却无益于创造我们的祂，那是不合情理的。祂造了你的眼睛：使它有益于祂，而非有益于魔鬼。但如何有益于祂？通过默观祂的受造物，赞美、光荣祂，并远离一切对女人的注视。祂造了你的手吗？为祂的用途保存它们，不为魔鬼，不伸手去抢劫、掠夺，而是为祂的诫命和善行，为恳切的祈祷，为伸手援助跌倒的人。祂造了你的耳朵吗？将这些献给祂，不要用于柔靡的曲调或可耻的故事；而要\BibleQuote{你们的一切言谈，都在至高者的法律中。}\BibleRef{德 9:15} 因为\BibleQuote{站，}他说，\BibleQuote{在长老的群体中，凡是智慧的人，都要紧靠他。}\BibleRef{德 6:34} 祂造了你的口吗？让它不做任何令祂不悦的事，而要咏唱圣咏、赞美诗与属神的歌曲。\BibleQuote{你们口中不可出任何败坏的话，}宗徒说，\BibleQuote{只要按需要说造就人的好话，使听的人获得恩宠。}\BibleRef{弗 4:29} 这是为了建立，而非拆毁；为了美言，而非诽谤和谋害他人，而是完全相反。祂造了你的脚，不是要你跑去作恶，而是行善。祂造了你的肚腹，不是要你填到爆裂，而是要实践智慧的教训。为了生儿育女，祂在你心中植入欲望，不是为了淫乱，也不是为了奸淫。祂赐给你理智，不是使你成为亵渎者或辱骂者，而是使你无伪。祂赐给我们金钱，用于适当场合，力气也用于适当场合。祂设立了技艺，使我们现世的存在得以维系，不是要我们与属神的事分离，不是要我们献身于卑贱的技艺，而是从事必要的技艺，使我们彼此服务，互相助益，而不是互相谋害。祂赐给我们屋顶，为了遮雨，仅此而已，不是要用黄金装饰，而穷人却饿死。祂赐给我们衣服遮盖身体，不是为了炫耀，不是要在这些物品上浪费大量黄金，而让基督赤身裸体地死去。祂赐给你一个避难所，不是要你独占，而是要也提供给别人。祂赐给你土地，不是要你切取其最大部分，将天主的美好恩赐花费在娼妓、舞女、戏子、笛手和竖琴手身上，而是花费在那些饥饿和贫困的人身上。祂赐给你海洋以供航行，使你不因旅途劳累，不是要你探测它的深处，从中打捞宝石和诸如此类的东西，也不是要你以此营生。\par

\Quote{那么为什么有宝石呢？}他说。不，请你告诉我，为什么这些石头是这样，为什么一类被认为很有价值，而另一类更有用？因为这些可用于建造，那些则毫无用途；而且这些比那些更坚硬。\Quote{但是它们，}他说\Quote{产生美的效果。}怎么会呢？这是出于幻想。它们更白吗？不，它们不比纯白大理石更白，甚至不及。它们更坚硬吗？这一点也无法说。那么，它们更有用吗？它们更大吗？这一点也无法说。那么，它们为何如此被赞赏，如果不是出于幻想？因为如果它们既不更美（因为我们会发现其他更闪亮、更白的石头），也不更有用，也不更坚硬，那么它们如何被如此赞赏？难道不是仅仅出于幻想吗？那么天主为什么赐予它们？它们不是祂的恩赐，而是你们自己的想象，认为它们是什么了不起的东西。\Quote{那么，怎么，}他回答\Quote{连圣经也赞赏它们呢？}到此为止，圣经是针对你们的幻想说的。就像一个老师在跟孩子说话时，常常称赞孩子所称赞的东西，当他想要吸引和引导孩子时。\par

为什么你们追求衣服的华丽？祂给你们穿上了衣服和凉鞋。但这些有什么理由呢？\Quote{天主的判断，}他说\Quote{比黄金更可爱；是的，比极多的纯金更可爱。}\BibleRef{咏 18:10}这些，亲爱的，毫无用处。如果它们有用，祂就不会命令我们鄙视它们。至于圣经，它是依照我们的观念说话的，这也是天主慈爱的表现。\Quote{那么，}他问\Quote{为什么祂赐予紫色和类似的东西？}这些是天主恩赐的产物。因为祂也愿意藉着其他事物显示祂自己的丰富。祂也赐予你们谷物本身；但你们用谷物制作许多东西，糕点和蜜饯，各种各类，享受很多快乐。享乐和虚荣产生了所有这些发明。你们喜欢把它们放在一切之前。因为如果一个外国人或乡下人，不熟悉这片土地，提出疑问，看到你们的赞赏，问说：\Quote{为什么你们赞赏这些？}你们有什么可说的？说它们好看吗？但并非如此。那么，让我们放弃这些想法；让我们抓住真正真实的事物。这些不是真实的，只是消逝，像河流一样流过。因此，我劝你们，让我们站立在磐石上，好使我们既不容易被动摇，又能获得将来的美善，藉着我们的主耶稣基督的恩宠和慈爱，与祂一同等。\par

\SourceRecord{Translated by John A. Broadus.；Nicene and Post-Nicene Fathers, First Series；Vol. 13.；Edited by Philip Schaff.；Buffalo, NY: Christian Literature Publishing Co.,；1889.；<http://www.newadvent.org/fathers/230210.htm>.}

% structure: p0001 source=h1 target=section reason=page-title

\section{论\WorkTitle{斐理伯书}第十一讲}

\BibleRef{斐 3:7-10}\par

\BibleQuote{但先前对我有益的，我为了基督，都看作损失。不但如此，我也将一切看作损失，因为我以认识我主基督耶稣为至宝；为了祂，我自愿损失一切，拿一切当废物，为赚得基督，并得以在祂内，不是有我那因法律而来的正义，而是有因信仰基督而来的正义，即出于天主、基于信德的正义。}\par

在我们与异端者的争辩中，我们必须以充沛的精力发起攻击，使他们能够仔细留意。因此，我要从上次结束的地方开始我现在的讲道。上次结束于何处？在枚举了犹太人的一切夸口——无论是出于出身还是出于选择——之后，他接着说：\BibleQuote{但先前对我有益的，我为了基督，都看作损失。因为我以认识我主基督耶稣为至宝；为了祂，我自愿损失一切，拿一切当废物，为赚得基督。}\par

如果这话说得直白，他们就会像在其他地方那样做：他们会删掉这些话，他们会否认圣经，因为他们根本不敢面对它。但就像捕鱼时，诱饵被隐藏起来让鱼游近，不显露出来；这里也发生了同样的事。他们说，\Person{保禄}称法律为\BibleQuote{废物}，称为\BibleQuote{损失}。他说，除非我\BibleQuote{蒙受}这\BibleQuote{损失}，否则就不可能有得着基督。所有这些都使异端者接受这段经文，以为是对他们有利；但一旦他们接受了，他就从四面用网将他们围住。因为他们自己说什么？看！法律是\BibleQuote{损失}，是\BibleQuote{废物}；那么你们怎么说它是出于天主的呢？\par

但这些话恰恰对法律有利，它们如何有利，下面便显明了。让我们仔细留意他确切的话。他没有说，法律是损失；而是说：\BibleQuote{我视为损失。} 但谈到利益时，他没有说，我视为，而是说\BibleQuote{它们曾是有益的。} 但谈到损失时，他说：\BibleQuote{我视为}；这是正确的；因为前者本是如此，而后者则因我的看法而成为如此。\Quote{那么，难道不是这样吗？}他说。是为了基督而成为损失。\Quote{那么，难道不是这样吗？}他说。是为了基督而成为损失。\par

法律如何成为收益？它并不是\Emph{被视为}收益，而是\Emph{本来就是}。试想这是多么伟大的事：将本性残暴的人转变为人的模样。如果法律不存在，恩宠就不会赐下。为什么？因为它成了一种桥梁；既然从深重的卑贱中无法上升到高处，便造了一个阶梯。但登上者不再需要阶梯，却并不轻视它，反而感激它。因为它使他处于不再需要它的境地。然而正因他不再需要它，他理当承认亏欠，因为他本不能飞翔。法律也是如此，它带领我们升到高处；因此它曾是收益，但今后我们视它为损失。如何？不是因为它是损失，而是因为恩宠大得多。就像一个穷人，饥饿时，只要有银子就能免于饥饿；但当他找到金子，又不允许同时保留两者时，他认为保留前者是损失，便扔掉它而取金币；这里也是如此；不是因为银子是损失（它不是），而是因为不能同时拿两者，必须舍弃一个。因此，不是法律本身是损失，而是人固守法律而离弃基督。所以，当它领我们离开基督时，它才是损失。但如果它把我们引向祂，就不再是这样了。为此他说\BibleQuote{损失}（原文是\BibleQuote{为了基督而损失}？）如果是为了基督，它就不是本性的损失。但法律为什么不允许我们来到基督面前？正是为此，它才被赐下。而基督是法律的满全，基督是法律的终向。\BibleQuote{因为基督是法律的终向。}遵守法律的人，便离开法律本身。如果我们留意法律，它是允许的；但如果我们不留意，它就不允许。\BibleQuote{不但如此，我也将一切看作损失。} 他意思是，我为什么说法律呢？世界不是好的吗？现世生命不是好的吗？但如果它们使我离开基督，我就视这些为损失。为什么？\BibleQuote{因为我以认识我主基督耶稣为至宝。} 因为太阳出现后，点着蜡烛就是损失；所以损失来自比较，来自另一者的优越。你看，保禄从优越性而非种类相异进行对比；因为优越的东西是优于与自己同类的某物。所以，他用同一方式表明那认识的关联，即通过对比得出优越性。\BibleQuote{为了祂，我自愿损失一切，拿一切当废物，为赚得基督。} 还不清楚他是否指法律，因为他很可能是指现世的事物。因为当他说，\BibleQuote{先前对我有益的，我为了基督，都看作损失；不但如此，}他接着说，\BibleQuote{我还将一切看作损失。} 虽然他说一切，但指的是现今的事物；即使你愿意也包括法律，也并非侮辱。因为废物来自小麦，小麦的力量就是废物，我指的是糠秕。正如废物在先前状态中有用，我们把它与小麦一起收集；没有废物就没有小麦，法律也是如此。\par

你看，他处处称它为\BibleQuote{损失}，并非本身是损失，而是为了基督。\BibleQuote{不但如此，我还将一切看作损失。} 为什么又如此？\BibleQuote{因为我以认识（祂）为至宝，为了祂，我自愿损失一切。} 又，\BibleQuote{所以我也将一切看作损失，为赚得基督。}\par

你看，他从各方面抓住基督作为他的根基，不容法律在任何地方暴露或受打击，而是从各方面保护它。\BibleQuote{并且为要我在他内被找到，不拥有我自己的义德，就是那出于法律的义德。}若那拥有义德的人，因自己的义德算不得什么，而奔向这另一种义德，那么那些没有义德的人，岂不更该奔向他吗？他说的很好：\BibleQuote{我自己的义德}，不是那由劳苦和辛劳获得的，而是那从恩宠中得到的。如此，若那如此卓越的人都是因恩宠得救，何况你们呢？因为很可能他们会说那从劳苦而来的义德更大，他就证明它在另一种义德面前如同粪土。否则，我在这义德中如此卓越，就不会抛弃它，而奔向另一种了。但那另一种是什么？是那出于对天主的信德的义德，也就是说，它也是由天主赐予的。这就是天主的义德；这完全是恩赐。而天主的恩赐远超过那些凭我们自己的勤勉而得的无价值的善工。\par

但\BibleQuote{借着信德，为要认识他}是什么意思？那么，知识是借着信德，没有信德就不可能认识他。怎么会呢？借着信德，我们必须\BibleQuote{认识他复活的德能。}有什么理性可以向我们证明复活呢？没有，只有信德。因为如果按肉体来说的基督的复活是凭信德认识的，那么天主圣言的受生怎能凭推理理解呢？因为复活小于受生。为什么？因为复活已有很多例子，但受生却从未有过；因为在基督之前有许多死人复活，虽然他们复活后又死了，但从未有人由童贞女所生。那么，如果我们必须借信德理解那低于按肉体的受生的事，那么那远远更大、无可比拟地更大的事，怎能凭理性理解呢？这些事造就义德；我们必须相信他能做到，但他如何能做到，我们无法证明。因为共受他的苦难是出于信德。但如何呢？如果我们不信，我们也不会受苦；如果我们不信\BibleQuote{我们若与他一同忍耐，也必与他一同为王}\BibleRef{弟后 2:12}，我们就不会忍受苦难。受生与复活都是由信德理解的。你看见吗，信德不能只是抽象地，而必须是借着善工；因为谁特别相信基督复活了，谁就同样将自己交付危险，与他的苦难有分。因为他与那复活的那位、与那生活的那位有分；因此他说：\BibleQuote{为要我在他内被找到，不拥有我自己的义德，就是那出于法律的义德，而是那借着对基督的信德而来的义德，也就是那由天主借着信德而来的义德：为要认识他和他复活的德能，并共受他的苦难，肖似于他的死亡；或许我得以达到从死者中的复活。}他说，肖似于他的死亡，就是有分；而由于他受人苦难，我也同样；因此他说\BibleQuote{变成肖似}，又在另一处说：\BibleQuote{在我肉身上补满基督苦难的欠缺}\BibleRef{哥 1:24}。这就是说，这些迫害和苦难塑造了他死亡的形像，因为他寻求的不是自己的利益，而是许多人的益处。\par

因此，迫害、苦难和困境不应使你困扰，反而应使你喜乐，因为我们藉着它们\BibleQuote{相似祂的死亡。}正如他所说的，我们被塑造为祂的样式；正如他在别处写道：\BibleQuote{身上常带着主耶稣的死。}\BibleRef{格后 4:10}这也出于伟大的信德。因为我们不仅相信祂复活了，也相信祂复活后拥有大能：因此我们走祂所走的路，即我们在这一点上也成为祂的弟兄。正如他所说，我们在这一点上成为\Person{基督}的。哦，苦难的尊荣何等伟大！我们相信，我们藉着苦难成为\BibleQuote{相似祂的死亡}！因为如在圣洗中，我们\BibleQuote{与祂的死亡相似而埋葬}，在这里，则是与祂的死亡本身。他在那里恰当地说\BibleQuote{相似祂的死亡}\BibleRef{罗 6:4}，因为在那里我们并非完全死亡，我们不是死于肉身、身体，而是死于罪。既然提到了两种死亡——祂确实死于肉身，而我们死于罪，并且在祂那里，是祂所取的那人（那在我们肉身内的）死了，而在这里，是罪人死了——因此他说\BibleQuote{相似祂的死亡}；但这里，不再是相似祂的死亡，而是祂的死亡本身。因为\Person{保禄}在他的迫害中，不再是死于罪，而是死于自己的身体。因此，他忍受了同样的死亡。\BibleQuote{或者}，他说，\BibleQuote{我也得以从死人中复活。}你说什么？所有人都将分享那复活。\BibleQuote{我们不是都要睡觉，而是都要改变}\BibleRef{格前 15:51}，并且所有人都将分享的不仅是复活，还有不朽坏。有些是为荣耀，但有些是为惩罚。如果所有人都分享复活，且不仅复活，还有不朽坏，他为何说\BibleQuote{或者我也得以}，好像要分享某种特殊的东西？\BibleQuote{因此}，他说，\BibleQuote{我忍受这些事，或者我也得以从死人中复活。}因为如果你没有死，你就不会复活。那么这是什么意思？似乎暗示了一件伟大的事。如此伟大，以至于他不敢公开断言，只说\BibleQuote{或者}。我相信祂和祂的复活，而且我还为祂受苦，但我仍不能对复活有信心。他在这里提到的是什么复活？那就是通往\Person{基督}本身的复活。我说，我相信\BibleQuote{祂和祂复活的大能}，并且我\BibleQuote{有分于祂的苦难}，并且我\BibleQuote{成为相似祂的死亡。}然而在这些之后，我丝毫没有信心；正如他在别处说：\BibleQuote{那自以为站得稳的，要小心，免得跌倒。}\BibleRef{格前 10:12}又说：\BibleQuote{我恐怕我传福音给别人，自己反被弃绝。}\BibleRef{格前 9:27}\par

12节。\BibleQuote{这并不是说我已经得着了，或已经成全了；而是说我竭力追求，或者可以得着基督耶稣之所以得着我的。}\BibleRef{斐 3:12}\par

\BibleQuote{这并不是说我已经得着了。}\BibleQuote{已经得着}是什么意思？他指的是奖赏，但如果那忍受了如此苦难、被逼迫、\BibleQuote{身上带着主耶稣的死，}的人，尚且对那复活没有信心，我们能说什么？\BibleQuote{或者可以得着}是什么意思？他先前所说的：\BibleQuote{或者也可以达到从死人中复活。}\BibleRef{格后 4:10}如果我可以得着，他说，祂的复活；即如果我能忍受如此大事，如果我能效法祂，如果我能成为相似祂的。例如，\Person{基督}遭受了许多事，祂被吐唾沫，被击打，被鞭打，最后祂所受的一切。这是整个历程。通过这一切，人必须忍受整个赛程，然后到达祂的复活。或者他指的是：如果我认为配得那荣耀的复活（这是有信心的事），以到达祂的复活。因为如果我能忍受所有的竞赛，我也将能拥有祂的复活，并荣耀地复活。他说，我尚未配得，但\BibleQuote{我竭力追求，或者可以得着。}我的生活仍是一场竞赛，我离终点仍远，我离奖赏仍远，我仍在奔跑，我仍在追赶。他没有说\Quote{我奔跑}，而是说\BibleQuote{我追赶。}因为你知道一个人追赶时的急切。他看不见任何人，他以大力推开所有打断他追赶的人。他聚集他的心思、视线、力量、灵魂和身体，除了奖赏之外不看别的。但如果\Person{保禄}如此追赶，如此受苦，尚且说\BibleQuote{或者可以得着}，我们这些放松努力的人又能说什么？然后为显示这事是出于亏欠，他说：\BibleQuote{基督耶稣所以得着我的，正是为此。}他说，我曾是丧亡者中的一员，我奄奄一息，我几乎死了，天主抓住了我。因为祂追赶我们，当我们从祂面前逃跑时，祂全力以赴。所以他指出了这一切；因为\BibleQuote{我被得着，}这句话，显示了那愿意得着我们的祂的恳切，以及我们极度的厌恶、我们的流浪、我们逃离祂。\par

因此，我们欠下了巨额的债务，却没有人在意，没有人哭泣，没有人叹息，所有人都回到了从前的状态。因为在基督出现之前，我们逃离天主，如今也是如此。我们能够逃离天主，不是地方上的逃离，因为祂无所不在；听先知所说：\Quote{我往何处，才能脱离祢的神能？我去哪里，才能逃避祢的圣容？}\BibleRef{咏 138:7}那么，我们如何能逃离天主呢？如同我们可以远离天主，可以远远地离开。\Quote{远离祢的人，必然沉沦。}它说。\BibleRef{咏 72:27}又说：\Quote{不是你们的罪恶在我和你们之间造成了隔阂吗？}\BibleRef{依 59:2}那么，这远离是如何发生的？这隔阂是如何发生的？在于心意和灵魂，因为不可能是地方上的。一个人如何能逃离那无所不在者呢？罪人于是逃离。这就是圣经所说的：\Quote{恶人无人追赶，也要逃跑。}\BibleRef{箴 28:1}我们急切地逃离天主，虽然祂总是追赶我们。宗徒急忙要靠近祂。我们却急忙要远离祂。\par

这些事难道不值得哀叹吗？难道不值得流泪吗？可怜可悲的人，你逃往何处？你逃离你的生命和你的救恩，逃往何处？如果你逃离天主，你投靠谁呢？如果你逃离光明，你向何处举目呢？如果你逃离你的生命，今后你靠什么生存呢？让我们逃离我们救恩的仇敌吧！每当我们犯罪，我们就逃离天主，我们如同逃亡者，我们前往异地，如同那消耗了父亲产业、前往异地、耗尽父亲一切所有、穷困潦倒的人。我们也有来自父亲的产业；这产业是什么？祂免除了我们的罪；祂白白赐给我们能力、行善的力量；祂白白赐给我们意愿、忍耐；祂在圣洗中白白赐给我们圣神；如果我们耗尽了这些，我们今后将陷入窘困。因为如同病人，只要他们仍受热病和体液失调的困扰，就不能起身、工作或做任何事；但如果有人使他们解脱，恢复健康，那时他们若不工作，这便源于他们自己的懒惰。我们也是如此。因为疾病沉重，热病过度。我们不是躺在病床上，而是躺在邪恶本身中，被罪恶丢弃，如同在粪堆上，满身疮痍，臭气熏天，污秽不堪，日渐消瘦，更像鬼魂而非活人。恶灵环绕我们，\OriginalTerm{今世之王}{Prince of this world}嘲笑并攻击我们；天主的独生子来了，发出祂显现的光线，立即驱散了黑暗。那坐在父宝座上的君王来到我们这里，离开了祂的父的宝座。当我说\Quote{离开}时，不要以为祂移动了，因为祂充满天地，我说的是\OriginalTerm{救恩经济}{economy}；祂来到一个仇敌面前，一个憎恨祂、转身逃跑、不能忍受观看祂、每日亵渎祂的人。祂看见他躺在粪堆上，被虫咬噬，遭受热病和饥饿之苦，患有各种疾病：热病折磨他，这是邪欲；发炎沉重地压着他，这是骄傲；饥饿啃咬着，这是贪婪；四围的疮痍溃烂，这是淫乱；眼目失明，这是崇拜偶像；口哑癫狂，这是叩拜木石并向它们说话；极度的丑陋，因为邪恶正是如此，面目可憎，并是一种最沉重的疾病。祂看见我们说话比疯人更愚妄，称木头为我们的天主，石头也一样；祂看见我们如此严重的罪过，并没有拒绝我们；没有发怒，没有转身离去，没有憎恨我们，因为祂是主，不能恨自己的受造物。但祂做了什么？如同一位最卓越的医生，祂准备了极贵重的药物，并亲自首先品尝了它们。因为祂首先亲自追随了德性，从而将它赐给了我们。祂先赐给我们洗涤，如同某种解毒剂，于是我们吐出了所有的罪过，一切同时逃逸，我们的炎症停止，热病熄灭，疮痍干愈。因为所有来自贪婪、愤怒及其他一切罪恶都被圣神消散了。我们的眼睛开了，耳朵开了，舌头说出了圣洁的话；我们的灵魂获得了力量，身体获得了如此美丽和容光，如同天主子从圣神的恩宠中应有的那样；如此荣耀，如同新生的王子，在紫色中养育所应有的那样。唉！祂赐予我们何等崇高的尊贵啊！\par

我们出生了，被养育了，为何又逃离我们的恩主呢？祂成就了这一切，也赐给我们力量，因为被疾病压弯的灵魂不可能承受，难道不是祂亲自赐予我们力量吗？祂赐予我们罪的赦免。我们消耗了一切。祂赐给我们力量，我们浪费了它。祂赐给我们恩宠，我们熄灭了它；如何？我们把它消耗在不恰当的事上，没有用于任何有益的目的。这些事毁灭了我们，而比一切更可怕的是，当我们在异乡，吃着豆荚时，我们不说：\Quote{我要回到父亲那里去，说：\InnerQuote{我得罪了天，也得罪了祢。}}\BibleRef{路 15:18}而且，我们有一位如此慈爱的父亲，祂热切渴望我们回头。只要我们回到祂那里，祂甚至不愿追究我们先前的事迹，只要我们离开它们。对祂来说，我们回头便是足够的道歉。祂不但自己不追究，如果别人追究，祂也堵住他的口，即使指控者是有声望的人。让我们回头吧！我们还要站立多远？让我们觉察我们的耻辱，让我们意识到我们的卑劣。罪恶使我们成为猪，罪恶给灵魂带来饥荒；让我们重新找回自己，再次清醒，回到我们从前的高贵出身，以获得将来的美好事物，在我们的主耶稣基督内，愿光荣、权能、尊崇归于祂及圣父及圣神，从今世直到永远，万世无疆。\par

\SourceRecord{Translated by John A. Broadus.；Nicene and Post-Nicene Fathers, First Series；Vol. 13.；Edited by Philip Schaff.；Buffalo, NY: Christian Literature Publishing Co.,；1889.；<http://www.newadvent.org/fathers/230211.htm>.}

% structure: p0001 source=h1 target=section reason=page-title

\section{斐理伯书第十二篇讲道}

\BibleRef{斐 3:13,14}\par

\BibleQuote{各位弟兄，我并不以为我已经夺得：但只有一件事，忘记背后，努力向前，向着标竿直跑，要得天主在基督耶稣里从高天召我来得的奖赏。}\par

没有什么能使我们的真实美德变得虚妄并吹散它们，如同记念自己所行的善事一样；因为这产生两种恶：既使我们懈怠，又使我们骄傲。因此，请看保禄，他知道我们的本性容易倾向懈怠，尽管他曾大大称赞斐理伯人，现在却通过其他许多上面的事，但主要藉着他当前的话来抑制他们的心思。这些话是什么呢？\BibleQuote{各位弟兄，我并不以为我已经夺得。} 但若保禄尚未夺得，对复活及将来之事并不自信，那么那些连他美德的极小部分都未达到的人，更难以自信了。也就是说，我认为我尚未获得全部美德，如同谈及一位赛跑者一样。他说，我尚未完成一切。若他在另一处说，\BibleQuote{这场好仗我已打完} \BibleRef{弟后 4:7}，但在这里却说，\BibleQuote{我并不以为自己已经夺得}；任何仔细阅读的人都会知道那些话和现在这些话的原因；因为不必一直停留在同一点上；并且他说这些话是在更早的时候，而另一些话则是在他临终前。但我只专注于\BibleQuote{一件事}，他说，\BibleQuote{努力向前，向着前面的事。} 但他又说，\BibleQuote{一件事，忘记背后，努力向前，向着标竿直跑，要得天主在基督耶稣里从高天召我来得的奖赏。} 因为使他向前追求前面的事的，正是他忘记背后的事。那么，凡以为一切都已完成，在成全美德上毫无欠缺的人，就可以停止奔跑，仿佛已夺得一切。但那些以为自己仍离目标尚远的人，将永不停止奔跑。因此我们应当时常思想这一点，即使我们行了无数的善事；因为若保禄在经历了无数的死亡、如此多的危险之后还这样想，我们岂不更当如此？他说，因为我虽跑了这么多仍然未能成功，但我没有气馁，也没有绝望，我仍然奔跑，我仍然努力。我只考虑这一件事：我确实能前进。我们也应如此行：忘记我们的成功，将它们抛在背后。因为赛跑者不会算他已跑了多少圈，而是算还剩多少圈。我们也应计算，不是我们在美德上前进了多少，而是我们还差多少。因为当所缺的尚未补足时，已完成的对我们有何益处？此外，他没有说我不计算，而是说我甚至不记念。因为当我们全力以赴于剩下的事，将其他一切都遗忘时，我们就变得热心了。\BibleQuote{努力向前}，他说；在我们到达之前，我们就努力获取。因为努力向前的人，即使双脚在奔跑，也尽力用身体的其他部分超越它们，将身体向前伸，伸出手臂，以便在赛程中多完成一些。这来自极大的热忱、极大的热情；因此赛跑者应当非常认真、非常热切地奔跑，毫不松懈。正如这样奔跑的人与仰卧的人之间的差别，保禄与我们之间的差别也是如此。他天天死，天天被认可，没有一时、没有一刻他的赛程不在前进。他不愿只取，而要抢夺奖赏；因为这样我们才能取得它。颁发奖赏者站在高处，奖赏安放在高天之上。\par

请看，这是何等遥远的路程！请看，这是何等高的攀登！我们必须藉圣神的翅膀向上飞，否则不可能越过此高度。我们必须带着身体前往那里，因为这是容许的。\Quote{我们的家乡在天上}\BibleRef{斐 3:20}，奖赏就在那里；你们看见那些赛跑者吗？他们如何按规矩生活，如何不接触任何消弱力量的事物，如何在师傅的指导下每天在角力场中锻炼，并遵守规则？效法他们，或更显出更大的热忱，因为奖赏并不相等：有许多人会阻碍你们；要按规矩生活：有许多事物会消弱你们的力量；要使你们的脚敏捷：因为这是可能的，不是出于本性，而是出于我们的意志。让我们使它轻快，免得其他事物的重量阻碍我们脚步的敏捷。教导你们的脚要稳健，因为有许多滑溜之处；你若跌倒，立刻失去很多。但即使你跌倒，要再站起来。即使如此，你仍能获胜。决不要尝试滑溜之物，你就不会跌倒；走在坚实的地上，抬起头来，抬起眼来；这些是教练向赛跑者发出的命令。这样，你的力量得到支持；但你若弯腰向下，你就跌倒，就松懈了。要向上看，那里有奖赏；看见奖赏，就增加了我们意志的决心。获得它的希望使人不感觉到劳苦，使路程显得短。这奖赏是什么？不是棕榈枝，而是什么？天国、永久的安息、与基督同享的荣耀、嗣业、兄弟情谊、无数的美善，其名不可胜数。那奖赏的美妙无法形容；只有拥有它的人知道，只有将要领受它的人知道。它不是金的，不以宝石镶嵌，它贵重得多。黄金与之相比是泥土，宝石与之相比是砖块。你若拥有它，并离开尘世前往天堂，你将能带着极大的尊荣在那里行走；天使将尊敬你，当你带着这奖赏时，你将能满怀信心地走近他们全体。\Quote{在基督耶稣内。}请看他的谦卑；他说，我这样做，\Quote{在基督耶稣内}，因为若无祂的推动，就不可能跨越如此巨大的距离：我们需要许多帮助，需要强大的同盟；祂愿意你在下面奋斗，在上面祂给你加冕。不像在这世上：冠冕不在这里，竞赛在这里，但冠冕在光明之处。你们岂不见，甚至在这里，最受尊敬的角力者和驾车者，不是在下方的赛场上被加冕，而是王召唤他们上去，在那里加冕他们？同样，在这里，你在天堂接受奖赏。\par

\BibleRef{斐 3:15}\Quote{所以，我们凡是成全的人，都应怀有这种心情，}他说。\Quote{若你们在什么事上存有另一种心情，天主也要将这事启示给你们。}是什么事？就是我们要\Quote{忘记背后的一切}。因此，成全的人的特点，就是不自以为成全。那么，你怎么说\Quote{凡是成全的人}呢？请告诉我，我们是否怀有你的心情？因为你若没有达到，也没有成为成全的，你怎么命令那些成全的人怀有像你——尚未成全的人——一样的心情呢？是的，他说，这就是成全。\Quote{若你们在什么事上存有另一种心情，天主也要将这事启示给你们。}也就是说，若有人以为自己已获得了一切美德。他警戒他们，不是直言，而是说什么？\Quote{若你们在什么事上存有另一种心情，天主也要将这事启示给你们。}请看他是多么谦卑地说这话！天主会教导你们，即天主会说服你们，不只是教导；因为保禄在教导，但天主会引领他们前进。他没有说\Quote{引领你们}，而是说\Quote{要启示}，这样这似乎更源自无知。这些话不是关于教义，而是关于生活的成全，以及我们不以为自己已成全；因为那以为自己已获得一切的人，其实一无所有。\par

\BibleRef{斐 3:16}\Quote{然而，我们达到了什么程度，就该照着同一规则行事，怀有同一心意。}\par

\Quote{然而，我们达到了什么程度。}这是什么意思？他说，我们要持守我们已经成功之处：爱、和睦与和平；因为在这方面我们已经成功了。\Quote{我们达到了什么程度：就要照着同一规则行事，怀有同一心意。}\Quote{我们达到了什么程度，}即在这方面我们已经成功了。你们看见吗？他愿意他的训诫成为我们的准则？准则既不容许增加，也不容许减少，因为那会破坏它作为准则的性质。\Quote{照着同一规则，}即照着同一信仰，在同一界限内。\par

\BibleRef{斐 3:17}\Quote{弟兄们，你们要一同效法我，也要留意那些按我们的榜样而行的人。}\par

他上面说过\Quote{提防犬类}，他把他们从那些人那里引开；他又把他们带到那些他们应当效法的人面前。他说，若有人愿意效法我，若有人愿意走同一条路，就让他留心他们；虽然我不在场，你们知道我的行走方式，即我生活中的行为。因为他不仅用言语教导，也用行为；如在歌咏团和军队中，其余的人必须效法歌咏团或军队的领队，从而秩序井然地前进。因为秩序可能因叛乱而解体。\par

宗徒们因此是一个\OriginalTerm{类型}{type}，并在整个过程中持守了某个\OriginalTerm{原型的典范}{archetypal model}。请思考他们的生活何等精准，以致他们被立为原型和榜样，如同活生生的法律。因为他们写在著作中的话，都在行为上向众人彰显了。这是最好的教导；如此，他才能带领他的门徒。但如果他嘴上说着高深的哲理，行为却相反，他就不再是老师了。因为仅仅是口头的哲学，连门徒也容易做到；但需要的是来自行为的教导和引领。因为这既使老师受人敬重，也使门徒预备好顺服。何以如此？当人看见他只在口头上谈论哲理，就会说他所命令的是不可能的事；他自己首先表明这些是不可能的，因为他没有实践。但如果他看见他的德行完全在行动中实行出来，他就不再能这样说了。然而，即使我们老师的生命是漫不经心的，我们也要留心自己，并听从先知的话：\BibleQuote{他们都要受天主的训诲。}\BibleRef{依 54:13}\BibleQuote{他们不再各教自己的邻人说：\InnerQuote{你们该认识上主}，因为从最小的到最大的，他们都要认识我。}\BibleRef{耶 31:34}你有一位不道德的老师吗？但你仍有那位真正是老师的，唯有他你才应称为老师。从他学习：他曾说：\BibleQuote{你们跟我学，因为我是良善心谦的。}\BibleRef{玛 11:29}所以，不要注意你的老师，而要留意他及其教训。从那里取得你的榜样，你有一个最好的典范，去效法它。在圣经中有无数的典范摆在你们面前，是圣善的生活；无论你愿意哪一个，来，追随老师之后，在门徒中找到它。一人以贫穷显出，另一人以财富；例如，\Person{厄里亚}以贫穷，\Person{亚伯拉罕}以财富。去那个你认为最容易、最适合自己去实践的典范。再次，一人以婚姻，另一人以童贞；\Person{亚伯拉罕}以婚姻，另一人以童贞。无论你跟随哪个，两者都通向天堂。一人以禁食发光，如\Person{若翰}；另一人没有禁食，如\Person{约伯}。再次，后者关心他的妻子、儿女、女儿、家庭，并拥有大财富；另一人除了毛衣之外一无所有。我为何要提家庭、财富、金钱呢？因为在王国中的人也可能持守德行，因为君王的家庭比任何私人家庭更充满烦恼。\Person{达味}就在他的王国中发光；紫袍和王冠并没有使他懈怠。另一个人被托付掌管全体人民，我是指\Person{梅瑟}，那是更困难的任务，因为那里的权力更大，所以困难也更大。你见过在财富中被赞许的人，也见过在贫穷中，在婚姻中，在童贞中；相反，看啊，有些人在婚姻和童贞中迷失，在财富和贫穷中迷失。例如，许多人在婚姻中丧亡，如\Person{三松}，但不是因婚姻，而是因他们自己的故意选择。同样在童贞中，如那五个童女。在财富中，如那轻视\Person{拉匝禄}的富人；在贫穷中，无数的穷人如今也丧亡。在王国中，我能指出许多丧亡的人，以及在治理人民中。你愿看见在士兵行列中被拯救的人吗？有\Person{科尔乃略}；在治理家庭中？有厄提约丕雅女王的太监。普遍如此。如果我们按适当的方式使用我们的财富，没有什么能毁灭我们；但如果不如此，一切都会毁灭我们，无论是王国、贫穷还是财富。但没有什么能伤害那保持警醒的人。\par

请告诉我，被掳有什么害处吗？毫无害处。因请你思考，\Person{若瑟}成了奴隶，却保持了德行。思考\Person{达内尔}和三个青年，他们成了俘虏，却越发闪耀，因为德行处处发光，不可战胜，没有什么能阻碍它。但我为何提贫穷、被掳、奴役；以及饥饿、疮痍和重病呢？因为疾病比奴役更难忍受。\Person{拉匝禄}如此，\Person{约伯}如此，\Person{弟茂德}也是如此，被\BibleQuote{屡次的软弱}\BibleRef{弟前 5:23}所困扰。你看见没有什么能胜过德行；无论是财富、贫穷、统治、服从、事务上的优先、疾病、轻蔑、遗弃。但德行将这一切留在地上和尘世，急奔向天堂。只要灵魂高贵，没有什么能阻止它成为有德行的。因为当工作的人充满活力时，外在无物能阻碍他；正如在技艺中，当工匠有经验、有毅力且精通其技艺时，即使疾病降在他身上，他仍然拥有它；即使他变穷了，他仍然拥有它；无论工具在手中或不在，无论他工作或不工作，他丝毫不会失去他的技艺：因为它的知识包含在他内。同样，有德行的人，献身于天主，若你把他投入财富或贫穷，投入疾病或健康，投入羞辱或大荣耀，他都会展现他的技艺。宗徒们岂非在各种境遇中行事，\BibleQuote{借着光荣和羞辱，借着美名和恶名}\BibleRef{格后 6:8}？这才是运动员，为一切准备好；因为德性的本性也是如此。\par

若你说：\Quote{我不能管理许多人，我应当过独居生活}，你便是侮辱了德性，因为德性能在任何状态中发挥作用，并透过一切而照耀；只要它存在于灵魂中。有饥荒吗？还是有富足？它表现出自己的力量，正如\Person{保禄}所说：\BibleQuote{我知道怎样处丰富，也知道怎样处缺乏。}\BibleRef{斐 4:12} 他需要工作吗？他不以为耻，反倒工作了两年。他需要忍受饥饿吗？他没有被压垮，也没有动摇。他需要承受死亡吗？他没有沮丧，反倒通过一切表现出他高贵的心志和技艺。因此让我们效法他，我们就不会再有悲伤的理由：因为告诉我，什么事能使这样的人悲伤？没有什么。只要没有人剥夺我们这门技艺，这人将成为所有人中最有福的，无论是今生还是来世。因为假设善人拥有妻子、儿女、财富和极大的尊荣，拥有这一切时他仍然同样有德。拿走这一切，他同样会有德，既不被不幸压倒，也不因顺境得意，而是如同磐石在汹涌的大海中与风平浪静时一样不动摇，既不被波浪打破，也不受平静影响，同样，坚固的心思在平静与风暴中都屹立不动。又如小孩子在船上航行时被抛来抛去，而舵手坐在一旁，笑着、不为所动，并欣喜地看着他们的慌乱；同样，真正明智的灵魂，当其他人都陷入混乱，或在不合时宜地因任何境况变化而微笑时，却仿佛坐在虔诚的舵柄和舵盘前，不为所动。因为告诉我，什么能扰乱虔诚的灵魂？死亡吗？那是更好生命的开始。贫穷吗？这帮助她走向德性。疾病吗？她不顾它的存在。她不顾安逸或痛苦；因为她预先对付它，已经自我痛苦。羞辱吗？世界已为她被钉在十字架上。失去子女吗？她不怕，当她完全相信复活时。那么什么能使她惊奇？这一切都不能。财富使她高涨吗？绝不，她知道金钱算不了什么。荣耀吗？她曾受教导：\BibleQuote{凡人的一切光荣都如草上的花。}\BibleRef{依 40:6} 奢侈吗？她曾听见\Person{保禄}说：\BibleQuote{那纵情享乐的，虽生犹死。}\BibleRef{弟前 5:6} 既然她既不燃烧也不受压，什么能与这样的健康相比呢？\par

其他灵魂却并非如此，而是比海或变色龙变化得更频繁，以致你大有理由微笑，当你看见同一个人时而大笑，时而哭泣，时而充满忧虑，时而又过度松弛和萎靡。为此，\Person{保禄}说：\BibleQuote{不要同化于这个世界。}\BibleRef{罗 12:2} 因为我们是天上子民，那里没有转变。不变的奖品向我们展示。让我们彰显我们这子民的身份，让我们从现在起就获得我们的善物。但为什么我们要将自己投入\Place{欧里波斯}，投入风暴，投入狂风，投入泡沫？让我们处于平静之中。这一切不取决于财富、贫穷、尊荣、羞辱、疾病、健康或软弱，而取决于我们自己的灵魂。如果它稳固并精通德性之学，一切对它将容易。甚至从现在起它就能看见它的安息，那个宁静的港湾，并且当它离开时，将在那里获得无量的善物，愿我们众人都能获得，赖我们的主\Person{耶稣基督}的恩宠和爱，借着祂，归于\Person{圣父}，偕同\Person{圣神}，获享光荣、权能、尊威，自今以至永远，及世之世。阿们。\par

\SourceRecord{Translated by John A. Broadus.；Nicene and Post-Nicene Fathers, First Series；Vol. 13.；Edited by Philip Schaff.；Buffalo, NY: Christian Literature Publishing Co.,；1889.；<http://www.newadvent.org/fathers/230212.htm>.}

% structure: p0001 source=h1 target=section reason=page-title

\section{斐理伯书第十三篇讲道词}

\BibleRef{斐 3:18-21}\par

\BibleQuote{因为有许多人，我时常告诉你们，现在又含泪告诉你们，他们是基督十字架的仇敌：他们的结局是丧亡，他们的神是自己的肚腹，他们的光荣在自己的羞辱，他们只思念地上的事。但是我们的家乡在天上，从那里我们也等待救主，主耶稣基督；他将改变我们卑贱的身体，使之相似他光荣的身体，按照他能将万有屈服于自己的大能。}\par

在基督徒身上，没有什么比追求安逸和休息更不协调、更与他的身份相悖的了；沉溺于现世生活，与我们所宣认的信仰和入伍的使命格格不入。你的主被钉在十字架上，你却追求安逸？你的主被钉子刺穿，你却生活奢华？这些事与一位高尚的士兵相称吗？因此保禄说：\BibleQuote{有许多人，我时常告诉你们，现在又含泪告诉你们，他们是基督十字架的仇敌。}因为有些人假装信奉基督教，却生活在安逸和奢华之中，这与十字架相悖，所以他这样说了。因为十字架属于灵魂在战斗岗位上，渴望死亡，不求安逸，而他们的行为却恰恰相反。因此即使他们说自己是基督的，仍然像是十字架的仇敌。因为他们若爱十字架，就会努力过被钉死的生活。你的主不是被挂在木头上吗？你当从另一面效法他。即使没有人钉死你，你也要钉死自己。钉死自己，不是要杀死自己（千万不可！那是邪恶的事），而是正如保禄所说：\BibleQuote{世界已被钉死于我，我也被钉死于世界。}\BibleRef{迦 6:14}你若爱你的主，就当死于他的死。要学习十字架的力量何其伟大！它成就了多少善事，并且仍在成就：它是我们生命的保障。万有都藉着它而成就。圣洗圣事藉着十字架，因为我们必须接受那印记。覆手礼藉着十字架。无论我们旅行，还是在家，无论我们在哪里，十字架都是一件巨大的善事，是救恩的铠甲，是不能被击倒的盾牌，是抵抗魔鬼的武器；当你与他敌对时，你就背负了十字架，不仅仅是在你以十字架划印时，更是当你承受属于十字架的一切时。基督认为我们的苦难配得上称为十字架。正如他说：\BibleQuote{除非人背起自己的十字架跟随我}\BibleRef{玛 16:24}，即除非他准备好死去。\par

但这些人是卑劣的，贪生爱命，迷恋自己的身体，他们是十字架的仇敌。凡是喜爱奢华、贪图现世平安的人，就是那十字架的仇敌——保禄以此为荣，他拥抱它，渴望与它结合。正如他说：\BibleQuote{我被钉死于世界，世界也被钉死于我。}但在这里他说：\BibleQuote{我现在含泪告诉你们。}为什么呢？因为邪恶紧迫，因为这样的人值得流泪。确实，奢华的人值得流泪，他们养肥了那裹着他们的东西——我指身体——却不关心那必须交账的灵魂。看，你生活奢华，你醉酒，今日明日，十年、二十年、三十年、五十年、一百年，这是不可能的；但如果你愿意，让我们假设。结局是什么？收益是什么？什么也没有。这样的生活难道不值得流泪和哀悼吗？天主把我们带进这场赛跑，为要给我们加冕，我们却在没有做任何高尚行动的情况下离开了。所以，别人欢笑享乐时，保禄却流泪。他是如此富有同情心，如此为众人挂虑。他说：\BibleQuote{他们的神是自己的肚腹。}他们竟以此为神！这就是说：\Quote{让我们吃喝吧！}你看到奢华是多么大的邪恶吗？对某些人来说，财富是神；对另一些人来说，肚腹是神。这些不也是偶像崇拜者，而且比普通人更糟糕吗？而且他们的\BibleQuote{光荣在自己的羞辱。}\BibleRef{格前 15:32}有人说这是指割损。我不这么认为，但意思就是：他们以自己本应羞愧的事为荣。做可耻的事是可怕的；然而做了并且感到羞愧，那只有一半可怕。但若一个人竟以此自夸，那就是极度的愚昧了。\par

这些话难道只适用于他们吗？难道在场的人就不受责难吗？难道没有人要为此交账吗？难道没有人以肚腹为天主，或以自己的羞辱为荣耀吗？我愿，我诚切希望，这些指控没有一条落在我们身上，我也希望不知道任何人牵涉在我所说的事中。但我恐怕这些话对我们要比对那些人更适用。因为当一个人将一生消耗在饮酒作乐中，只在穷人身上花费一点小钱，却将大部分消耗在肚腹上，这些话岂不正好适用于他吗？没有什么话比这些话更能引人注意，更刺人心扉了：\BibleQuote{他们的天主是肚腹，他们的荣耀在于他们的羞辱。}这些人是谁？他说，是那些思念地上事的人。\Quote{让我们建造房屋。}我问：在哪里？在地上，他们回答。让我们购买田产：又在地上。让我们获得权力：又在地上。让我们获取荣耀：又在地上。让我们致富：这一切都在地上。这些人的天主就是他们的肚腹；因为他们没有属神的思想，而将所有财产放在此处，并思念这些事，他们以肚腹为天主是合理的，正如他们说：\Quote{我们吃喝吧，因为明天我们要死了。}至于你的身体，你因它出于尘土而悲伤，请告诉我，即便如此你也没有受到任何伤害。但你的灵魂，你却把它拖到地上，而你本应使你的身体也成为属神的；因为只要你愿意，你就能做到。你领受了一个肚腹，是为了喂养，而不是撑大它；是为了管理它，而不是让它作你的主母；是为了它服务于你，滋养其他部分，而不是你服务于它，不是让你超出界限。大海越出其边界时，不会造成像肚腹加于我们身体和灵魂那样多的祸害。前者淹没整个地面，后者淹没整个身体。给肚腹设定节制为界限，如同天主为大海设定了沙滩。然后，如果它的波浪兴起，狂怒不止，就用你内在的力量斥责它。看，天主如何尊荣了你，使你能够效法祂，你却不效法；但你看见肚腹泛滥，毁坏并压倒你的整个本性，却不敢约束或节制它。\par

\BibleQuote{他们的天主，}他说，\BibleQuote{是肚腹。}让我们看看保禄如何服侍天主；让我们看看贪食者如何服侍他们的肚腹。他们不是经历成千上万这样的死亡吗？他们不是害怕违背肚腹的任何命令吗？他们不是为肚腹效力不可能的事吗？他们岂不比奴隶更不堪？\BibleQuote{但我们的家乡，}他说，\BibleQuote{是在天上。}因此，我们不要在这里寻求安逸；我们在那里发光，那里也是我们的家乡。\BibleQuote{从那里，}他说，\BibleQuote{我们等候救主，}主耶稣基督：\BibleQuote{他将改变我们卑微的身体，使其相似于他荣耀的身体。}他一步一步地将我们提升。他说\Quote{从天上}和\Quote{我们的救主，}从地点和人物显示主题的尊贵。\BibleQuote{他将改变我们卑微的身体，}他说。身体如今遭受许多苦难：被锁链捆绑，被鞭打，承受无数灾祸；但基督的身体也遭受了同样的苦难。他暗示这一点时说：\BibleQuote{使其相似于他荣耀的身体。}因此，身体是同一个，但穿上了不朽。\Quote{将改变。}因此，形状不同了；或者，他可能比喻性地谈论了这种改变。\par

他说，\BibleQuote{我们卑微的身体，}因为它如今卑微，受制于败坏、痛苦，因为它似乎无价值，与其他动物没有区别。\BibleQuote{使其相似于他荣耀的身体。}什么？我们这身体将变得像那坐在圣父右边的那一位，像那受天使朝拜、无形体之能力侍立在他面前的、超越一切统治、权势和德能的那一位？如果整个世界都为那些已从这希望中堕落的人哀哭叹息，它能配得上这样的哀哭吗？因为，当赐给我们身体相似于他的应许时，它却仍与魔鬼一同离去。我从此不在乎地狱；不论能说什么，既然从如此大的荣耀中堕落，现在及今后，请把地狱视为与这堕落相比微不足道的东西。你说什么，保禄？相似于他？是的，他回答；然后，免得你不信，他加了一个理由：\BibleQuote{按照他能使万有归顺自己的大能。}他说，他有能力使万有归顺自己，包括败坏和死亡。或者说，他用同样的能力成就这事。告诉我，哪一样需要更大的能力：使魔鬼、天使、总领天使、革鲁宾、色辣芬归顺，还是使身体成为不朽和不死的？后者当然比前者更大；他显示了更大能力的工作，使你也相信这些。因此，虽然你看到这些人喜乐并受尊荣，仍要站稳，不要被他们绊倒，不要动摇。这些希望足以振作最懒惰和懈怠的人。\par

第四章 第一节。\BibleQuote{因此，}他说，\BibleQuote{我所亲爱和想念的弟兄，我的喜乐和我的冠冕，你们要这样在主内站稳，亲爱的。}\BibleRef{斐 4:1}\par

\Quote{这样。}如何？不动摇。看，他在劝勉之后如何加上赞美：\Quote{我的喜乐和我的冠冕，}不单是喜乐，也是荣耀，不单是荣耀，也是我的冠冕。这荣耀无物可比，因为它是保禄的冠冕。\Quote{你们要这样在主内站稳，亲爱的，}即在天主的希望中。\par

第2、3节。\BibleQuote{我劝告厄敖狄雅，也劝告辛提赫，要在主内同心合意。是的，我也请求你，真正的同侪，帮助这些妇女。}\par

有人说保禄在此劝勉自己的妻子；但事实并非如此，而是某个别的妇人，或是其中一位的丈夫。\BibleQuote{你们要帮助这些妇人，因为她们曾与我在福音上一同劳苦，并与克肋孟以及其余的同工，他们的名字都写在生命册上。}你是否看见他对她们的德行作了何等大的见证？因为正如基督对祂的宗徒们所说：\BibleQuote{不要因为魔鬼屈服于你们而欢喜，却要因为你们的名字已登记在天上而欢喜}\BibleRef{路 10:20}；同样，保禄也向她们作证说：\BibleQuote{她们的名字写在生命册上。}在我看来，这些妇人是当地教会的主要成员，他将她们托付给一位他称之为\OriginalTerm{同轭者}{Syzygus}的显要人物，或许他惯常将她们托付给他，如同托付给一位同工、战友、兄弟与伙伴，正如他在《罗马书》中所说：\BibleQuote{我把我们的姊妹非比推荐给你们，她是耕格勒教会的仆役。}\BibleRef{罗 16:1}\BibleQuote{同轭者；}要么是她们的一位兄弟，或是其中一位的丈夫；好像他曾说：如今你是一位真正的兄弟，如今你是一位真正的丈夫，因为你已成为一个肢体。\BibleQuote{因为她们曾与我在福音上一同劳苦。}这保护来自家庭，并非出于友情，而是出于善行。\BibleQuote{与我一同劳苦。}你说什么？妇人与你一同劳苦？是的，他回答，她们也贡献了不小的部分。虽然有许多人与他一同工作，但这些妇人也与他在众人中行动。那时教会因这而大得造就，因为有许多益处产生，当那些被认可的，无论是男人还是女人，都从其余人获得如此的尊敬。因为首先，其余人被激发起同样的热忱；其次，他们也因所显示的尊敬而获益；第三，这使得那些人更加热忱和恳切。因此，你看保禄处处关心此事，并称赞这样的人以示重视。正如他在《格林多前书》中所说：\BibleQuote{他们是阿哈雅初结果子的人，}\BibleRef{格前 16:15}有人说\BibleQuote{同轭者}（Syzygus）是一个专名。那又如何呢？不论是否如此，我们不必精确探究，只需注意他下令要使这些妇人受到极大的保护。\par

他说，我们所有的一切都在天上：我们的救主、我们的城、凡能命名的：\BibleQuote{我们等待救主，主耶稣基督从那里降来。}这是祂的恩慈与对人的爱。祂自己再次来到我们这里，并不将我们拖到那里，而是携带着我们，就这样与我们一同离开。这是极大的荣耀；因为如果我们还是仇敌时祂就来到我们这里，那么现在我们成为朋友，祂岂不更来临？祂不把这任务交给天使，也不交给仆人，而是亲自来召唤我们进入祂的宫殿。看，我们也将\BibleQuote{被提到云中}\BibleRef{得前 4:17}，以表对祂的尊敬。\par

那么，谁能被找着是\BibleQuote{忠信而精明的仆人}？谁是那些堪当如此福乐的人？那些失败的人是多么可怜！即使我们永远哭泣，难道我们能做出与这情景相称的事吗？因为如果你提及无数的地狱，你也说不出比灵魂在下列情境中承受的痛苦更大的事：当全世界陷入混乱，当号角吹响，当天使们冲上前去，第一队，然后第二队，然后第三队，然后成千上万的队伍倾泻到地上；然后\OriginalTerm{革鲁宾}{Cherubim}（众多且无穷）；\OriginalTerm{色辣芬}{Seraphim}；当祂自己带着不可言喻的光荣降临时；当那些曾去聚集被选者到中央的人迎接祂时；当保禄及其同伴，以及所有在他那时代被认可的人，被加冕，被公开宣告，被君王在祂所有天军面前尊荣时。因为即使没有地狱，这是一件何等可怕的事：一部分人被尊荣，另一部分人被羞辱！我承认，地狱是难以忍受的，是的，极难以忍受，但比地狱更难忍受的是丧失天国。试想：若有一位君王或王子，出发后历经无数战争大获全胜，成为众人钦佩的对象，如今率领全军返回某城，乘着战车，带着战利品，带着无数金盾的队伍，带着长矛手，所有卫兵围着他，全城装饰着花环，世上所有统治者都陪同他，所有异国军队都作为俘虏跟随他，还有总督、省长，当着所有统治者及那一切辉煌，他接纳前来迎接的市民，亲吻他们，伸手赐予他们自由晋见，与他们交谈，所有旁立者如同朋友，告诉他们那次旅程全是为他们而行的，然后引他们进入王宫，与他们分享，即使其余人不被惩罚，这岂不是与极重的惩罚相等？但若在人间，失去如此光荣已是痛苦之事，何况在天主面前，当所有天上的掌权者与君王同在，当魔鬼被捆绑、低头，连魔鬼本身也被锁链牵来，所有敌对势力，当诸天的掌权者，当祂自己乘云降临时，更是如此。\par

请相信我，我甚至无法说完我的话，因为述说此事时攫取我灵魂的悲伤。想一想，当我们本可以不被剥夺却要被剥夺何等的荣耀！这正是悲惨：我们受这些苦，而我们本可以不受苦。当祂接纳一部分人，带领他们到天堂的圣父那里，而拒绝另一部分人——天使们抓住他们，不顾他们的意愿，哭泣着、垂头丧气地拖向地狱之火，他们先被示众给全世界——你认为，那时是何等悲伤？那么，趁着还有时间，让我们赶紧，为自己的救恩深思熟虑。我们将有多少话要说，如同那富人一样？\Quote{若有人现在容忍我们，我们必会考虑有益之事！} 但没有人容忍我们。我们将会这样说，这不仅是藉着他显明，也藉着许多其他人。为使你们明白这一点，有多少人曾发烧，并说：\Quote{我们若痊愈，必不再陷入同样的状态。} 那时我们也会说许多这样的话，但我们会如那富人一般被回答：有深渊隔开，我们在此已享用了我们的福乐。\BibleRef{路 16:25} 那么，我恳求你们，让我们痛苦地呻吟吧；不如说，让我们不仅呻吟，还要追求德行；让我们现在为救恩哀哭，以免那时徒然哀哭。让我们现在哭泣，不要在那时为我们的厄运哭泣。这哭泣是出于德行，那哭泣是无益的悔改；让我们现在受苦，以免那时受苦；因为在此处和在那里受苦并不相同。在此处，你受苦片刻，更确切地说，你感觉不到自己的痛苦，因为你知道受苦是为你好。但在那里，痛苦更为剧烈，因为没有希望，没有逃脱，而是无限制且永无止境。\par

但愿我们都能摆脱这一切，获得赦免。但让我们祈祷并勤勉，好使我们获得赦免。我恳求你们，让我们勤勉；因为如果我们勤勉，我们甚至能藉祈祷得胜：如果我们恳切祈祷，天主必会允准我们的请求；但如果我们不求祂，也不认真地做这类事，也不努力，我们怎么可能成功呢？靠睡觉吗？绝不可能。因为即使藉着奔跑、伸开双手、效法祂的死亡，正如保禄所说，我们才能成功，更不用说睡觉了。他说：\Quote{我或许能到达。} 但如果保禄说：\Quote{我或许能到达，} 我们又该说什么呢？因为靠睡觉，连世俗的事务都无法完成，更不用说超性的事了。靠睡觉，甚至连朋友那里也得不到什么，更何况从天主呢？连父亲都不尊重那些睡觉的人，更何况天主呢？让我们劳苦一小会儿，好得永远的安息。我们无论如何都必须受苦。如果我们不在此受苦，那里就有苦难等着我们。为何我们不选择在此受苦，好使我们在那里得享安息，并获得那不可言宣的祝福，在基督耶稣内，愿光荣、权能和尊荣归于祂，并与圣父及圣神，从今世直到永远。阿们。\par

\SourceRecord{Translated by John A. Broadus.；Nicene and Post-Nicene Fathers, First Series；Vol. 13.；Edited by Philip Schaff.；Buffalo, NY: Christian Literature Publishing Co.,；1889.；<http://www.newadvent.org/fathers/230213.htm>.}

% structure: p0001 source=h1 target=section reason=page-title

\section{斐理伯书第十四篇讲道}

斐理伯书 4:4–7 \BibleRef{斐 4:4}\par

\BibleQuote{你们要常在主内喜乐；我再说，你们要喜乐。让你们的宽仁为众人所知。主已临近。你们什么也不要挂虑，但在一切事上，借祈祷和恳求，怀着感恩之心，使你们的请求呈于天主。那样，天主那超越一切理解的平安，必在基督耶稣内，固守你们的心思念虑。}\par

\Quote{\Emph{哀恸}的人是有福的，} 并 \Quote{你们有祸了，你们现今欢笑的人} \BibleRef{玛 5:4}，基督说。那么保禄怎能说 \BibleQuote{你们要常在主内喜乐} 呢？基督说：\BibleQuote{你们有祸了，你们现今欢笑的人}，这是指属于现世事物的欢笑。祂也祝福那些哀恸的人，并非仅因失去亲人而哀恸，而是那些心中刺痛的人，他们为自己的过错哀伤，并清算自己的罪，甚至为别人的罪而哀伤。这种喜乐并不与那哀恸相悖，反而正是从哀恸生出的。因为那为自己的过错忧伤并承认的人，是喜乐的。此外，我们为自己的罪忧伤，却能在基督内喜乐。既然他们因所受的苦难而困扰，\BibleQuote{因为你们蒙恩，不仅为相信基督，也要为祂受苦} \BibleRef{斐 1:29}，因此他说：\BibleQuote{你们要常在主内喜乐。} 这只能意味着：如果你们活出这样的生命，就能喜乐。或者，当你们与天主的共融不受阻碍时，就喜乐。又或者，\BibleQuote{在} 这个词可解释为 \BibleQuote{与}，仿佛他说：与主同在。\BibleQuote{常在我再说，你们要喜乐。} 这是安慰之辞；例如，那在天主内的人常喜乐。是的，即使他受患难，无论遭遇什么，这样的人总是喜乐的。听听路加所说：\BibleQuote{他们离开公议会，心里欢喜，因为他们配为这名字受辱。} \BibleRef{宗 5:41} 如果鞭打和锁链，这些看来最痛苦的事，能带来喜乐，还有什么能使我们忧愁呢？\par

\BibleQuote{我再说，你们要喜乐。} 他重复这句话，做得很好。因为事物的本性带来忧伤，他借重复表明，他们无论如何都应喜乐。\par

\Quote{让你们的宽仁为众人所知。} 他在前面说：\BibleQuote{他们的天主是肚腹，他们的光荣在于他们的羞耻}，以及他们 \BibleQuote{思念地上的事} \BibleRef{斐 3:19}。他们很可能与恶人结仇；因此他劝勉他们不要与他们同流合污，却要以全部宽仁对待他们，不仅对待弟兄，也对待仇敌和反对者。\Quote{主已临近，什么也不要挂虑。} 因为，请告诉我，他们起来反对吗？如果你见他们生活奢华，为什么忧愁？审判已近了；他们不久就要为他们的行为交账。你们受困苦，他们享奢华？但这些事不久就要结束。他们图谋反对你们，威胁你们？\Quote{什么也不要挂虑。} 审判已在眼前，那时这些都将反转。\Quote{什么也不要挂虑。} 如果你们善意对待那些谋害你们的人，但最终这不会对他们有益。报偿已在眼前，无论你们遭遇贫穷、死亡或其他可怕的事。\Quote{但在一切事上，借祈祷和恳求，怀着感恩之心，使你们的请求呈于天主。} 这里有第一个安慰：\Quote{主已临近。} 又有：\BibleQuote{我天天与你们同在，直到世界的终结。} \BibleRef{玛 28:20} 看，另一个安慰，一种医治忧伤、困苦和一切痛苦的良药。这是什么？祈祷，凡事感恩。祂愿意我们的祈祷不仅是请求，也要为我们已得的感恩。因为那不为过去感恩的人，怎能祈求将来？\Quote{在一切事上，借祈祷和恳求。} 所以我们应为一切事感恩，甚至那些看来痛苦的事，因为这是真正感恩之人的本分。在别的情况中，事物的本性要求如此；但这出自感恩的灵魂，及那热切向往天主的人。天主接纳这样的祈祷，其他的祂不承认。献上可蒙悦纳的祈祷；因为祂安排一切为我们的益处，即使我们不知道。这就是它大有益处的证明，即我们不知道。\Quote{那超越一切理解的天主的平安，必在基督耶稣内固守你们的心思念虑。} 这是什么意思？\Quote{天主的平安} 即祂对人所成就的和平，超越一切理解。因为谁能期待，谁能希望如此美好的事发生？它们不仅超出人的言语，也超出人的理解。因为祂为自己的仇敌、为那恨祂的人、为那决意离弃祂的人，不惜交付祂的独生子，为使我们与祂和好。因此，这平安，即和好、天主的爱，必固守你们的心思念虑。\par

因为这正是教师的本分，不仅劝勉，也要祈祷，并藉恳求协助，使他们不致被试探所压倒，也不致被诡计所迷惑。就如他曾说：\Quote{愿那以人心无法理解的方式拯救了你们的那一位，亲自守护你们，保全你们，使你们不受任何损害。}他或者是指这个，或者是指基督所说的那种平安：\Quote{我把平安留给你们，我将我的平安赐给你们}\BibleRef{若 14:27}；这平安将守护你们，因为这平安超越一切人的理解。如何超越？当他吩咐我们要与仇敌、与那亏待我们的人、与那些与我们作战为敌的人和好时，这岂非超越人的理解吗？但让我们先看前一点。如果平安超越一切理解，那么大赐平安的天主本身更超越一切理解，不仅是我们的理解，也超越天使和天上万有的理解。\Quote{在基督耶稣内}是什么意思？将在祂内守护我们，使你们得以坚定，不致从祂的信德中坠落。\par

第8节。\Quote{最后，弟兄们，凡是可敬的，凡是真实的，凡是公义的。} \Quote{最后}是什么意思？相当于\Quote{我已说尽一切}。这是急于前行、不涉今世之人的口吻。\BibleRef{斐 4:8}\par

\Quote{最后，弟兄们，凡是可敬的，凡是真实的，凡是公义的，凡是纯洁的，凡是可爱的，凡是有美名的；若有什么德行，若有什么可称赞的，你们要思念这些事。}\par

第9节。\Quote{你们所学习、所领受、所听见、所看见的，在我身上的一切。}\par

\Quote{凡是可爱的}是什么意思？对信友是可爱的，对天主是可爱的。\Quote{凡是真实的}。德行是真实的，恶习是虚假的。因为它的快乐是虚假，它的荣耀是虚假，世上的一切都是虚假。\Quote{凡是纯洁的}。这与\Quote{思念地上的事}相对。\Quote{凡是可敬的}。这与\Quote{他们的天主是肚腹}相对。\Quote{凡是公义的}，他说，就是\Quote{凡是有美名的}。\Quote{若有什么德行，若有什么可称赞的}。这里他也愿意他们顾及关乎人的事。\Quote{你们要思念这些事}，他说。你看，他切愿驱除我们灵魂中一切恶念，因为恶行出自恶念。\Quote{你们所学习、所领受的}。这是教训：在所有劝勉中，他以自己为模范；正如他在另一处说：\Quote{正如你们以我们为榜样}\BibleRef{斐 3:17}。而在此处又说：\Quote{你们所学习、所领受的}，即由口头所教导的，\Quote{所听见、所看见的，在我身上}：无论是我的言语、行为或品行。你看见吗？关于每件事，他都给我们这些命令？因为不可能详细列举一切——我们的进来、出去、言谈、举止、交往（因为基督徒需要思虑所有这些事），所以他简短地、几乎以总结的方式说：\Quote{你们在我身上所听见、所看见的。}我以言行引导了你们。\Quote{你们要实行这些事}，不单是言语，也要\Emph{实行}。\Quote{这样，赐平安的天主必与你们同在}，意思是：你们将在平安中，在极大的安全中，不受任何痛苦或违愿之事。因为当我们与祂和好时——而这和好是通过德行达成的——祂更要与我们和好。因为那如此爱我们，甚至在我们违逆时仍施恩于我们的，岂不更在我们奔向祂时，更加显示祂对我们的爱吗？\par

没有任何东西像恶习那样是我们本性的仇敌。从许多事上可见，恶习如何与我们为敌，德行如何与我们为友。你愿意我谈论奸淫吗？它使人遭受羞辱、贫穷、嘲笑，在众人眼中可鄙，正如仇敌对待他们。它常使人陷入疾病和危险；许多人因情妇而丧生或受伤。若奸淫产生这些，通奸就更甚了。而施舍也是如此吗？绝非如此。反之，就像慈爱的母亲使儿子处于极大的得体、秩序、美誉中，并让他从容从事必要的工作，施舍并不释放我们或使我们脱离必要的工作，反而使灵魂更加明智。因为没有什么比情妇更愚蠢了。\par

但你愿意看看贪婪吗？它也像仇敌对待我们。如何？它使我们被众人憎恨。它使众人都自夸胜过我们；既包括那些曾受我们不公对待的人，也包括那些未曾受害的，他们同情前者，并为自己担忧。众人都视我们为公敌，如同野兽，如同魔鬼。处处有无数的控告、谋害、嫉妒，都是仇敌的行为。但公义相反，使众人成为朋友，众人和善，众人善待我们；众人为我们祈祷；我们的事完全安全，没有危险，没有猜疑。甚至睡眠也带着完全的安全无忧地来临，没有忧虑，没有哀哭。\par

这种生活是多么美好！还有什么呢？最好是嫉妒，还是彼此同乐？让我们探究这一切，就会发现德行如同一位真正仁慈的母亲，使我们处于安全之中；而邪恶则是背信弃义、充满危险的东西。请听先知的话，他说：\BibleQuote{主是敬畏祂者的堡垒，祂的盟约是为了指示他们。}\BibleRef{咏24:14}一个人若心中没有意识到任何邪恶，就谁也不怕；相反，生活在罪恶中的人从不自信，反而在自己的家人面前颤抖，用怀疑的眼光看着他们。何止家人？他连自己良心的法庭也承受不了。不仅外在的人，而且他内心的思想也同样影响他，使他不得安宁。那么，\Person{保禄}说什么呢？我们应当依靠赞美而活吗？他没有说要看重赞美，而是要行可赞美的事，但并非为了赞美。\par

\Quote{凡是真实的}——因为我们所说的是虚假的。\Quote{凡是可敬的。}\Quote{可敬的}属于外在的德行，\Quote{纯洁的}属于灵魂。他说，不要给人绊脚的理由，也不要给人指责的把柄。因为他曾说：\Quote{凡是美名的}——免得你以为他仅指在人眼中看为美的事——他又接着说：\Quote{若有什么德行，若有什么可称赞的，你们要思念这些事}——要行这些事。他愿意我们常在这些事上，关心这些事，思念这些事。因为如果我们彼此和睦，天主也会与我们同在；但如果我们挑起战争，和平的天主就不会与我们同在。因为没有什么比邪恶更敌视灵魂。也就是说，和平与德行使灵魂处于安全之中。因此，我们必须在我们这一方开始，然后才能吸引天主到我们身边。\par

天主不是战争和争斗的天主。使战争和争斗止息，无论是敌对祂的，还是敌对邻人的。你要与众人和睦相处，思想天主以怎样的品性拯救你。\BibleQuote{缔造和平的人是有福的，因为他们要称为天主的子女。}\BibleRef{玛5:9} 这样的人常常效法天主子：你也要效法祂。要和睦。你的兄弟越与你争斗，你的赏报就越大。因为请听先知的话：\BibleQuote{我与厌恶和平的人，仍是和平的。}\BibleRef{咏119:7} 这就是德行，这超越人的理解，这使我们接近天主；没有什么比不记仇更令天主喜悦。这释放你脱离你的罪恶，这解除对你的控告；但如果我们争斗和殴打，我们就远离天主：因为敌意由冲突产生，而从敌意中生出记仇。\par

砍断根，就没有果实。这样我们就会学会轻视今生的事物，因为在属灵的事上并没有冲突，完全没有；但无论你看到什么，无论是冲突、嫉妒，或任何人能提及的，这一切都源自今生的事物。每一个冲突，其开端要么是贪财，要么是嫉妒，要么是虚荣。因此，如果我们和睦相处，我们就会学会轻视地上的事物。有人偷了我们的钱？他并没有损害我们，只要他不偷我们上面的宝藏。他妨碍了你的荣耀？但这荣耀不是来自天主的，而是不足挂齿的。因为这根本不是荣耀，只是荣耀的虚名，甚至是一种羞辱。他偷了你的荣誉？这不是你的，而是他自己的。因为行不义的人，与其说是施加不义，不如说是承受不义；同样，谋害邻人的人，首先毁灭自己。\par

因为\BibleQuote{挖坑的，反掉在其中。}\BibleRef{箴26:27}因此，我们不要图谋害人，免得伤害自己。当我们破坏他人声誉时，要想到我们是在伤害自己，我们是在谋害自己。因为或许在人间我们若有能力，确实能伤害他，但在天主面前，我们却因激怒祂而伤害了自己。所以我们不要伤害自己。因为正如我们伤害邻人时是在伤害自己，同样，加惠他们就是加惠自己。所以，如果你的敌人伤害了你，你若明智，他其实有益于你，因此不要以同样的事回报他，反而要向他行善。但你说，打击仍然严重。那么想一想：你并不是在帮助他，而是在惩罚他，并且使你自己得益，这样你很快就会向他行善。那么怎样？我们应当出于这个动机行动吗？我们本不该出于这个动机，但如果你心不接受其他理由，他说，即使凭借这个理由，你也可以迅速说服它放弃敌意，并且将来会像对朋友一样向敌人行善，并会获得将来的美好事物。愿天主使我们众人都在基督耶稣内获得这些。阿们。\par

\SourceRecord{Translated by John A. Broadus.；Nicene and Post-Nicene Fathers, First Series；Vol. 13.；Edited by Philip Schaff.；Buffalo, NY: Christian Literature Publishing Co.,；1889.；<http://www.newadvent.org/fathers/230214.htm>.}

% structure: p0001 source=h1 target=section reason=page-title

\section{斐理伯人书第十五篇讲道}

\BibleRef{斐 4:10-14}\par

\BibleQuote{但我在主内大大喜乐，因为如今你们对我的思念终于再次萌发；你们固然一向思念，只是缺少机会。我并不是因缺乏而说这话，因为我已学会了，无论在什么境况，都可以知足。我知道怎样处卑贱，也知道怎样处富余；在各种事上，我学会了秘诀：饱足与饥饿，富余与缺乏。我靠着加给我力量的，凡事都能做。然而，你们却作得好，因为你们分担了我的患难。}\par

我屡次说过，施舍不是为了领受者的缘故，而是为了施与者的缘故，因为后者才是获得最大益处的人。\Person{保禄}在此也表明了这一点。怎么表明呢？斐理伯人经过很长时间后，送了些东西给他，并将这托付给\Person{厄帕夫洛狄托}。请看，当他即将派\Person{厄帕夫洛狄托}作为这封书信的送信人时，他称赞他们，并指出这个行动不是为了领受者的需要，而是为了施与者的需要。他这样做，既使那施恩于他的人不因骄傲而自高，也使他们在行善上更加热忱，因为他们反而使自己受益；也使那些领受的人不至于无所畏惧地冲上前去领受，以免遭受谴责。因为主说：\BibleQuote{施予比领受更为有福。}\BibleRef{宗 20:35}那么，他为什么说\Quote{我在主内大大喜乐}呢？他说，不是带着世俗的喜乐，也不是今生的喜乐，而是在主内。不是因为我自己得了安慰，而是因为你们长进了；因为这是我的安慰。因此他又说\Quote{大大}；因为这喜乐不是属肉体的，也不是为了自己的安慰，而是为了他们的长进。\par

且看他如何温和地责备了他们以往的时间后，又迅速将这事遮掩，并教导他们要常常、始终持守善行。他说：\Quote{因为如今终于}。\Quote{终于}，这词表明过了很长时间。\Quote{你们再次萌发}，如同曾发芽、枯萎后又发芽的果实。这里他表明，他们起初茂盛，然后衰败，现在再次发芽。因此，\Quote{再次萌发}一词既有责备，也有称赞。因为枯萎的人能够再次发芽，这不是小事。他也指出，这一切发生在他们身上是由于懈怠。但这里他表明，即使在以前，他们也惯于在这些事上热心。因此他加上：\Quote{你们对我的思念，你们固然一向思念。}免得你以为他们在其他事上也更为热心，然后枯萎，唯独在这件事上没有，看他如何补充说：\Quote{你们对我的思念。}我将\Quote{如今终于}一词仅用于此事；因为在其他事上并非如此。\par

在此有人可能会问：当他曾说\BibleQuote{施予比领受更为有福}\BibleRef{宗 20:25}，以及\Quote{我这两只手常供给我和同人的需用}，并且写信给格林多人时又说\BibleQuote{我宁可死，也不叫人使我所夸的落了空}\BibleRef{格前 9:15}，他如何容许自己所夸的落空呢？怎么落空？通过领受。因为如果他所夸的是他不领受，他现在怎能容忍这样作呢？那么，这是什么意思呢？大概，他那时不领受是因为假宗徒的缘故，他说：\BibleQuote{好使他们在所夸的事上，与我们相同}\BibleRef{格后 11:12}。他不是说\Quote{是}，而是说\Quote{所夸的}；因为他们暗中领受。因此他说\Quote{他们所夸的}。所以他又说：\BibleQuote{没有人能阻止我这夸口}\BibleRef{格后 11:10}。他不是简单地说\Quote{不能阻止}，而是\Quote{在\Place{阿哈雅}地区}。又说：\BibleQuote{我抢夺了别的教会，向他们取了工价来，为给你们服役}\BibleRef{格后 11:8}。这里他表明他确实领受了。但\Person{保禄}确实正当地领受，因为他有如此大的工作；如果他真的领受了的话。但那些不工作的人，如何能领受呢？有人说：\Quote{但我祈祷。}但没有工作。因为这事可以与工作一同进行。\Quote{但我禁食。}这也不是工作。请看这位有福的人，在多地传道，同时也在工作。\Quote{你们却缺少机会。}\Quote{缺少机会}是什么意思？他说，这不是出于懈怠，而是出于必须。你们手中没有，也不富足。这就是\Quote{你们缺少机会}的意思。多数人这样说话，当今生的事物不丰裕地流向他们，且缺乏的时候。\par

\BibleQuote{这不是因我有所缺乏} 他说，我说，\BibleQuote{如今终于}，我责备你们，不是寻求自己的益处，也不是为此而指责你们，似乎我有所缺乏：因为我并非为此而寻求。\Person{保禄}啊，这是为什么呢，你为何不徒然夸耀？他对\Person{格林多人}说：\BibleQuote{我们给你们写的，无非是你们所读的，或甚至所承认的}。\BibleRef{格后 1:13} 在这种情况下，他不会对他们说话以致被定罪，他若夸口，就不会这样说了。他是对知道事实的人说话，被他们识破将是更大的耻辱。\BibleQuote{我已学会了}，他说，\BibleQuote{在我所处的任何境况中，知足}。因此，这是需要操练、练习和注意的对象，因为不容易达到，而是非常困难且崭新的事。\BibleQuote{在我所处的任何境况中}，他说，\BibleQuote{知足。我知道如何受卑贱，也知道如何处富余。我已在一切事上学习了秘诀}。这就是说，我知道如何使用少量，忍受饥饿和缺乏。\BibleQuote{既能处富余，也能处缺乏}。\Quote{但是，有人说，处富余不需要智慧或德行}。这需要极大的德行，与其他情况一样。因为缺乏使我们做许多恶事，富余也是如此。因为许多人一旦进入富余，就变得懒惰，不知道如何承受好运。许多人以此为由不再工作。但\Person{保禄}没有这样，因为他所领受的，都用在别人身上，并为他们倒空自己。这就是知道。他丝毫没有放松，也不因富余而自夸；在缺乏和富余中他是同样的，既不被重压，也不成为夸口者。\BibleQuote{既能饱足}，他说，\BibleQuote{又能饥饿，既能处富余，又能处缺乏}。许多人不知道如何饱足，例如以色列子民，\BibleQuote{吃了，就踢跳} \BibleRef{申 32:15}，但我在一切事上同样有序。他表明他现在既不得意，以前也不悲伤：或者他悲伤，是为了他们，不是为自己，因为他自己也有同样的感受。\par

\BibleQuote{在一切事上}，他说，\BibleQuote{并在所有的事上，我学习了秘诀}，也就是说，在这长久的时间里，我经历了万事，这些事在我都成功了。但既然在这里似乎有夸口的余地，你看他如何迅速制止，并说：\BibleQuote{我依靠那坚固我的基督，凡事都能做}。成功不是我的，而是那赐我力量者的。但那些施恩惠的人，当他们看见受惠者对他们不怀好感，反而轻视礼物时，他们自己就变得更加懈怠（因为他们认为自己是在施恩和安慰），因此如果保禄轻视这安慰，他们将必然懈怠。为了不让这事发生，你看他如何又医治了。他通过上面的话已压低了他们骄傲的思想，通过下面的话又使他们重新奋发，他说：\BibleQuote{然而你们做得很好，因为你们与我分担了我的患难}。你看，他是如何先疏远自己，然后又与它们联合。这是真正属神友谊的部分。不要想，他说，因为我并不缺乏，所以我不需要你们这一行动。我需要它是为了你们的缘故。那么，他们如何分担了他的患难呢？是这样。正如他在被囚时所说：\BibleQuote{你们都是与我同享恩宠的人}。\BibleRef{斐 1:7} 因为为基督受苦是恩宠，正如他在另一处所说：\BibleQuote{因为天主赐给你们的不但是相信祂，而且是为祂受苦}。\BibleRef{斐 1:29} 因为前面的话本身可能使他们不关心，为此他安慰他们，接纳他们，又赞美他们。而且用有分寸的话。因为他没有说\Quote{给了}，而是\Quote{分担了}，以表明他们也因成为他劳苦的参与而获益。他没有说\Quote{你们减轻了}，而是\Quote{你们与我共患难}，这是更高尚的说法。你看见保禄的谦卑吗？你看见他高贵的品性吗？当他表明他不需要他们的礼物是为了自己之后，他随后自由地使用了这样卑微的话语，如同那些请求的人一样：\Quote{因为你们素来是乐意给予的}。因为他既不拒绝做，也不拒绝说任何事。就是说：\Quote{不要以为我的话是出于羞耻，在这事上我责备你们，并说\InnerQuote{如今终于你们又关心了}，或是出于需要之人的话；我这样说不是因为需要，但为什么？是出于我对你们极大的信任，而你们自己也是这信任的原由}。\par

你看见他如何安抚他们吗？你们怎么是原由？因为你们在所有人之前急忙做了这工作，并给我信心提醒你们这些事。注意他的超然：当他们未送时，他不指责他们，免得他显得看重自己的利益；但当他们送后，他才责备他们过去的怠慢，而他们接受了，因为此后他不能再显得看重自己的利益。\par

第15节。\BibleQuote{你们自己也知道，\Place{斐理伯}人，当我在福音的起初，离开\Place{马其顿}时，没有别的教会与我交通有关给予和收受，惟独你们。} \BibleRef{斐 4:15}\par

看，他的称赞何其伟大！因为格林多人、罗马人听到他说这些话而激动起来，而斐理伯人却是没有其他教会开始就做了这事。他说：\Quote{在福音的开端}，他们就对圣宗徒显示出如此的热忱，以至在没有任何榜样的情况下，他们自己首先开始结出这果实。没有人能说他们做这些事是因为他住在他们中间，或是为他们自己的益处；因为他说：\Quote{当我离开马其顿时，没有教会与我相交，在给与收的事上，只有你们而已。}\Quote{收}是什么意思，\Quote{相交}又是什么意思？为什么他不说\Quote{没有教会给了我}，而是说\Quote{在给与收的事上与我相交}？因为这是一种互通。他说：\BibleQuote{如果我们给你们播种了属灵的事物，收获你们的属肉的事物还算大事吗？}\BibleRef{格前 9:11}又说：\BibleQuote{你们的富馀可以补他们的缺乏。}\BibleRef{格后 8:14}他们如何相通？在给与属肉的事物、领受属灵的事物上。正如买卖的人彼此相通，互相给予所拥有的（这就是相通），同样，这里也是如此。因为没有什么比这种贸易和买卖更有利。它在地上进行，但在天上完成。买的人在地上，但他们买卖并约定属天的事物，同时付出属地的代价。\par

但不要灰心；属天的事物不是用钱能买到的，财富不能购买这些事物，而是出于那给钱之人的意志、他的真智慧、他对属世事物的超脱、他对人的爱、他的怜悯。因为如果钱能买到它，那投入两个小钱的女人就不会得到什么大的收获。但既然不是钱，而是意志有效，那显出全然心志的人就得到了一切。那么，我们不要说得天国可以用钱买到；它不是靠钱，而是靠心志，而这心志是通过钱来显明的。因此，有人会回答：不需要钱吗？不是不需要钱，而是需要准备；如果你有这准备，你甚至可以用两个小钱买到天堂；如果没有这准备，即使一万他连得的金子也不能做到那两个小钱所能做的。为什么呢？因为如果你有很多却只投入一小部分，你确实给了施舍，但不及那寡妇所给的那么多；因为你没有像她那样心甘情愿地投入。因为她剥夺了自己所有的一切，或者不如说她没有剥夺，而是当作给予自己的礼物全部给了出去。不是为了一杯凉水，天主应许了天国，而是为了心的准备；不是为了死亡，而是为了心志。因为这实在不是一件大事。因为献出一条性命算什么？那只是交出一个人；但是一个人的价值并不足够。\par

第16节。\BibleQuote{因为甚至在得撒洛尼，你们也曾一次再次地供应我的需要。}\BibleRef{斐 4:16}\par

这里又是极大的称赞，即当他住在大都会时，竟由一个小城供养。并且，免得他常常以缺乏的假设而退避（如我起初所说的），使他们出差错，他先前用许多证据证明自己并不缺乏，现在只用一句话做到，就是说\Quote{需要}。他没有说\Quote{我的}，而是绝对地——关乎尊严。不仅这一点，还有随后的话，因为他意识到这是一件非常卑微的事，于是再次加以保障，补充更正说：\par

第17节。\BibleQuote{我不是求礼物，}\BibleRef{斐 4:17}\par

正如他上面所说：\Quote{我不是因缺乏而说}；那句话比这个更强。因为缺乏而不求是一回事，缺乏而不以为自己缺乏是另一回事。他说：\Quote{我不是求礼物，乃是求那归入你们账上的果子。}不是为我自己的。你们看见吗？那果子是为他们结出的？他说，我说这话是为你们，不是为我自己，是为了你们的救恩。因为我收取时一无所获，但恩宠属于施予者，因为赏报是在那边为施予者存留的，而礼物是在这里被收取者消耗的。再一次，连他的愿望也与赞美和同情结合在一起。\par

当他说了\Quote{我不求}之后，免得他再次使他们松懈，他补充说：\par

\BibleRef{斐 4:18} \BibleQuote{但我样样都有了，并且富足。}也就是说，藉这赠礼，你们补足了所缺欠的，这会使他们更加热切。因为施惠者越明智，就越求受惠者的感恩。换言之，你们不仅补足了从前的不足，而且超出了。免得这些话显得是在责备他们，看他如何封住了一切。他说了：\BibleQuote{我并不寻求赠礼}和\BibleQuote{如今终于}，并表明他们的行为是还债，因为\BibleQuote{我样样都有了}的意思正是如此，然后他又表明他们做得超过了应尽的，便说：\BibleQuote{我样样都有了，并且富足，我已充足。}我说这话不是出于偶然，也不仅出于我内心的感受，而是为什么？\BibleQuote{我从\Person{厄帕夫洛狄托}收到了你们所送来的馨香之气，可悦纳的祭献，中悦天主的。}看，他把他们的赠礼提升到了何处；不是我，他说，收了，而是天主藉着我收了。所以，尽管我不需要，也不必在意，因为天主并不需要，却从他们手中那样地收了，以至圣经毫不犹豫地说：\BibleQuote{天主闻到了馨香之气}\BibleRef{创 8:21}，这是指一位喜悦者。因为你们知道，的确你们知道，我们的灵魂如何受馨香之气的感动，如何喜悦，如何欢愉。因此圣经毫不犹豫地将如此人性化、如此卑微的词应用于天主，为向人表明他们的祭品已蒙悦纳。因为不是脂油，不是烟熏，使它们蒙悦纳，而是献祭者的心意。否则，\Person{加音}的祭品也必蒙悦纳。所以，圣经说祂甚至喜悦，并说了祂如何喜悦。因为人若不如此，就无法得知。那位无所需者，祂说祂如此喜悦，为使人不致因祂无需要而懈怠。后来，当他们不关心其他德行，只依靠献祭时，看，祂又如何藉着说\BibleQuote{我岂要吃公牛的肉，喝公山羊的血吗？}\BibleRef{咏 49:13}来纠正他们。\Person{保禄}也说了这话；他说：\BibleQuote{我并不寻求赠礼。}\par

\BibleRef{斐 4:19} \BibleQuote{我的天主必要以其光荣中的丰富，在\Person{基督耶稣}内，满足你们一切的需用。}\par

看他如何如同穷人一般祝福他们。但如果连\Person{保禄}都祝福施舍者，何况我们呢？我们在领受时不要以此为耻。我们不要领受得仿佛自己需要，不要为自己喜乐，而要为了施予者喜乐。这样，我们这些领受的人也将得到赏报，若我们为他们而喜乐。这样，当人不给时，我们不会感到难受，反而要为他们而忧伤。这样我们将使他们更为热切，如果我们教导他们说，我们这样做不是为了自己，\BibleQuote{但是我的天主}满足你们一切的需用，或一切恩宠，或一切喜乐。若是第二种，即\BibleQuote{一切恩宠}，他不仅指地上的施舍，更指一切卓越之事。若是第一种，即\BibleQuote{你们一切的需用}（我认为这更应如此读），这就是他要表明的意思。如他所说：\BibleQuote{你们缺乏机会}，他在这里作了补充，如同他在\BibleBook{格林多}{后书}中所说：\BibleQuote{那位给播种者供应种子，又赐予饼作食物的，必会供应你们的种子并使之增多，且增加你们义德的果实}\BibleRef{格后 9:10}。他祝福他们，使他们丰足，并有可以播种的。他也祝福他们，不仅简单地使他们丰足，而是\BibleQuote{按祂的丰富}，所以这也在衡量的限度内。因为如果他们像他那样，那样真智慧，那样被钉死，他不会这样做；但既然他们是工匠，贫穷，有妻子，养育儿女，管理家庭，且从微薄的财产中献出这些礼物，并对此世之物有某些欲望，他便适当地祝福了他们。因为对这样使用它们的人祈求足够与富足并非不当。再看他说了什么。他没有说：愿祂使你们富有，大大丰足；而是说什么？\BibleQuote{愿祂满足你们一切的需用}，使你们不致缺乏，而能拥有所需之物。因为\Person{基督}在教我们祈祷时，也在祈祷中加上了这一句，教导我们说：\par

\BibleQuote{求祢今日赐给我们日用的食粮}\BibleRef{玛 6:11}\par

\BibleQuote{按祂的丰富。}即按祂的白白恩赐，也就是说，对祂来说是容易的、可能的、快速的。既然我提到了需用，不要以为祂会使你陷入窘迫。因此他又加了：\BibleQuote{在\Person{基督耶稣}内，按祂在光荣中的丰富。}这样一切都要丰足地赐予你们，使你们能用以光荣祂。或者，你们一无所缺（因为经上记载：\BibleQuote{在他们众人身上大有恩宠，也没有任何人缺少什么}\BibleRef{宗 4:33}）。或者，使你们为祂的光荣做一切事，就如他曾说：使你们用你们的丰足来光荣祂。\par

\BibleRef{斐 4:20} \BibleQuote{愿光荣归于我们的天主和\Person{圣父}，至于无穷之世。阿们。}因为他所说的光荣不仅属于\Person{圣子}，也属于\Person{圣父}，因为当\Person{圣子}受光荣时，\Person{圣父}也受光荣。因为当他说这事是为\Person{基督}的光荣而做时，免得有人以为这仅是为祂的光荣，他便接着说：\BibleQuote{愿光荣归于我们的天主和\Person{圣父}}，即那归于\Person{圣子}的光荣。\par

\BibleRef{斐 4:21} \BibleQuote{你们要在\Person{基督耶稣}内问候每一位圣人。}这也不是小事。因为藉着信函问候他们，是极大善意的证明。\BibleQuote{与我在一起的弟兄们问候你们。}然而你曾说：你\BibleQuote{没有谁与我同心，真正关心你们的事}。那么你现在怎么说\BibleQuote{与我在一起的弟兄们}？他或者这样说是因为在那些与他在一起的人中，没有同心的人（他不是指城里的人，因为他们怎么被迫承担\Person{宗徒}们的事务？），或者他不拒绝称这些人为弟兄。\par

第二十二、二十三节。\BibleQuote{诸位圣人问候你们，尤其是凯撒家里的人。主耶稣基督的恩宠与你们的灵同在。}\par

他提升并坚固了他们，表明他的宣讲甚至传到了国王的家中。因为如果那些在王宫里的人为着天上君王的缘故轻看一切，那么他们更应该这样做。这也是\Person{保禄}之爱的一个明证，他告诉了他们许多事，并为他们说大话，以致他甚至引导那些在宫里的人也渴慕他们，以至于那些从未见过他们的人向他们致意。尤其因为信友当时处于患难中，他的爱是伟大的。那些彼此分离的人紧密相连，如同真正的肢体。穷人对富人和富人对穷人也有同样的心，没有优越感，因为他们都同样被憎恨、被驱逐，且为同一原因。因为就像从不同城市被掳的俘虏来到同一城镇，他们热切地拥抱彼此，共同的不幸将他们联结在一起；同样，那时他们彼此怀有极大的爱，苦难和迫害的共融使他们合一。\par

道德教训。因为患难是不可断绝的纽带，是爱的增长，是痛悔和虔敬的契机。请听\Person{达味}的话：\BibleQuote{我受苦是好的，为使我学习你的法令。}\BibleRef{咏 118:71} 又有另一位先知说：\BibleQuote{人年轻时负轭，原是好的。}\BibleRef{哀 3:27} 又说：\BibleQuote{上主，你所管教的人是有福的。}\BibleRef{咏 93:12} 另一位说：\BibleQuote{不可轻视上主的管教。}\BibleRef{箴 3:11} 又说：\BibleQuote{你若来事奉上主，当预备你的灵魂接受试探。}\BibleRef{德 11:1} 基督也对祂的门徒说：\BibleQuote{在世上你们将有苦难，但你们放心。}\BibleRef{若 16:33} 又说：\BibleQuote{你们将要痛哭哀号，世界却要欢乐。}\BibleRef{若 16:20} 又说：\BibleQuote{窄而狭的路。}\BibleRef{玛 7:14} 你岂不见患难处处被称颂，处处被视为我们所必需的吗？因为若在世上的竞赛中，没有人不经此就得冠冕，除非他藉劳苦、戒绝美味、按规律生活、守夜和无数其他事来坚固自己，那么在这里更当如此。因为你将以谁为例？君王？连他也不是无忧无虑地生活，而是肩负许多患难和焦虑。不要看他的冠冕，而要看他的忧愁之海，冠冕由此为他产生。也不要看他的紫袍，而要看他的灵魂，比那紫袍更黑暗。他的冠冕不像焦虑那样紧束着他的额头。也不要看他的众多持戟卫士，而要看他的众多不安。因为不可能找到一所私宅像王宫那样承载如此多的忧虑。每天都有暴死期待，坐下吃喝时看见血影。我们不能说他在夜间受了多少惊扰，跳起来，被幻象困扰。这一切还在和平中；但若战争临到他，还有比这更可怜的生活吗！他因自己人——我意指那些在他管辖下的人——所受的祸害何其多！而且，的确，王宫的地面常是血淋淋的，是他自己亲属的血。你若愿意，我也将叙述一些事例，你立刻就会知道；主要是古代的事——但也有我们时代发生的一些事——仍存留在记忆中。据说，有一人怀疑妻子通奸，将她裸体绑在骡子上，暴露给野兽，尽管她已经为他生了多位王子。你们想，那人过的是什么样的生活？因为若不是深受那种痛苦，他不会发泄出那样的报复。此外，这人杀了自己的儿子，或者不如说是他兄弟杀的。他的儿子中，一个被暴君擒获时自杀了，另一个杀了他的堂兄——他的王国同僚，即他任命的人；还看见他的妻子因坐药而死去，因为她不生育时，一个可怜可悲的女人（她的确是那样，以为凭自己的智慧可以代替天主的恩赐）给了她坐药，毁了王后，自己也与她同归于尽。据说这人还杀了自己的兄弟。另一个人，他的继任者，被毒药所害，他的杯不再是饮料，而是死亡。他的儿子因为害怕将要发生的事而被挖出一只眼睛，尽管他并无过错。不便提及另一人如何悲惨地结束生命。之后，有一人被烧死，像可怜虫一样，在马匹、木料和各样东西中，留下妻子守寡。因为他活着时所受的祸害无法叙述，当他起来造反时。现在统治的那位，从他接受冠冕时起，岂不就在劳苦、危险、悲伤、沮丧、不幸中，遭受阴谋吗？天国却不是这样，接受之后，有平安、生命、喜乐、欢愉。但正如我所说，没有痛苦的生活是不可能的。因为若在这世上的事中，连那被认为最幸福的，若君王都被如此多的不幸所负，你以为私人生活会是怎样的？我无法说出还有多少别的祸患！在这些问题上，有多少故事曾被编织！因为几乎舞台上的所有悲剧以及神话故事都以君王为主题。因为这些故事大多是由真实事件形成的，这样才令人愉悦。例如，\Person{提厄斯忒斯}的宴席，以及那整个家族因不幸而毁灭。\par

这些事我们是从外邦作家知道的；但若你愿意，我也要从圣经中引证例子。\Person{撒乌耳}是第一个君王，你知道他经历了无数灾祸后如何丧亡。此后，\Person{达味}、\Person{撒罗满}、\Person{阿彼雅}、\Person{希则克雅}、\Person{约史雅}，也大致如此。因为不经历忧苦和辛劳，不经历心灵的沮丧，就无法度过今生。但我们的心灵不应为这些事沮丧——君王为之忧伤的那些事——而应为那些使我们大得利益的事沮丧。\BibleQuote{因为依着天主而来的忧苦，能产生非导致后悔的悔改，以致得救。}\BibleRef{格后 7:10} 我们应为这些事忧伤，为这些事痛苦，为这些事心中刺痛；\Person{保禄}就是这样为罪人忧伤，就是这样哭泣。\BibleQuote{因为我从极大的忧苦和心中的痛苦，流着许多泪给你们写了信。}\BibleRef{格后 2:4} 因为他本人并无忧伤的缘由，却为他人如此，或者说，至少在忧伤方面，他将那些事也视为自己的。别人跌倒，他就焦心；别人软弱，他就软弱。这样的忧伤是好的，超越一切世俗的喜乐。我喜爱这样忧伤的人胜过所有人；或者更确切地说，主亲自称那些忧伤、富有同情心的人为有福的。我不那么钦佩他在危难中的表现——或者说不那么钦佩他每日冒死的危难——但这一点更令我着迷。因为它出自一颗奉献于天主、充满爱心的灵魂；出自基督自己所寻求的爱；出自兄弟般的、父亲般的同情，或者更超乎两者之上。我们应当这样被感动，应当这样哭泣；这样的眼泪饱含极大的喜乐；这样的忧伤是喜乐的基础。\par

不要对我说：我所忧伤的人能得到什么益处呢？即使我们无法使所忧伤的人得益，至少我们会使自己得益。因为为他人这样忧伤的人，更会为自己忧伤；为别人的罪这样哭泣的人，不会不为自己过犯哭泣，或者说，他不会轻易犯罪。但可怕的是，当我们受命为犯罪的人这样忧伤时，我们甚至对自己所犯的罪毫不悔改，反而在犯罪时毫无感觉，只关心别的事，而不顾及自己的罪。因此，我们以无价值的喜乐为乐，那是世俗的喜乐，转瞬即逝，并生出无数的忧伤。让我们以忧伤为忧，那是喜乐之母；不要以喜乐为乐，那是生出忧伤的喜乐。让我们流泪，那是大喜乐的种子；不要发出那笑声，那会为我们带来咬牙切齿的哀号。让我们受苦，那是安逸的源头；不要追求奢侈，那是大忧苦与痛苦的根源。让我们在地上劳苦片刻，好能在天上享受永恒。让我们在这短暂的今生受苦，好能在无尽的来世得享安息。让我们不要在这短暂的生命中懈怠，免得在那无尽的生命中叹息。\par

你们岂不见许多人为世俗之事在此受苦吗？你要想到自己也是其中之一，并要忍受你的忧苦和痛苦，以对将来之事的希望为食粮。你并不比\Person{保禄}或\Person{伯多禄}更强，他们从未得享安逸，终生忍饥挨渴、衣不蔽体。你若愿与他们得到同样的一切，为何走相反的道路？你若愿到达那座他们堪当进入的城，就当走通往那里的路。安逸之路不通向那里，只有受苦之路通达。前者是宽阔的，后者是狭窄的；让我们走这条路，好能获得在基督耶稣我们的主内的永生；愿光荣、权能、尊威归于他，偕同圣父及圣神，从现在直到永远，世世无尽。阿们。\par

\SourceRecord{Translated by John A. Broadus.；Nicene and Post-Nicene Fathers, First Series；Vol. 13.；Edited by Philip Schaff.；Buffalo, NY: Christian Literature Publishing Co.,；1889.；<http://www.newadvent.org/fathers/230215.htm>.}
